% concatenated from otSolverNips2026-arxiv.tex
\documentclass{article}

% if you need to pass options to natbib, use, e.g.:
     \PassOptionsToPackage{numbers, compress}{natbib}
% before loading neurips_2025

% The authors should use one of these tracks.
% Before accepting by the NeurIPS conference, select one of the options below.
% 0. "default" for submission
 \usepackage[preprint]{neurips_2026}
% the "default" option is equal to the "main" option, which is used for the Main Track with double-blind reviewing.
% 1. "main" option is used for the Main Track
%  \usepackage[main]{neurips_2025}
% 2. "position" option is used for the Position Paper Track
%  \usepackage[position]{neurips_2025}
% 3. "dandb" option is used for the Datasets & Benchmarks Track
 % \usepackage[dandb]{neurips_2025}
% 4. "creativeai" option is used for the Creative AI Track
%  \usepackage[creativeai]{neurips_2025}
% 5. "sglblindworkshop" option is used for the Workshop with single-blind reviewing
 % \usepackage[sglblindworkshop]{neurips_2025}
% 6. "dblblindworkshop" option is used for the Workshop with double-blind reviewing
%  \usepackage[dblblindworkshop]{neurips_2025}

% After being accepted, the authors should add "final" behind the track to compile a camera-ready version.
% 1. Main Track
 % \usepackage[main, final]{neurips_2025}
% 2. Position Paper Track
%  \usepackage[position, final]{neurips_2025}
% 3. Datasets & Benchmarks Track
 % \usepackage[dandb, final]{neurips_2025}
% 4. Creative AI Track
%  \usepackage[creativeai, final]{neurips_2025}
% 5. Workshop with single-blind reviewing
%  \usepackage[sglblindworkshop, final]{neurips_2025}
% 6. Workshop with double-blind reviewing
%  \usepackage[dblblindworkshop, final]{neurips_2025}
% Note. For the workshop paper template, both \title{} and \workshoptitle{} are required, with the former indicating the paper title shown in the title and the latter indicating the workshop title displayed in the footnote.
% For workshops (5., 6.), the authors should add the name of the workshop, "\workshoptitle" command is used to set the workshop title.
% \workshoptitle{WORKSHOP TITLE}

% "preprint" option is used for arXiv or other preprint submissions
 % \usepackage[preprint]{neurips_2025}

% to avoid loading the natbib package, add option nonatbib:
%    \usepackage[nonatbib]{neurips_2025}


\usepackage[utf8]{inputenc} % allow utf-8 input
\usepackage[T1]{fontenc}    % use 8-bit T1 fonts
\usepackage{hyperref}       % hyperlinks
\usepackage{url}            % simple URL typesetting
\usepackage{booktabs}       % professional-quality tables
\usepackage{amsfonts}       % blackboard math symbols
\usepackage{nicefrac}       % compact symbols for 1/2, etc.
\usepackage{microtype}      % microtypography
\usepackage{xcolor}         % colors
\usepackage{graphicx}       % graphics and images
\usepackage{subcaption}     % subfigure environment
\usepackage{amsmath}
\usepackage{amsthm}
\usepackage{algorithm,algorithmic}
\usepackage{amssymb}
\newcommand{\dd}{\,{\rm d}}

\theoremstyle{plain}
\newtheorem{theorem}{Theorem}
\newtheorem{definition}{Definition}
\newtheorem{lemma}{Lemma}
\newtheorem{assumption}{Assumption}
\newtheorem{corollary}{Corollary}
\newtheorem{proposition}{Proposition}
\newtheorem{proposition*}{Proposition}
\newtheorem*{theorem*}{Theorem}
\newtheorem*{lemma*}{Lemma}

\theoremstyle{remark}
\newtheorem*{remark}{Remark}

\newcommand{\dual}[1]{\left\langle {#1} \right\rangle}

\usepackage{comment}
\newcommand{\ZX}[1]{\textcolor{blue}{#1}}
\newcommand{\LC}[1]{\textcolor{red}{#1}}
% Note. For the workshop paper template, both \title{} and \workshoptitle{} are required, with the former indicating the paper title shown in the title and the latter indicating the workshop title displayed in the footnote. 
\title{Accelerating Sinkhorn for Entropy-Regularized Optimal Transport}


% The \author macro works with any number of authors. There are two commands
% used to separate the names and addresses of multiple authors: \And and \AND.
%
% Using \And between authors leaves it to LaTeX to determine where to break the
% lines. Using \AND forces a line break at that point. So, if LaTeX puts 3 of 4
% authors names on the first line, and the last on the second line, try using
% \AND instead of \And before the third author name.


\author{%
  Zeyi Xu \\
  Department of Mathematics\\
  University of California, Irvine\\
  Irvine, CA 92697 \\
  \texttt{zeyix1@uci.edu} \\
  % examples of more authors
 \And
 Long Chen \\
 Department of Mathematics\\
  University of California, Irvine\\
  Irvine, CA 92697 \\
  \texttt{chenlong@math.uci.edu} \\
  % \AND
  % Coauthor \\
  % Affiliation \\
  % Address \\
  % \texttt{email} \\
  % \And
  % Coauthor \\
  % Affiliation \\
  % Address \\
  % \texttt{email} \\
  % \And
  % Coauthor \\
  % Affiliation \\
  % Address \\
  % \texttt{email} \\
}


\begin{document}


\maketitle

\begin{abstract}
We propose Acc-Sinkhorn, a simple accelerated variant of Sinkhorn for entropy-regularized optimal transport (EOT). The method is derived from a bilevel optimization view: Sinkhorn row scaling solves the inner variable $u$ exactly and defines the reduced dual objective $f(v)=\min_u F(u,v)$, while the remaining column scaling is a unit-step dual mirror descent step in $v$. This structure yields a Hessian-driven Nesterov acceleration that keeps Sinkhorn's scaling form and per-iteration cost, using only extrapolated combinations of Sinkhorn iterates. We prove an $\mathcal{O}(1/k^2)$ rate under a verifiable stability condition. For an $\varepsilon$-approximation of unregularized OT, the resulting complexity is $\widetilde{\mathcal{O}}(n^2/\varepsilon)$, improved from $\widetilde{\mathcal{O}}(n^2/\varepsilon^2)$ for Sinkhorn. On synthetic problems, color transfer, and word alignment, Acc-Sinkhorn gives a $10\times$--$30\times$ speedup over Sinkhorn at small regularization.
\end{abstract}

\section{Introduction}

\paragraph{Entropic optimal transport.}
Given two discrete distributions $a \in \mathbb{R}^n_+$ and $b \in \mathbb{R}^m_+$, the \emph{entropy-regularized optimal transport} (EOT) problem is
\begin{equation}\label{eq:EOT}
\min_{P \in \Pi(a,b)} \langle C, P \rangle + \varepsilon \sum_{i,j} P_{ij}(\log P_{ij}-1),
\end{equation}
where $C \in \mathbb{R}^{n \times m}_+$ is a cost matrix, $\varepsilon>0$ controls the strength of the entropic regularization, and
$$
\Pi(a,b):=\Big \{P\in\mathbb{R}^{n\times m}_+:\; \sum_{j} P_{ij}=a_i,\; i=1,\ldots,n,\quad \sum_{i} P_{ij}=b_j,\; j=1,\ldots,m\Big \}
$$
is the transportation polytope. 
Since its introduction to machine learning by \citet{cuturi2013sinkhorn}, EOT has become a standard tool in generative modeling, domain adaptation, image processing, computer vision, and natural language processing~\citep{genevay2018learning,courty2017joint,solomon2015convolutional,xu2019learning,alvarez2020unsupervised}.

The original OT problem
\begin{equation}\label{eq:OT}
\min_{P \in \Pi(a,b)} \langle C, P \rangle
\end{equation}
is a linear program and can be solved by standard LP methods, including simplex and interior-point methods~\citep{peyre2019computational}. These methods scale poorly for large problems: $P$ has $nm$ entries, interior-point methods require large KKT solves, and simplex-type methods may need many pivots on degenerate transport polytopes. In contrast, EOT has a unique positive minimizer and admits simple $\mathcal O(nm)$-per-iteration algorithms that exploit the transport structure and are easy to parallelize on GPUs.

\paragraph{Sinkhorn algorithm.}
The Sinkhorn algorithm \citep{sinkhorn1964} alternately rescales the rows and columns of the kernel matrix $K=\exp(-C/\varepsilon)$ to enforce the marginal constraints. Each iteration costs $\mathcal O(nm)$ and is easy to parallelize. Its convergence has been studied through Hilbert's projective metric, block coordinate descent, sharp $\varepsilon$-dependent estimates, and mirror descent in continuous and discrete time~\citep{franklin1989scaling,carlier2022linear,ghosal2025convergence,chizat2026sharper,pmlr-v238-reza-karimi24a,mishchenko2019sinkhorn,leger2021gradient,aubin2022mirror,pmlr-v235-chopin24a}.

Despite these advances, a central question remains: \emph{can Sinkhorn be accelerated to improve its accuracy dependence without losing its simplicity?}

% \paragraph{Existing approaches for accelerating Sinkhorn.}
\paragraph{Related Work.}
We call a method accelerated if its error decays as $\mathcal{O}(1/k^2)$ rather than the $\mathcal{O}(1/k)$ rate of Sinkhorn, where $k$ is the number of full Sinkhorn-type iterations. For simplicity, we discuss complexity for $m=n$. For an $\varepsilon$-approximation of the original OT problem, the total cost is often $\mathcal{O}(n^p/\varepsilon^q)$ for positive $p$ and $q$. Existing acceleration methods fall into three broad categories.

\smallskip
\noindent\textbf{(i) Coordinate-wise and stochastic methods.}
Coordinate-wise and stochastic variants, such as Greenkhorn, reduce the cost of each iteration by updating only part of the variables or by using stochastic estimates~\citep{altschuler2017near,lin2019efficient,genevay2016stochastic}. These methods can improve practical efficiency, but they retain the $\mathcal{O}(1/k)$ convergence behavior of Sinkhorn and thus do not improve its dependence on $\varepsilon$.

\smallskip
\noindent\textbf{(ii) Nesterov acceleration and primal-dual methods.}
Accelerated primal-dual methods, such as APDAGD, improve the $\varepsilon$-dependence of Sinkhorn by applying Nesterov-type acceleration or accelerated mirror descent to the EOT dual problem~\citep{dvurechensky2018computational,lin2019efficient}. However, they typically require line search or adaptation to local smoothness and are more complex to implement than Sinkhorn.

\smallskip
\noindent\textbf{(iii) Sinkhorn variants with modified update rules.}
Sinkhorn variants with modified update rules, such as overrelaxed Sinkhorn and annealed Sinkhorn, keep the scaling structure of Sinkhorn while changing the update or regularization schedule~\citep{thibault2021overrelaxed,chizat2024annealed}. They can improve local convergence or practical behavior, but they do not give a global $\mathcal{O}(1/k^2)$ acceleration with a parameter-free update.

\smallskip
In short, existing methods do not simultaneously achieve an accelerated $\mathcal{O}(1/k^2)$ rate, $\mathcal O(n^2/\varepsilon)$ complexity, and a parameter-free simple implementation; see Table~\ref{tab:comparison}.

\begin{table}[htp]
\centering
\caption{Comparison of algorithms for entropy-regularized optimal transport ($n=m$).}
\label{tab:comparison}
\renewcommand{\arraystretch}{1.3}
\begin{tabular}{lccc}
\toprule
\textbf{Method}
  & \textbf{Rate}
  & \textbf{Total complexity}
  & \textbf{Parameters $L,\mu$} \\
\midrule
Sinkhorn \citep{cuturi2013sinkhorn,altschuler2017near}
  & $\mathcal{O}(1/k)$
  & $\widetilde{\mathcal O}(n^2/\varepsilon^2)$
  & free \\

Greenkhorn \citep{altschuler2017near,lin2019efficient}
  & $\mathcal{O}(1/k)$
  & $\widetilde{\mathcal O}(n^2/\varepsilon^2)$
  & free \\

Overrelaxed Sinkhorn \citep{thibault2021overrelaxed}
  & $O(\rho^k)$ (local)
  & $\widetilde{\mathcal O}(n^2/\varepsilon)$ (local)
  & spectral parameter \\

APDAGD \citep{dvurechensky2018computational,lin2019efficient}
  & $\mathcal{O}(1/k^2)$
  & $\widetilde{\mathcal O}(n^{5/2}/\varepsilon)$
  & linesearch for $L$ \\

APDAMD \citep{lin2019efficient}
  & $\mathcal{O}(1/k^2)$
  & $\widetilde{\mathcal O}(n^2\sqrt{\delta}/\varepsilon)$
  & linesearch for $L$ \\

Annealed Sinkhorn \citep{chizat2024annealed}
  & $\mathcal{O}(1/\sqrt{k})$
  & heuristic
  & schedule \\

\midrule
\textbf{Acc-Sinkhorn (ours)}
  & $\mathcal{O}(1/k^2)$
  & $\widetilde{\mathcal O}(n^2/\varepsilon)$
  & $L=1$; homotopy $\mu$ \\
\bottomrule
\end{tabular}
\end{table}

\paragraph{Our contribution.}
The main contribution of this paper is \textbf{Acc-Sinkhorn}, a simple accelerated variant of Sinkhorn that keeps the $\mathcal O(nm)$ per-iteration cost of Sinkhorn while improving the convergence rate from $\mathcal{O}(1/k)$ to $\mathcal{O}(1/k^2)$. As shown in Table~\ref{tab:comparison}, Acc-Sinkhorn is the only method in the table that combines accelerated convergence, Sinkhorn-level per-iteration cost, and a simple update rule without line search. Its smoothness parameter is fixed by the geometry, namely $L=1$, and the strong-convexity parameter $\mu$ is handled by a homotopy schedule rather than by manual tuning. Numerically, Acc-Sinkhorn achieves a $10\times$--$30\times$ speedup over Sinkhorn on synthetic datasets, color transfer, and word alignment at small $\varepsilon$.

The acceleration is built on a bilevel view of Sinkhorn. Let $F(u,v)$ be the dual objective of \eqref{eq:EOT}, where $(u,v)$ are the Lagrange multipliers for the marginal constraints. For each fixed $v$, the inner problem
$$
u(v)=\arg\min_u F(u,v)
$$
is exactly the row-scaling step of Sinkhorn. This defines the reduced outer objective
$$
f(v)=F(u(v),v)=\min_u F(u,v),
$$
and Sinkhorn can be viewed as an exact inner solve in $u$ followed by an outer descent step in $v$.

With the mirror function $\phi(v)=\sum_j b_j(e^{v_j}-v_j)$, the outer Sinkhorn step can be written as the unit-step dual mirror descent update
\begin{equation}
%\tag{\textsc{Sinkhorn}}
v_{k+1}=v_k-\nabla\phi^*(\nabla f(v_k)).
\end{equation}
Moreover, the dual relative smoothness constant is $L=1$:
$$
\nabla^2 \phi^*(\nabla f(v)) \preceq \nabla^2 f^*(\nabla f(v)).
$$

Using this bilevel structure, Acc-Sinkhorn (Algorithm~\ref{alg:A2MD}) applies a Hessian-driven Nesterov acceleration gradient flow~\citep{chen2019orderoptimizationmethodsbased} with nonlinear preconditioning to the reduced outer problem. The method keeps two sequences $(x_k,y_k)$ and uses essentially one Sinkhorn step per iteration; the extra cost is only a few vector operations. A homotopy outer loop (Algorithm~\ref{alg:A2MD-homotopy}) decreases the unknown strong-convexity parameter $\mu$ by a predefined schedule and gives an accelerated $\mathcal{O}(1/k^2)$ rate. Acc-Sinkhorn needs $\mathcal{O}(\tau^{-1/2})$ Sinkhorn steps to reach accuracy $\tau$, whereas Sinkhorn needs $\mathcal{O}(\tau^{-1})$ steps.

For the unregularized OT problem, this inner complexity is combined with the standard entropic approximation error $\mathcal O(\varepsilon\log n)$. Balancing optimization and regularization errors by setting $\tau=\varepsilon^2$ gives an overall complexity of order $\mathcal O(n^2/\varepsilon)$ for an $\varepsilon$-approximation of the OT cost.

%\paragraph{Contributions.}
%Our contributions are threefold:
%\begin{enumerate}% [leftmargin=16pt,itemsep=2pt]
%\item \textbf{Dual mirror descent view of Sinkhorn.}
%We prove that Sinkhorn is a dual mirror descent method with relative smoothness constant $L=1$ (Section~3), which yields an $\mathcal{O}(1/k)$ convergence rate.
%
%\item \textbf{Acc-Sinkhorn with $\mathcal{O}(1/k^2)$ convergence.}
%We propose Algorithm~\ref{alg:A2MD} and Algorithm~\ref{alg:A2MD-homotopy}, prove a global accelerated convergence rate, and show that the step size is governed by the KL asymmetry of the marginals (Section~4).
%
%\item \textbf{Empirical validation.}
%\end{enumerate}
\paragraph{Limitations.}
In our experiments, Acc-Sinkhorn uses the step size $\alpha=\sqrt{2\mu}$ and shows stable accelerated behavior. The current proof is more conservative: it gives the full accelerated guarantee under a stability condition, which holds for sufficiently small $\alpha$ and can be enforced by linesearch. The gap comes from controlling the metric-variation term $\|y_k-x^\star\|_{\mathcal D_{k+1}-\mathcal D_k}^2$, which is small in practice but hard to bound uniformly. Removing this condition and justifying the step size $\alpha=\sqrt{2\mu}$ remain future work.

%\paragraph{Organization.}
%Section~2 reviews the background. Section~3 establishes the dual mirror descent formulation of Sinkhorn. Section~4 presents Acc-Sinkhorn and its convergence analysis. Section~5 reports numerical results.

\section{Preliminaries}
\label{sec:prelim}

\paragraph{Notations.}
For two vectors $x,y\in\mathbb{R}^n$, we write $x./y$ and $x.*y$ for entrywise division and multiplication, and $\exp(x)$ and $\log(x)$ for the entrywise exponential and logarithm. For a vector $x\in\mathbb{R}^n$, $\mathrm{diag}(x)\in\mathbb{R}^{n\times n}$ denotes the diagonal matrix with diagonal $x$. For a positive definite matrix $A$, we write $\|x\|_A:=(x^\top A x)^{1/2}$ for the $A$-weighted norm.

We write $\mathbf{1}_n\in\mathbb{R}^n$ for the all-ones vector, dropping the subscript when the dimension is clear from context. Let $N=\mathrm{span}\{\mathbf{1}_n\}$, and $N^\perp$ be its orthogonal complement, then $N^\perp=\{x\in\mathbb{R}^n:\langle x,\mathbf{1}_n\rangle=0\}$. Let $P_N=\mathbf{1}_n\mathbf{1}_n^\top/n$ be the orthogonal projection onto $N$, and $P_{N^\perp}=I-P_N$ be the projection onto $N^\perp$. 

%For a pair of vectors $u\in\mathbb{R}^n$ and $v\in\mathbb{R}^m$, we denote by $P(u,v)$ the matrix with entries $P_{ij}(u,v)=\exp(u_i+v_j-C_{ij})$, and by.

For a convex function $\phi$, we write $\phi^*$ for its Fenchel conjugate, $\nabla\phi$ for its gradient, $\nabla^2\phi$ for its Hessian when exists, and the induced Bregman divergence
$$
D_\phi(x,y):=\phi(x)-\phi(y)-\langle \nabla\phi(y),x-y\rangle \ge 0.
$$
When $\phi$ is of Legendre type, the maps $\nabla\phi$ and $\nabla\phi^*$ are inverses of each other. In particular,
$$
\nabla\phi^*(\nabla\phi(x))=x,
\qquad
\nabla\phi(\nabla\phi^*(\xi))=\xi.
$$
Therefore, whenever $\nabla f(x)$ lies in the domain of $\nabla\phi^*$, we  have
\begin{equation}\label{eq:dphiphistar}
\nabla\phi(\nabla\phi^*(\nabla f(x)))=\nabla f(x).
\end{equation}



% The KL divergence between nonnegative vectors $p,q\in\mathbb{R}^n_{++}$
% is $\mathrm{KL}(p\|q):=\sum_i p_i\log(p_i/q_i)-(p_i-q_i)$,
% which reduces to the standard KL divergence when $\mathbf{1}^\top p
% =\mathbf{1}^\top q$.
% We use $O(\cdot)$ and $\widetilde{O}(\cdot)$ for standard and
% logarithm-hiding asymptotic notation.

\paragraph{Dual Formulation.}
First, note that the minimizer of \eqref{eq:EOT} is unchanged if we divide the cost matrix $C$ by $\varepsilon$. We therefore assume $\varepsilon=1$ in \eqref{eq:EOT} for simplicity and numerical stability. The complexity bounds, however, depend on this rescaling and will be discussed later.

Introducing Lagrange multipliers $u\in\mathbb{R}^n$ and $v\in\mathbb{R}^m$ for the row and column marginal constraints, we write the Lagrangian as
$$
\mathcal{L}(P,u,v)=\langle C,P\rangle+\sum_{i,j}P_{ij}(\log P_{ij}-1)+\langle u,a-P\mathbf{1}_m\rangle+\langle v,b-P^\top\mathbf{1}_n\rangle.
$$
Solving the stationarity condition $\partial_P \mathcal{L}(P,u,v)=0$ gives
\begin{equation}\label{eq:P-from-dual}
P_{ij}(u,v)=\exp\Bigl(u_i+v_j-C_{ij}\Bigr).
\end{equation}
Substituting \eqref{eq:P-from-dual} into $\mathcal{L}$ yields the dual objective
\begin{equation}\label{eq:dual-obj}
F(u,v):=- \mathcal L (P(u,v), u, v) = \sum_{i,j}\exp(u_i+v_j-C_{ij})-\langle u,a\rangle-\langle v,b\rangle,
\end{equation}
so EOT \eqref{eq:EOT} is equivalent to the unconstrained dual minimization problem
\begin{equation}\label{eq:dual-min}
\min_{u\in\mathbb{R}^n,\,v\in\mathbb{R}^m} F(u,v).
\end{equation}

By the chain rule and the stationarity condition $\partial_P\mathcal{L}(P(u,v),u,v)=0$, the dependence of $P(u,v)$ drops out when differentiating the reduced Lagrangian $\mathcal L(P(u,v),u,v)$. Thus
\begin{equation}\label{eq:grad-F}
\begin{aligned}
\partial_u F(u,v)=r_P(u,v)-a, \quad
\partial_v F(u,v)=c_P(u,v)-b,
\end{aligned}
\end{equation}
where $r_P(u,v)$ and $c_P(u,v)$ are row-sum and column-sum vectors of $P$. Hence $\nabla F(u,v)=0$ is exactly the marginal-matching condition.

Since $\mathbf{1}^\top a=\mathbf{1}^\top b=1$, for any $c\in\mathbb{R}$ we have
\begin{equation}\label{eq:Fc}
F(u+c\mathbf{1}_n,v-c\mathbf{1}_m)=F(u,v)+c\,\mathbf{1}^\top a-c\,\mathbf{1}^\top b=F(u,v).
\end{equation}
This is the only ambiguity. Hence, after fixing a gauge, for example by imposing $\langle u,\mathbf{1}_n\rangle=0, \langle v,\mathbf{1}_m\rangle=0$, the minimizer of \eqref{eq:dual-min} is unique.

\paragraph{Sinkhorn as Exact Coordinate Minimization}

A key property of $F$ is that it is separable in each block. For fixed $v$, the function $F(\cdot,v)$ is strictly convex in $u$, and its minimizer is available in closed form. Indeed, fixing $v$ and setting $\partial_u F(u,v)=0$ gives
$$
r_P(u,v)=a
\;\iff\;
\exp(u_i)\sum_j \exp(v_j-C_{ij})=a_i
\;\iff\;
u_i=\log a_i-\log\!\bigl(K\exp(v)\bigr)_i,
$$
where $K:=\exp(-C)$, and $K\exp(v)$ is a matrix-vector product. Rewriting this minimizer in iterative form gives
\begin{equation}\label{eq:u-update}
u^{k+1}=u^k+\log(a\,./\,r_P(u^k,v^k)).
\end{equation}
The same argument, applied to the $v$-block, gives the exact minimizer of $F(u,\cdot)$ for fixed $u$:
\begin{equation}\label{eq:v-update}
v^{k+1}=v^k+\log(b\,./\,c_P(u^{k+1},v^k)).
\end{equation}

Equations \eqref{eq:u-update}--\eqref{eq:v-update} are exactly the Sinkhorn algorithm: alternating exact minimization of the dual objective $F(u,v)$ over the two blocks $u$ and $v$.
% Because $F$ is strictly convex in each block and its level sets
% are bounded (modulo the invariance $F(u+c,v-c)=F(u,v)$),
% standard block coordinate descent theory guarantees convergence
% to a dual minimizer $(u^\star,v^\star)$, and hence to the
% primal solution $P^\star$ via \eqref{eq:P-from-dual}.

% \paragraph{Reduced dual objective.}
% Since the $u$-update is available in closed form, we can
% \emph{eliminate} $u$ and work with the reduced objective
% \begin{equation}\label{eq:reduced-dual}
% f(v) := \min_u F(u,v) = F(u(v),\,v),
% \end{equation}
% where $u(v)$ is the unique minimizer of $F(\cdot,v)$.
% By the envelope theorem, $\partial_u F(u(v),v)=0$, so
% \begin{equation}\label{eq:grad-f}
% \nabla f(v) = \nabla_v F(u(v),v) = c_P(u(v),v)-b,
% \end{equation}
% and the Sinkhorn $v$-update \eqref{eq:v-update}
% is a descent step on $f$.
% The Hessian of $f$ is given by the Schur complement formula
% \begin{equation}\label{eq:hessian-f}
% \nabla^2 f(v)
% = \mathrm{diag}(c_P) - P^\top\mathrm{diag}(r_P)^{-1}P
% \;\succeq\; 0,
% \end{equation}
% which is positive semidefinite with $\ker(\nabla^2 f(v))
% =\mathrm{span}\{\mathbf{1}\}$, reflecting the
% invariance $f(v+c\mathbf{1})=f(v)$.

%\paragraph{Dual Mirror Descent}
%
%We now recall the dual mirror descent for solving convex optimization $\min_v f(v)$ that will be used
%throughout the paper.
%Let $\phi:\mathbb{R}^m\to\mathbb{R}$ be a strictly convex,
%differentiable function.
%% Its \emph{Bregman divergence} is
%% \[
%% D_\phi(x,y):=\phi(x)-\phi(y)-\langle\nabla\phi(y),x-y\rangle\ge 0.
%% \]
%The \emph{mirror descent} iteration for minimizing $f$ is
%\begin{equation}\label{eq:MD}
%% v^{k+1}
%% = \arg\min_{v}\Bigl\{
%% \langle\nabla f(v^k),v\rangle + \tfrac{1}{\eta}D_\phi(v,v^k)
%% \Bigr\}
%% = (\nabla\phi)^{-1}\!\bigl(\nabla\phi(v^k)-\eta\nabla f(v^k)\bigr).
%\nabla \phi(v^{k+1}) = \nabla\phi(v^k) - \alpha\nabla f(v^k).
%\end{equation}
%For $\phi=\tfrac{1}{2}\|\cdot\|^2$, \eqref{eq:MD} reduces to
%standard gradient descent.
%
%An alternative preconditioned formulation works in the
%\emph{dual space} via the Fenchel conjugate
%$\phi^*(s):=\sup_v\{\langle s,v\rangle-\phi(v)\}$.
%Assume additionally that $\nabla\phi^*(0)=0$.
%The \emph{dual mirror descent} iteration is
%\begin{equation}\label{eq:DMD}
%v^{k+1} = v^k - \nabla\phi^*\!\bigl(\nabla f(v^k)\bigr).
%\end{equation}
%The two formulations are equivalent when we switch the roles of variables and operators: \begin{equation}\label{eq:equivalence}
%f\leftrightarrow \phi^*, \phi\leftrightarrow f^*, x\leftrightarrow\chi=\nabla f(x).
%\end{equation}
%
%\paragraph{Convergence Analysis.}
%A key concept for the convergence analysis is
%\emph{relative smoothness} \citep{lu2018relatively}.
%\begin{definition}\label{def:rel-smooth}
%A differentiable function $f$ is $L$-\emph{dual smooth relative to}
%$\phi$ if
%\begin{equation}\label{eq:rel-smooth}
%D_f(x,y)\leq L D_\phi(x,y)
%\qquad\forall\,x,y.
%\end{equation}
%Equivalently, when $f$ and $\phi$ are twice differentiable,
%\begin{equation}\label{eq:rel-smooth-hess}
%\nabla^2 f(v) \preceq L\,\nabla^2\phi(v)
%\qquad\forall\,v.
%\end{equation}
%\end{definition}
%Under the role change \eqref{eq:equivalence}, we get the definition of dual relative smoothness.
%
%% \paragraph{Dual relative smoothness.}
%% An equivalent dual condition \citep{lu2018relatively} states that
%% $f$ is $L$-smooth relative to $\phi$ if and only if
%% \begin{equation}\label{eq:dual-rel-smooth}
%% D_{\phi^*}(\xi,\eta) \leq L\,D_{f^*}(\xi,\eta)
%% \qquad\forall\,\xi,\eta,
%% \end{equation}
%% or equivalently, when both are twice differentiable,
%% \begin{equation}\label{eq:dual-rel-smooth-hess}
%% \nabla^2\phi^*(\xi) \preceq L\,\nabla^2 f^*(\xi)
%% \qquad\forall\,\xi.
%% \end{equation}
%
%\begin{proposition}[Convergence of Mirror Descent, \citep{lu2018relatively}]\label{thm:MD-convergence}
%Suppose $f$ is $L$-smooth relative to $\phi$ and let
%$\{v^k\}$ be generated by mirror descent \eqref{eq:MD} with
%step size $\eta=1/L$.
%Then for any minimizer $v^\star$ of $f$,
%\begin{equation}\label{eq:MD-rate}
%f(v^k) - f(v^\star)
%\leq \frac{L\,D_\phi(v^\star,v^0)}{k}.
%\end{equation}
%% Moreover, if $f$ is additionally $\mu$-strongly convex relative
%% to $\phi$, i.e.\
%% $\nabla^2 f(v)\succeq\mu\,\nabla^2\phi(v)$,
%% then mirror descent converges linearly:
%% \begin{equation}\label{eq:MD-linear-rate}
%% D_\phi(v^\star,v^k)
%% \leq \left(1-\frac{\mu}{L}\right)^k D_\phi(v^\star,v^0).
%% \end{equation}
%\end{proposition}
%
%% Theorem~\ref{thm:MD-convergence} motivates the design of a mirror
%% function $\phi$ that (a) makes the relative smoothness constant $L$
%% as small as possible, and (b) admits a cheap closed-form update
%% \eqref{eq:DMD}.
%% In Section~\ref{sec:sinkhorn-dmd} we show that the mirror function
%% \begin{equation}\label{eq:phi}
%% \phi(v) = \sum_{j=1}^m b_j(e^{v_j}-v_j)
%% \end{equation}
%% achieves both goals simultaneously for the reduced dual
%% objective $f$, with $L=1$, and that the resulting DMD iteration
%% \eqref{eq:DMD} is \emph{exactly} the Sinkhorn $v$-update
%% \eqref{eq:v-update}.

\section{Sinkhorn as Dual Mirror Descent}
\label{sec:sinkhorn-dmd}
In this section, we interpret Sinkhorn as dual mirror descent in the variable $v$, which provides the basis for the acceleration scheme in Section~4. This geometric view also gives sublinear and linear convergence of Sinkhorn; the details are deferred to Appendix~\ref{app:proof-Sinkhorn} as our focus is acceleration.

We first rewrite the dual problem from a bilevel point of view. The variable $u$ is the inner variable and $v$ is the outer variable. For each fixed $v$, the inner problem
$$
u(v)=\arg\min_u F(u,v)
$$
has a closed-form solution, and its optimality condition $\partial_u F\bigl(u(v),v\bigr)=0$ is exactly the $u$-update \eqref{eq:u-update}. This defines the reduced outer objective
\begin{equation}
f(v):=F\bigl(u(v),v\bigr)=\min_u F(u,v).
\end{equation}
Thus Sinkhorn solves the inner problem in $u$ exactly and then updates the outer variable $v$ by a gradient-type iteration. Since $\partial_u F\bigl(u(v),v\bigr)=0$, the chain rule gives
\begin{equation}\label{eq:grad-f}
\nabla f(v)
=
\partial_u F\bigl(u(v),v\bigr)\,\partial_v u(v)
+\nabla_v F\bigl(u(v),v\bigr)
=
\nabla_v F\bigl(u(v),v\bigr)
=
c_P\bigl(u(v),v\bigr)-b.
\end{equation}
Hence $\nabla f(v)$ is the column marginal residual after the row marginal has been matched exactly.

We now show that the outer Sinkhorn update is a \emph{dual mirror descent} step on $f$. The Hessian of $f$ has the Schur-complement form
\begin{equation}\label{eq:HessianL}
\nabla^2 f(v)=\mathrm{diag}(c_P)-P^\top \mathrm{diag}(r_P)^{-1}P\preceq \mathrm{diag}(c_P).
\end{equation}
Let $\xi=\nabla f(v)=c_P-b$. We define the mirror geometry by
$$
\nabla^2\phi^*(\xi):=\mathrm{diag}(b+\xi)^{-1}.
$$
Since $b+\xi=c_P$, \eqref{eq:HessianL} gives
$$
\nabla^2\phi^*(\nabla f(v))=\mathrm{diag}(c_P)^{-1}\preceq \nabla^2 f(v)^{-1}=\nabla^2 f^*(\nabla f(v)).
$$
Thus the dual relative smoothness inequality holds with constant $L=1$ \citep{maddison2021dual}; equivalently,
$$
D_{\phi^*}(\xi,\eta)\le D_{f^*}(\xi,\eta)
$$
for all admissible $\xi,\eta$. Integrating $\nabla^2\phi^*$ with $\nabla\phi^*(0)=0$ gives
$$
\nabla\phi^*(\xi)=\log(\mathbf{1}+\xi./b),
\qquad
\phi(v)=\sum_{j=1}^m b_j(e^{v_j}-v_j).
$$
Therefore, unit-step dual mirror descent gives
\begin{equation}\label{eq:md-general}
v^{k+1}=v^k-\nabla\phi^*(\nabla f(v^k))
=v^k+\log \Bigl(b./c_P(u^{k+1},v^k)\Bigr),
\end{equation}
which is exactly the Sinkhorn $v$-update \eqref{eq:v-update}.
This immediately implies the standard $\mathcal{O}(1/k)$ sublinear convergence of Sinkhorn \cite{maddison2021dual}. 

Linear convergence is more subtle. The shift invariance of $f$ follows from the corresponding invariance of the full dual objective \eqref{eq:Fc}:
$$
f(v+c\mathbf{1})
=
\min_u F(u,v+c\mathbf{1})
=
\min_u F(u+c\mathbf{1},v)
=
f(v).
$$
Hence $f$ is not strongly convex. We remove this degeneracy by the normalized Sinkhorn iteration
\begin{equation}\label{eq:normalized-sinkhorn}
v_0\in N^\perp,\qquad
v^{k+1} = \text{Sinkhorn}(v^k):=
v^k-P_{N^\perp}\nabla\phi^*\!\bigl(P_{N^\perp}\nabla f(v^k)\bigr).
\end{equation}
It differs from the plain Sinkhorn iteration only by an additive multiple of $\mathbf{1}$, and therefore gives the same primal update $P^k$. The normalized iterates are uniformly bounded \citep{carlier2022linear}.  After the normalization, one proves a Polyak--{\L}ojasiewicz inequality and obtain linear convergence; see Appendix~\ref{app:proof-Sinkhorn}.


%The following one-step estimate is the key descent bound used later.
%\begin{lemma}\label{lem:sinkhorn-descent}
%Let
%$
%v^+:=v-\nabla\phi^*(\nabla f(v))
%$
%be the Sinkhorn update. Then
%\begin{equation}\label{eq:sinkhorn-descent}
%f(v^+)-f(v)
%\le
%-D_{\phi^*}(0,\nabla f(v))
%-D_{\phi^*}(\nabla f(v^+),0).
%\end{equation}
%\end{lemma}
% \begin{proof}
% Let $g:=\nabla f(v)$ and $g^+:=\nabla f(v^+)$. Since $v-v^+=\nabla\phi^*(g)-\nabla\phi^*(0)$, expansion at $v^+$ and Bregman duality give
% $$
% f(v^+)-f(v)
% =
% \langle g^+,v^+-v\rangle-D_f(v,v^+)
% =
% -\langle g^+,\nabla\phi^*(g)-\nabla\phi^*(0)\rangle-D_{f^*}(g^+,g).
% $$
% By the three-point identity,
% $$
% -\langle g^+-0,\nabla\phi^*(g)-\nabla\phi^*(0)\rangle
% =
% -D_{\phi^*}(g^+,0) - D_{\phi^*}(0,g) + D_{\phi^*}(g^+,g).
% $$
% Using dual relative smoothness, $D_{\phi^*}(g^+,g)\le D_{f^*}(g^+,g)$, we obtain \eqref{eq:sinkhorn-descent}.
% %$$
% %f(v^+)-f(v)\le -D_{\phi^*}(\nabla f(v^+),0)-D_{\phi^*}(0,\nabla f(v)).
% %$$
% \end{proof}

%\LC{add reference}. 

% \LC{check the $v^k$ and $v_k$. be consistent on the subscript.}
%\begin{equation}\label{eq:md-vector}
%v^{k+1}=v^k-\log\!\Bigl(\mathbf{1}+\nabla f(v^k)./b\Bigr)=v^k-\log\!\Bigl(c_P\bigl(u^{k+1},v^k\bigr)./b\Bigr),
%\end{equation}
%which is exactly the Sinkhorn $v$-update with $u^{k+1}=u(v^k)$.

% \begin{proof}
% Since $f$ is dual relatively smooth with respect to $\phi$ with constant $L=1$, the descent inequality for dual mirror descent gives
% $$
% \phi^*(\nabla f(v^+))\le \phi^*(0)+f(v)-f(v^+)-D_{\phi^*}(0,\nabla f(v)).
% $$
% Because $\phi^*(0)=0$ and $\phi^*\ge 0$, we obtain
% $$
% 0\le \phi^*(\nabla f(v^+))\le f(v)-f(v^+)-D_{\phi^*}(0,\nabla f(v)),
% $$
% which implies \eqref{eq:sinkhorn-descent}.
% %$$
% %f(v^+)-f(v)\le -D_{\phi^*}(0,\nabla f(v)).
% %$$
% %This is exactly .
% \end{proof}

% \LC{change to a proof of K-L condition and from the sufficient decay to give a linear convergence. Treat this is as a the main theorem and Theorem 1 as a colloary which is secondary.}


% \begin{remark}
% Since $f$ is invariant under shifts along $N$, i.e.
% $
% f(v+c\mathbf{1})=f(v), \forall\,c\in\mathbb{R},
% $
% $f$ is not strongly convex on $\mathbb{R}^n$. It is well known that $f$ is strictly convex on $N^\perp$, and that the only flat direction of $f$ is the constant-shift mode $N$; see, e.g., \citet{carlier2022linear}. Hence $f$ admits a unique minimizer $v^*$ on $N^\perp$. We can restrict the domain to $N^\perp$ and obtain the corresponding \textit{normalized Sinkhorn} algorithm
% $$
% v_0\in N_{\perp},\qquad v_{k+1}=P_{N^\perp}(\mathrm{Sinkhorn}(v_k)) = v_k-P_{N^\perp}\nabla\phi^*\!\bigl(P_{N^\perp}\nabla f(v_k)\bigr).
% $$
% Noticibly, the iterates of the algorithm differ from the original Sinkhorn iterates only by a multiple of the vector $\mathbf{1}$, so they are equivalent in the primal space.
% After restricting to the orthogonal subspace, one obtains a Polyak--Łojasiewicz (P-L) inequality, which together with the sufficient decay lemma yields linear convergence of Sinkhorn; see Appendix~\ref{app:normalized-sinkhorn} for detailed proofs. 
% \end{remark}


%
%%$$
%%:=\{v\in\mathbb R^n:\langle v,\mathbf 1\rangle=0\}.
%%$$
%On this subspace, the only flat direction is removed. In particular, on every bounded sublevel set
%$$
%\mathcal S:=\{v\in N^\perp:\ f(v)\le f(v_0)\},
%$$
%the function $f$ satisfies the following mirror Polyak--{\L}ojasiewicz inequality.
%
%\begin{proposition}[Mirror Polyak--{\L}ojasiewicz inequality on $N^\perp$]\label{prop:PL}
%For every bounded sublevel set $\mathcal S\subset N^\perp$, there exists $0<\mu < 1$ such that
%\begin{equation}\label{eq:mirror-PL}
%f(v)-f(v^*)
%\le
%\frac1\mu D_{\phi^*}(0,\nabla f(v)),
%\qquad
%v\in\mathcal S.
%\end{equation}
%\end{proposition}
%
%
%\begin{lemma}[Lemma 3.1 \cite{carlier2022linear}]\label{lem:bounded}
%Let $\{v^k\}_{k\ge0}\subset N^\perp$ be generated by \eqref{eq:normalized-sinkhorn}. Then there exists $C>0$ such that
%$$
%\|v^k\|_\infty\le C,\qquad k\ge0.
%$$
%Consequently, the iterates remain in a compact sublevel set of $f$.
%\end{lemma}
%
%Combining Lemma~\ref{lem:sinkhorn-descent}, Proposition~\ref{prop:PL}, and Lemma~\ref{lem:bounded} gives linear convergence.
%
%\begin{theorem}[Linear convergence of Sinkhorn]\label{thm:sinkhorn-linear-conv}
%Let $\{v^k\}$ be generated by the normalized Sinkhorn iteration \eqref{eq:normalized-sinkhorn}. Then the linear convergence holds
%\begin{equation}\label{eq:sinkhorn-linear}
%f(v^k)-f(v^*)
%\le
%(1-\mu )^k\bigl(f(v^0)-f(v^*)\bigr).
%\end{equation}
%The same rate holds for the usual Sinkhorn iterates modulo constant shifts.
%\end{theorem}
%
%\begin{corollary}\label{thm:Sinkhorn}
%The Sinkhorn iteration satisfies
%$$
%f(v^k)-f(v^*)=\mathcal{O}(1/k).
%$$
%\end{corollary}


% \LC{Move the normalized Sinkhorn to appendix. Just leave one sentence remark here}.


%Then the $x$-update in \eqref{eq:scheme} becomes
%$$
%x^{k+1}-x^k=\alpha(y^k-x^{k+1})-\nabla\phi^*(\nabla f(x^k)),
%$$
%%and hence
%%$$
%%x^{k+1}=\frac{1}{1+\alpha}\Bigl(x^k+\alpha y^k-\nabla\phi^*(\nabla f(x^k))\Bigr)=\frac{1}{1+\alpha}\Bigl(\alpha y^k+\mathrm{Sinkhorn}(x^k)\Bigr).
%%$$
%%To simplify the notation, define
%%$$
%%Then the $x$-update reads
%$$
%x_{k+1}=\frac{1}{1+\alpha}\Bigl(w_k+\mathrm{Sinkhorn}(x_k)\Bigr).
%$$
%
%The $y$-update in \eqref{eq:scheme} gives
%$$
%(1+\alpha)y^{k+1}=y^k+\alpha x^{k+1}-\frac{\alpha}{\mu}\nabla\phi^*(\nabla f(x^{k+1})).
%$$
%Multiplying by $\alpha$ and using
%$$
%\nabla\phi^*(\nabla f(x^{k+1}))=x^{k+1}-\mathrm{Sinkhorn}(x^{k+1}),
%$$
%we obtain
%$$
%(1+\alpha)w_{k+1}=w_k+\alpha^2 x^{k+1}-\frac{\alpha^2}{\mu}\Bigl(x^{k+1}-\mathrm{Sinkhorn}(x^{k+1})\Bigr).
%$$
%Now set $\alpha=\sqrt{2\mu}$, so that $\alpha^2/\mu=2$. Then
%$$
%w_{k+1}=\frac{1}{1+\alpha}\Bigl(w_k+(\alpha^2-2)x^{k+1}+2\,\mathrm{Sinkhorn}(x^{k+1})\Bigr).
%$$

\section{Accelerating Sinkhorn}
\label{sec:accelerated}
Motivated by Hessian-driven Nesterov accelerated gradient (HNAG) \citep{chen2019orderoptimizationmethodsbased}, we propose the following accelerated dual mirror descent scheme for minimizing $f$:
\begin{equation}\label{eq:scheme}
\begin{aligned}
\frac{x_{k+1}-x_k}{\alpha} &= y_k-x_{k+1}-\frac{1}{\alpha} \nabla \phi^*(\nabla f(x_k)),\\
\frac{y_{k+1}-y_k}{\alpha} &= x_{k+1}-y_{k+1}-\frac{1}{\mu}\nabla\phi^*(\nabla f(x_{k+1})).
\end{aligned}
\end{equation}
Here $x_k$ is the main iterate, corresponding to $v^k$ in Sinkhorn, $y_k$ is an auxiliary iterate.

Following \citet{chen2025hnagsuperfastacceleratedgradient}, we set $\alpha=\sqrt{2\mu}$ and introduce $w_k:=\alpha y_k$. This gives the equivalent simplified updates in Algorithm~\ref{alg:A2MD}. We also give a practical homotopy version in Algorithm~\ref{alg:A2MD-homotopy}, where $\mu$ is decreased to $0$ by a prescribed schedule.

% Essentially, when $x_k, y_k\in N^\perp$, we have $x_{k+1}, y_{k+1}\in N^\perp$ too, so once we take $x_0\in N^\perp$, $w_0\in N^\perp$, Algorithm~\ref{alg:A2MD} is exactly operating in the subspace $N^\perp$, minimizing $f|_{P_{N^\perp}}$.

\begin{algorithm}[htbp]
\caption{Accelerated Dual Mirror Descent for Sinkhorn (Acc-Sinkhorn)}
\label{alg:A2MD}
\begin{algorithmic}[1]
\STATE \textbf{Input:} $x_0\in N^\perp$, $w_0\in N^\perp$, $0<\mu<1$, and $m\geq 1$
\STATE \textbf{Set} $\alpha=\sqrt{2\mu}$
\FOR{$k=0,1,\dots,m-1$}
    \STATE $x_{k+1}=\frac{1}{1+\alpha}\left(w_k+\text{Sinkhorn}(x_k)\right)$
    \STATE $w_{k+1}=\frac{1}{1+\alpha}\left(w_k+(\alpha^2-2)x_{k+1}+2\,\text{Sinkhorn}(x_{k+1})\right)$
\ENDFOR
\STATE \textbf{Output:} $(x_m,w_m)$
\end{algorithmic}
\end{algorithm}

The Sinkhorn step computed in the update of $w_{k+1}$ is reused in the next update of $x_{k+2}$. Therefore, each iteration requires only one normalized Sinkhorn step. The extra cost is negligible, since it consists only of vector additions and scalar multiplications.

\begin{algorithm}[htbp]
\caption{Acc-Sinkhorn with Homotopy}
\label{alg:A2MD-homotopy}
\begin{algorithmic}[1]
\STATE \textbf{Input:} $x_0\in\mathbb{R}^n$, $w_0\in\mathbb{R}^n$, $0<\mu_0<1$, $m_0\geq 1$, and maxIt
\FOR{$k=0,1,\dots,\text{maxIt}$}
    \STATE $(x_{k+1},w_{k+1})=\text{Acc-Sinkhorn}(x_k,w_k,\mu_k,m_k)$
    \STATE $\mu_{k+1}=\mu_k/2$, \quad $m_{k+1}=\lfloor\sqrt{2}\,m_k\rfloor+1$
\ENDFOR
\end{algorithmic}
\end{algorithm}

% \begin{lemma}
% The matrix $A(g)$ can be computed in closed form as
% \[A(g) = \mathrm{Diag}\Big(\frac{1}{g_j}\log\frac{b_j+g_j}{b_j}\Big)_{j=1}^n.\]    
% \end{lemma}
% \begin{proof}
%     Since $\nabla\phi^*$ is $C^1$ on its domain and $\nabla\phi^*(0)=0$, we have
% \[
% \nabla\phi^*(g)-\nabla\phi^*(0)
% =\int_0^1 \nabla^2\phi^*(t g)\,g\dd t
% =\Big(\int_0^1 \nabla^2\phi^*(t g)\dd t\Big)\,g.
% \]

% Since $\phi$ is separable, we can compute $\nabla^2\phi^*(g)$ in closed form as a diagonal matrix. For $v\in\mathbb{R}^n$,
% \[
% \phi(v)=\sum_{j=1}^n b_j\big(e^{v_j}-v_j\big),
% \qquad
% \nabla\phi^*(g)=\Big(\log\big(1+g_j/b_j\big)\Big)_{j=1}^n,
% \]
% hence
% \[
% \nabla^2\phi^*(g)=\mathrm{Diag}\Big(\frac{1}{b_j+g_j}\Big)_{j=1}^n.
% \]
% Then
% \[
% \int_0^1 \nabla^2\phi^*(t g)\dd t
% =
% \mathrm{Diag}\Big(\int_0^1 \frac{1}{b_j+t g_j}\dd t\Big)_{j=1}^n
% =
% \mathrm{Diag}\Big(\frac{1}{g_j}\log\frac{b_j+g_j}{b_j}\Big)_{j=1}^n,
% \]
% \end{proof}
% \begin{remark}
%     When $g_j\rightarrow 0$, the $j$th diagonal entry of $A(g)$ converges to $\frac{1}{b_j}$, which is exactly the $j$th diagonal entry of $\nabla^2\phi^*(0)$. This is consistent with the fact that $A(0)=\nabla^2\phi^*(0)$. So we can continuously extend $A(g)$ to $0$ by defining $A(0)=\nabla^2\phi^*(0)$.
% \end{remark}

\paragraph{Convergence Analysis.}
We analyze the convergence of \eqref{eq:scheme} through its continuous-time flow and discretization effects. Define the Lyapunov function
\begin{equation}
\mathcal{E}(x,y;\mu,D):=f(x)-f(x^\star)+\frac{\mu}{2}\|y-x^\star\|_{D(p(x))}^2,
\qquad
p(x):=\nabla\phi^*(\nabla f(x)),
\end{equation}
where $x^\star\in N^\perp$ is the optimal solution. The diagonal matrix $D(z)$ is
\begin{equation}\label{eq:D}
D(z):=\mathrm{diag}\!\left(b.*(\exp(z)-\mathbf{1})./z\right),
\end{equation}
where the quotient is entrywise and is understood by continuous extension at $z=0$, with $D(0)=\nabla^2\phi(0)$. Equivalently, $D(z)z=\nabla\phi(z)$.

\begin{lemma}\label{assump:boundedness}
There exists $R>0$ such that, for all iterates $(x_k,y_k)$ of \eqref{eq:scheme} with $\alpha$ sufficiently small,
\begin{equation}\label{eq:boundedness}
\|x_k-x^\star\|_{D(p(x_k))}\le R,\qquad \|y_k-x^\star\|\le R,\qquad k\ge 0.
\end{equation}
\end{lemma}
The proof is given in Appendix~\ref{app:boundedness-iterates}, using a different Lyapunov function.
%Assumption 1 is mild and easy to satisfy in practice. We refer to Appendix C for a detailed discussion on what conditions ensure this boundedness, and Appendix~\ref{app:boundedness-iterates} for a theoretical verification of this assumption. In short, $\alpha=c\sqrt{\mu}$ with $c\in[1,2)$ is sufficient to ensure the boundedness of iterates given $\mu$ is small enough. 
Although the boundedness can be established for $\alpha < \sqrt{2\mu}$, in practice, we choose the exact upper bound $\alpha=2\sqrt{\mu}$, and the iterates are always observed to be bounded in our experiments.
% \LC{move the discussion of Assumption into Appendix}
% The main result is the following.

%  Again, this is a mild assumption, and we refer to Appendix~\ref{app:boundedness-iterates} for a theoretical verification.

% \begin{assumption}\label{assumption:boundedness-y}
% There exists a constant $R_y>0$ such that for all iterates $y_k$, it holds $\|y_k-x^\star\|\le R_y$.
% \end{assumption}
% For simplicity, we take $R$ as a bound for both Assumption \ref{assump:boundedness} and \ref{assumption:boundedness-y} for simplicity.
Under Lemma~\ref{assump:boundedness}, our main result is as follows.
% \LC{What is $D$ in Theorem? Is it fixed or changing?}
\begin{theorem}
    \label{thm:homotopy}
Choose $(x_0,y_0)$ and $\mu_0$ such that
$$
\mathcal E(x_0,y_0;\mu_0,D(p(x_0)))\le (R^2+1)\mu_0.
$$
Let $(x_k,y_k,\mu_k)$ be generated by Algorithm~\ref{alg:A2MD-homotopy} and assume \eqref{eq:boundedness} holds. Then
\begin{equation}\label{eq:Ek-ek}
\mathcal E(x_k,y_k;\mu_k,D(p(x_k)))\le (R^2+1)\mu_k,
\qquad\forall\,k\ge 0.
\end{equation}
Moreover, let
$
M_k:=\sum_{i=0}^k m_i
$
be the total number of inner iterations after the $k$th outer loop, and
$
C_*:=\frac{\sqrt{2}-1}{(\sqrt{2L_F}+\sqrt{2\mu_0})\ln\bigl(2(R^2+1)\bigr)}
$
be a constant. Then
\begin{equation}% \label{eq:rate}
\mathcal E(x_k,y_k;\mu_k,D(p(x_k)))\le \frac{R^2+1}{\bigl(C_*M_k+\mu_0^{-1/2}\bigr)^2}
\qquad\forall\,k\ge 0.
\end{equation}
In particular, it takes $M_k=O(\tau^{-1/2})$ iterations to reach accuracy
$
\mathcal E(x_k,y_k;\mu_k,D(p(x_k)))\leq C \tau.
$
\end{theorem}

We give an outline of the proof in the rest of this section and leave details to Appendix~\ref{app:proofs-accelerated}.
\paragraph{Continuous-time ODE}
We first consider the continuous-time ODE of the scheme \eqref{eq:scheme}
\begin{equation}
\begin{aligned}
x' = y - x - \beta \nabla \phi^*(\nabla f(x)),\qquad 
y' = x - y - \frac{1}{\mu} \nabla\phi^*(\nabla f(x)),
\end{aligned}
\end{equation}
where $\beta>0$ is a parameter. The parameter $\mu$ is omited in the Lyapunov function when it is fixed and clear from the context.
%We will use the Lyapunov function
%\begin{equation}
%\mathcal{E}(x,y;\mu,D) = f(x)-f(x^\star) + \frac{\mu}{2}\|y-x^\star\|_{D(p(x))}^2.
%\end{equation}
%When $x_j\to 0$, the $j$th diagonal entry of $D(x)$ satisfies
%$$
%D_{jj}(x) = b_j\frac{e^{x_j}-1}{x_j} = b_j(x_j + x_j^2 ...)/x_j \to b_j = \partial_{jj}\phi(0).
%$$

%\begin{lemma}
%Let $\phi(x)=(b,e^x)-(b,x)$. Define $D(x)$ by $D(x)\,x=\nabla\phi(x)$. Then
%$$
%D(x)=\mathrm{diag}\!\left(b.\!*\frac{e^x-\mathbf{1}}{x}\right),
%$$
%where the fraction is entrywise (with the continuous extension at $x_j=0$). Consequently, with $p(x)=\nabla\phi^*(\nabla f(x))$,
%$$
%D(p(x))\,p(x)=\nabla f(x).
%$$
%\end{lemma}
%
%\begin{proof}
%Since $\nabla\phi$ is $C^1$ and $\nabla\phi(0)=0$, the fundamental theorem of calculus gives
%$$
%\nabla\phi(x)=\int_0^1 \nabla^2\phi(tx)\,x\dd t=\Bigl(\int_0^1 \nabla^2\phi(tx)\dd t\Bigr)x.
%$$
%Because $\phi$ is separable,
%$$
%\nabla\phi(x)=b.\!*\,e^x-b,
%\qquad
%\nabla^2\phi(x)=\mathrm{diag}\bigl(b.\!*\,e^x\bigr).
%$$
%Therefore,
%$$
%\int_0^1 \nabla^2\phi(tx)\dd t
%=\mathrm{diag}\!\Bigl(\int_0^1 b.\!*\,e^{tx}\dd t\Bigr)
%=\mathrm{diag}\!\left(b.\!*\frac{e^x-\mathbf{1}}{x}\right).
%$$
%Defining $D(x):=\int_0^1 \nabla^2\phi(tx)\dd t$ yields $D(x)\,x=\nabla\phi(x)$ and the stated closed form.
%
%For the second claim, assume $\nabla f(x)>-b$ so that $p(x)=\nabla\phi^*(\nabla f(x))$ is well defined. Then
%$$
%D(p(x))\,p(x)=\nabla\phi(p(x))=\nabla\phi\bigl(\nabla\phi^*(\nabla f(x))\bigr)=\nabla f(x).
%$$
%\end{proof}

\begin{comment}
\textbf{Added a version for $\psi$ to replace the above lemma.}
\begin{proposition}[Diagonal Representation for $\psi$]\label{prop:psi-diagonal}
Let $b \in \mathbb{R}^n_{>0}$, $N = \mathrm{span}(\mathbf{1})$, and let $\psi: N^\perp \to \mathbb{R}$ be defined by
\[
\psi(x_\perp) := \min_{\alpha \in \mathbb{R}} \sum_{j=1}^n b_j \big( e^{x_\perp^j + \alpha} - (x_\perp^j + \alpha) \big), \quad x_\perp \in N^\perp.
\]

Define the diagonal operator $D_\perp(x_\perp) \in \mathbb{R}^{n \times n}$ acting on $N^\perp$ by
\[
D_\perp(x_\perp)\, x_\perp := \nabla \psi(x_\perp),
\]
with the fraction interpreted entrywise as the continuous extension at zero. Then
\[
D_\perp(x_\perp) = \mathrm{diag} \left(\left(\frac{B b.*e^{x_\perp}}{\sum_j b_j e^{x_\perp^j}}-b\right)./{x_\perp}\right),  \quad B = \sum_j b_j.
\]

Consequently, if $p_\perp(x) := \nabla \psi^*(\nabla f(x))$ is well-defined for $\nabla f(x) \in \mathrm{range}(\nabla \psi)$, then
\[
D_\perp(p_\perp(x))\, p_\perp(x) = \nabla f(x).
\]
\end{proposition}

\begin{proof}
Recall from Proposition \ref{prop:psi} that $$\nabla \psi(x_\perp) = B\frac{b .* e^{x_\perp}}{\sum_j b_j e^{x_\perp^j}} - b\in N^\perp.$$ Therefore, the diagonal operator $D_\perp(x_\perp)$ is given by
\[
D_\perp(x_\perp) = \mathrm{diag} \left(\left(\frac{B b.*e^{x_\perp}}{\sum_j b_j e^{x_\perp^j}}-b\right)./{x_\perp}\right),
\]

Let $p_\perp(x) := \nabla \psi^*(\nabla f(x))$, assuming it is well-defined. Then
\[
\nabla \psi(p_\perp(x)) = \nabla f(x) \implies D_\perp(p_\perp(x))\, p_\perp(x) = \nabla f(x),
\]
by the same argument as for $\phi$ using the fundamental theorem of calculus.
\end{proof}
\end{comment}

Let $z=(x,y)$ and let $\mathcal{G}(z)$ denote the right-hand side of the ODE. We have the following identity for the time derivative of the Lyapunov function along the trajectory of the ODE. The decay is exponential up to positive perturbation terms.

\begin{lemma}
Let $z(t)=(x(t),y(t))$ be a trajectory of $z'=\mathcal{G}(z)$. Define
$\mathcal D(t):=D(p(x(t))),$
where $D(\cdot)$ is the diagonal map \eqref{eq:D}. Then, for all $t\ge 0$, the following identity holds:
$$
\begin{aligned}
\frac{\dd}{\dd t}\mathcal{E}(x(t),y(t);\mathcal D(t))
%={}&\langle \nabla \mathcal{E}(x,y;\mathcal D),\mathcal{G}(z)\rangle\\
={}& -\mathcal{E}(x,y;\mathcal D) - D_f(x^\star,x) -\beta\|p(x)\|_{\mathcal D(t)}^2 -\frac{\mu}{2}\|x-y\|_{\mathcal D(t)}^2
\\
&\quad + \frac{\mu}{2}\|x-x^\star\|_{\mathcal D(t)}^2+\frac{\mu}{2}\|y-x^\star\|_{\mathcal D'(t)}^2.
\end{aligned}
$$
\end{lemma}

% \begin{proof}
% Since $D=\mathcal D(t)$ varies with $t$, differentiating the Lyapunov function along the trajectory gives
% $$
% \mathcal{E}'(x,y;\mathcal D)
% =\langle \nabla f(x),x'\rangle+\mu\langle y-x^\star,y'\rangle_{\mathcal D(t)}+\frac{\mu}{2}\|y-x^\star\|_{\mathcal D'(t)}^2.
% $$
% Substituting $x'=y-x-\beta p(x)$ and $y'=x-y-\tfrac{1}{\mu}p(x)$, we obtain
% $$
% \begin{aligned}
% {}& \langle \nabla f(x),y-x-\beta p(x)\rangle + \mu\langle y-x^\star,x-y-\tfrac{1}{\mu}p(x)\rangle_{\mathcal D(t)}
% \\
% ={}& \langle \nabla f(x),y-x\rangle-\beta\langle \nabla f(x),p(x)\rangle+\mu\langle y-x^\star,x-y\rangle_{\mathcal D(t)}
% -\langle y-x^\star, \mathcal D(t)p(x)\rangle.
% \end{aligned}
% $$
% By definition of $D(\cdot)$ and $p(x)=\nabla\phi^*(\nabla f(x))$, we have
% $$
% \mathcal D(t)p(x(t))=D(p(x(t)))\,p(x(t)))=\nabla\phi(p(x(t)))=\nabla\phi\bigl(\nabla\phi^*(\nabla f(x(t)))\bigr)=\nabla f(x(t)),
% $$
% and hence
% $$
% \langle \nabla f(x),y-x\rangle-\beta\langle \nabla f(x),p(x)\rangle-\langle y-x^\star, \mathcal D(t)p(x)\rangle
% =-\beta\|p(x)\|_{\mathcal D(t)}^2-\langle \nabla f(x),x-x^\star\rangle.
% $$
% Moreover, the polarization identity yields
% $$
% \mu\langle y-x^\star,x-y\rangle_{\mathcal D(t)}
% =-\frac{\mu}{2}\|y-x^\star\|_{\mathcal D(t)}^2-\frac{\mu}{2}\|x-y\|_{\mathcal D(t)}^2+\frac{\mu}{2}\|x-x^\star\|_{\mathcal D(t)}^2.
% $$
% Using
% $$
% \langle \nabla f(x),x-x^\star\rangle = f(x)-f(x^\star) + D_f(x^\star,x),
% $$
% and recalling $\mathcal{E}(x,y;\mathcal D)=f(x)-f(x^\star)+\frac{\mu}{2}\|y-x^\star\|_{\mathcal D(t)}^2$, we conclude the claim.
% \end{proof}

\paragraph{Discretization and convergence analysis}

%The continuous-time dynamics show that the Lyapunov function $\mathcal{E}$ decays exponentially, up to a positive perturbation term $\frac{\mu}{2}\|x-x^\star\|_{D}^2$ and a negative term $\beta\|p(x)\|_{D}^2$. To control the positive perturbation term, we use Lemma~\ref{assump:boundedness}. 
%The metric also changes with the iterate, producing the additional term $\frac{\mu}{2}\|y-x^\star\|_{\mathcal D_{k+1}-\mathcal D_k}^2$ in the Lyapunov difference. 
We define this discrete metric by
$$
p_k:=\nabla\phi^*(\nabla f(x_k)),
\qquad
\mathcal D_k:=D(p_k),
$$
where $D(\cdot)$ is the diagonal map defined in \eqref{eq:D}. 
%Thus $\mathcal D_k$ is the metric used in the Lyapunov function at the iterate $x_k$.

\begin{lemma}
\label{thm:single-step}
Let $z_k=(x_k,y_k)$ be the iterates generated by \eqref{eq:scheme} and assume \eqref{eq:boundedness} holds. Then
$$
\begin{aligned}
\mathcal{E}(z_{k+1};\mathcal D_{k+1}) - \mathcal{E}(z_k;\mathcal D_{k+1})
\le{}& -\alpha \mathcal{E}(z_{k+1};\mathcal D_{k+1}) + \frac{\alpha\mu}{2}R^2 - D_{\phi^*}(0,\nabla f(x_k))
\\
&\quad + \frac{\alpha^2}{2\mu}\|\nabla \phi^*(\nabla f(x_{k+1}))\|_{\mathcal D_{k+1}}^2 - D_{\phi^*}(\nabla f(x_{k+1}),0).
\end{aligned}
$$
\end{lemma}


% \begin{proof}
% First, consider the implicit Euler scheme
% \begin{equation}
% \label{eq:implicit-euler}
% \frac{z_{k+1}-z_k}{\alpha}=\mathcal{G}(z_{k+1}).
% \end{equation}
% Expanding the Lyapunov difference at $z=z_{k+1}$ yields
% $$
% \begin{aligned}
% \mathcal{E}(z_{k+1};\mathcal D_{k+1}) - \mathcal{E}(z_k;\mathcal D_{k+1})
% ={}& \langle \nabla \mathcal{E}(z_{k+1};\mathcal D_{k+1}), z_{k+1}-z_k\rangle - D_{\mathcal{E}}(z_{k},z_{k+1};\mathcal D_{k+1})
% \\ 
% ={}& \alpha \langle \nabla \mathcal{E}(z_{k+1};\mathcal D_{k+1}), \mathcal{G}(z_{k+1})\rangle - D_{f}(x_{k},x_{k+1}) - \frac{\mu}{2}\|y_{k+1}-y_k\|_{\mathcal D_{k+1}}^2.
% \end{aligned}
% $$
% Comparing \eqref{eq:implicit-euler} with the actual scheme \eqref{eq:scheme}, the $x$-update differs by two terms. This produces additional error terms in the Lyapunov difference:
% $$
% \begin{aligned}
% \mathcal{E}(z_{k+1};\mathcal D_{k+1}) - \mathcal{E}(z_k;\mathcal D_{k+1})
% \leq{}& \alpha \langle \nabla \mathcal{E}(z_{k+1};\mathcal D_{k+1}), \mathcal{G}(z_{k+1})\rangle - D_{f}(x_{k},x_{k+1}) - \frac{\mu}{2}\|y_{k+1}-y_k\|_{\mathcal D_{k+1}}^2
% \\
% &\quad + \alpha\langle \nabla f(x_{k+1}), y_k-y_{k+1}\rangle
% + \alpha\beta\langle \nabla f(x_{k+1}), \nabla \phi^*(\nabla f(x_{k+1})) - \nabla \phi^*(\nabla f(x_{k}))\rangle.
% \end{aligned}
% $$
% By the continuous-time analysis,
% $$
% \langle \nabla \mathcal{E}(z_{k+1};\mathcal D_{k+1}), \mathcal{G}(z_{k+1})\rangle \leq - \mathcal{E}(z_{k+1};\mathcal D_{k+1}) +\frac{\mu}{2}R^2-\beta\|\nabla\phi^*(\nabla f(x_{k+1}))\|_{\mathcal D_{k+1}}^2.
% $$
% For the first cross term, the Cauchy--Schwarz and Young inequalities give
% $$
% \begin{aligned}
% \alpha\langle \nabla f(x_{k+1}), y_k-y_{k+1}\rangle
% \leq{}& \frac{\alpha}{\sqrt{\mu}}\|\nabla f(x_{k+1})\|_{\mathcal D_{k+1}^{-1}} \cdot \sqrt{\mu} \|y_k-y_{k+1}\|_{\mathcal D_{k+1}}
% \\
% \leq{}& \frac{\alpha^2}{2\mu}\|\nabla f(x_{k+1})\|_{\mathcal D_{k+1}^{-1}}^2 + \frac{\mu}{2}\|y_k-y_{k+1}\|_{\mathcal D_{k+1}}^2.
% \end{aligned}
% $$
% Since $\mathcal D_{k+1}\nabla\phi^*(\nabla f(x_{k+1}))=\nabla\phi(\nabla\phi^*(\nabla f(x_{k+1})))=\nabla f(x_{k+1})$, we have
% $$
% \|\nabla f(x_{k+1})\|_{\mathcal D_{k+1}^{-1}}^2 = \|\nabla \phi^*(\nabla f(x_{k+1}))\|_{\mathcal D_{k+1}}^2.
% $$
% For the second cross term, the three-point identity for the Bregman divergence yields
% $$
% \begin{aligned}
% &\langle \nabla f(x_{k+1}), \nabla \phi^*(\nabla f(x_{k+1})) - \nabla \phi^*(\nabla f(x_{k}))\rangle
% \\
% ={}& D_{\phi^*}(\nabla f(x_{k+1}),\nabla f(x_{k})) - D_{\phi^*}(\nabla f(x_{k+1}),0) - D_{\phi^*}(0,\nabla f(x_{k})).
% \end{aligned}
% $$
% Using $\phi^*(0)=0$ and the definition of the Bregman divergence,
% $$
% D_{\phi^*}(\nabla f(x_{k+1}),0)=\phi^*(\nabla f(x_{k+1}))-\langle \nabla f(x_{k+1}),0\rangle-\phi^*(0)=\phi^*(\nabla f(x_{k+1})).
% $$
% Combining the above bounds gives
% $$
% \begin{aligned}
% \mathcal{E}(z_{k+1};\mathcal D_{k+1}) - \mathcal{E}(z_k;\mathcal D_{k+1})
% \leq{}& -\alpha \mathcal{E}(z_{k+1};\mathcal D_{k+1}) + \frac{\alpha\mu}{2}R^2  - \alpha\beta D_{\phi^*}(0,\nabla f(x_{k}))
% \\
% &\quad + \frac{\alpha^2}{2\mu}\|\nabla \phi^*(\nabla f(x_{k+1}))\|_{\mathcal D_{k+1}}^2 - \alpha\beta D_{\phi^*}(\nabla f(x_{k+1}),0)
% \\
% &\quad +\alpha\beta D_{\phi^*}(\nabla f(x_{k+1}),\nabla f(x_{k})) - D_{f}(x_{k},x_{k+1}).
% \end{aligned}
% $$
% By dual relative smoothness, $ D_{\phi^*}(\nabla f(x_{k+1}),\nabla f(x_{k})) \leq D_f(x_{k},x_{k+1})$. Therefore, it suffices to choose $\alpha$ and $\beta$ such that $\alpha\beta=1$.
% \end{proof}




% In the above Lyapunov difference, there is one positive perturbation term and two negative perturbation terms
% $$
% \frac{\alpha^2}{2\mu}\|\nabla \phi^*(\nabla f(x_{k+1}))\|_{\mathcal D_{k+1}}^2
% \qquad\text{and}\qquad
% - D_{\phi^*}(\nabla f(x_{k+1}),0) - D_{\phi^*}(0,\nabla f(x_{k})).
% $$
% We use the negative terms to offset the positive term. 
% The following lemma makes the relation between these two terms explicit.

% \begin{lemma}
% For any $k\ge 1$, it holds that
% $$
% \frac{1}{2}\|\nabla \phi^*(\nabla f(x_{k}))\|_{D_{k}}^2
% =\frac{1}{2}\mathrm{KL}(c_{P_k}\|b)+\frac{1}{2}\mathrm{KL}(b\|c_{P_k}),
% $$
% and
% $$
% D_{\phi^*}(\nabla f(x_{k}),0)=\mathrm{KL}(c_{P_k}\|b),
% $$
% where $c_{P_k}:=P_k^\top\mathbf{1}_n$ is the column-sum vector of $P_k$, and $P_k$ is the optimal transport matrix computed by Sinkhorn at $x_k$.
% \end{lemma}

% \begin{proof}
% Let $g_k:=\nabla f(x_k)=c_{P_k}-b$. Since $\nabla \phi^*(\eta)=\log\!\bigl(\mathbf{1}+\eta./b\bigr)$, we have
% $$
% \nabla \phi^*(g_k)=\log(c_{P_k}./b).
% $$
% By the definition of $\mathcal D_k$ and $x^\phi=\nabla\phi^*(\nabla f(x))$, we have
% $$
% \mathcal D_k\,\nabla \phi^*(g_k)=g_k.
% $$
% Therefore,
% $$
% \|\nabla \phi^*(g_k)\|_{\mathcal D_k}^2=\big\langle \nabla \phi^*(g_k),\,g_k\big\rangle.
% $$
% Using the Bregman three-point identity with $(\eta,\xi)=(g_k,0)$,
% $$
% D_{\phi^*}(g_k,0)+D_{\phi^*}(0,g_k)
% =\big\langle \nabla\phi^*(g_k)-\nabla\phi^*(0),\,g_k-0\big\rangle
% =\big\langle \nabla\phi^*(g_k),\,g_k\big\rangle,
% $$
% since $\nabla\phi^*(0)=0$. Hence
% $$
% \frac{1}{2}\|\nabla \phi^*(\nabla f(x_k))\|_{\mathcal D_k}^2
% =\frac{1}{2}\Big(D_{\phi^*}(g_k,0)+D_{\phi^*}(0,g_k)\Big).
% $$
% By the identification $g_k=c_{P_k}-b$ and equation \eqref{eq:Dphistar-KL-special}, we have
% $$
% D_{\phi^*}(g_k,0)=\mathrm{KL}(c_{P_k}\|b),
% \qquad
% D_{\phi^*}(0,g_k)=\mathrm{KL}(b\|c_{P_k}),
% $$
% which gives
% $$
% \frac{1}{2}\|\nabla \phi^*(\nabla f(x_k))\|_{\mathcal D_k}^2
% =\frac{1}{2}\mathrm{KL}(c_{P_k}\|b)+\frac{1}{2}\mathrm{KL}(b\|c_{P_k}).
% $$
% \end{proof}

% To ensure
% $$
% \frac{\alpha^2}{2\mu}\|\nabla \phi^*(\nabla f(x_{k}))\|_{D_{k}}^2 -  D_{\phi^*}(\nabla f(x_{k}),0) \leq 0,
% $$
% it suffices to choose $\alpha$ such that
% $$
% \alpha=\sqrt{\mu \frac{\mathrm{KL}(c_{P_k}\|b)+\mathrm{KL}(b\|c_{P_k})}{2\,\mathrm{KL}(c_{P_k}\|b)}}.
% $$
% The ratio
% $$
% \frac{\mathrm{KL}(c_{P_k}\|b)+\mathrm{KL}(b\|c_{P_k})}{2\,\mathrm{KL}(c_{P_k}\|b)}
% $$
% measures the relative asymmetry of the KL divergence. In practice, we set $\alpha=\sqrt{\mu}$, which amounts to ignoring this asymmetry. When the marginal deviation of $c_{P_k}$ from $b$ is small, the asymmetry is also small, so $\alpha=\sqrt{\mu}$ is a good approximation.

% \begin{remark}
%     In the above analysis, the negative term $-D_{\phi^*}(0,\nabla f(x_k))$ is not used. Inspired by \citet{chen2025hnagsuperfastacceleratedgradient}, we can enlarge the step size $\alpha$, and offset the extra positivity using this term. More precisely, we want to choose $\alpha$ such that
%     $$\frac{\alpha^2}{2\mu}\|\nabla \phi^*(\nabla f(x_{k}))\|_{D_{k}}^2 -  D_{\phi^*}(\nabla f(x_{k}),0) \leq (1+\alpha) D_{\phi^*}(0,\nabla f(x_{k})).$$
%     Omit the extra positivity $\alpha D_{\phi^*}(0,\nabla f(x_{k}))$ on the right-hand side, we can choose 
%     $$\alpha=\sqrt{2\mu}\sqrt{\frac{D_{\phi^*}(0,\nabla f(x_{k}))+D_{\phi^*}(\nabla f(x_{k}),0)}{\|\nabla \phi^*(\nabla f(x_{k}))\|_{D_{k}}^2}}=\sqrt{2\mu},$$
%     where the last equality is due to the KL formulation of the three terms.
%     This gives a cleaner step size of $\alpha$.
    
% \end{remark}

% \begin{theorem}[Local bound of $\alpha$ in terms of marginal deviations]\label{thm:alpha-bound}
% Let $P_k$ be the optimal transport matrix computed by Sinkhorn at iterate $x_k$, with column sum $c_{P_k}:=P_k^\top \mathbf{1}_n$ and target marginal $b>0$. Define
% $$
% g_k:=c_{P_k}-b,\qquad
% \eta:=\left\|g_k./b\right\|_\infty\le \tfrac{1}{2},
% $$
% and set
% $$
% \alpha:=\sqrt{\mu\,\frac{\mathrm{KL}(c_{P_k}\|b)+\mathrm{KL}(b\|c_{P_k})}{2\,\mathrm{KL}(c_{P_k}\|b)}}.
% $$
% Then
% $$
% \sqrt{\mu}\!\left(1-\frac{\eta}{8}\right)\le \alpha \le \sqrt{\mu}\!\left(1+\frac{\eta}{8}\right).
% $$
% \end{theorem}

% \begin{proof}
% Write $c_{P_k,j}=b_j(1+t_j)$ with $t_j:=(g_k)_j/b_j$ and $\eta=\max_j|t_j|\le \tfrac12$. Define
% $$
% A:=\sum_j b_j t_j^2>0,\qquad B:=\sum_j b_j t_j^3.
% $$

% \noindent\textbf{KL expansions.}
% Using $\log(1+t)=t-\tfrac12 t^2+\tfrac13 t^3-\cdots$ and $\sum_j b_j t_j=(\mathbf{1},g_k)=0$, we obtain
% $$
% \mathrm{KL}(c_{P_k}\|b)=\tfrac12 A-\tfrac16 B+O(\eta^2 A),
% \qquad
% \mathrm{KL}(b\|c_{P_k})=\tfrac12 A-\tfrac13 B+O(\eta^2 A).
% $$

% \noindent\textbf{Ratio.}
% $$
% \frac{\mathrm{KL}(c_{P_k}\|b)+\mathrm{KL}(b\|c_{P_k})}{2\,\mathrm{KL}(c_{P_k}\|b)}
% =\frac{A-\tfrac12 B}{A-\tfrac13 B}+O(\eta^2)
% =1-\frac{\tfrac16 B}{A-\tfrac13 B}+O(\eta^2).
% $$

% \noindent\textbf{Bounding the correction.}
% Since $|B|\le \eta A$ and $A-\tfrac13|B|\ge A(1-\eta/3)>0$, we have
% $$
% \left|\frac{\tfrac16 B}{A-\tfrac13 B}\right|
% \le \frac{\eta/6}{1-\eta/3}\le \frac{\eta}{4}
% \qquad (\eta\le\tfrac12),
% $$
% and therefore
% $$
% 1-\frac{\eta}{4}\le \frac{\mathrm{KL}(c_{P_k}\|b)+\mathrm{KL}(b\|c_{P_k})}{2\,\mathrm{KL}(c_{P_k}\|b)}\le 1+\frac{\eta}{4}.
% $$

% \noindent\textbf{Conclusion.}
% Applying $\sqrt{1-x}\ge 1-\tfrac12 x$ and $\sqrt{1+x}\le 1+\tfrac12 x$ for $x\in[0,1]$, and then multiplying by $\sqrt{\mu}$, yields
% $$
% \sqrt{\mu}\!\left(1-\frac{\eta}{8}\right)\le \alpha \le \sqrt{\mu}\!\left(1+\frac{\eta}{8}\right). \qedhere
% $$
% \end{proof}

% To summarize, when the marginal deviation $\eta$ is small, we can choose $\alpha$ close to $\sqrt{\mu}$ so that the negative term $- D_{\phi^*}(\nabla f(x_{k+1}),0)$ offsets the positive term $\frac{\alpha^2}{2\mu}\|\nabla \phi^*(\nabla f(x_{k+1}))\|_{\mathcal D_{k+1}}^2$. This also explains the common practice in Sinkhorn-type methods of assuming $\min_j b_j>0$: a strictly positive target marginal limits the relative deviation $g_k./b$ and thus permits a stable step size as in Theorem~\ref{thm:alpha-bound}. We formalize this in the following assumption.

% \begin{assumption}[Marginal deviation]\label{assump:small-deviation}
% There exists a constant $\eta_0>0$ such that for all iterates $x_k$,
% $$
% \eta_k:=\left\|(c_{P_k}-b)./b\right\|_\infty\le \eta_0,
% $$
% where $c_{P_k}:=P_k^\top\mathbf{1}_n$ and $P_k$ is the optimal transport matrix computed by Sinkhorn at $x_k$. Moreover, $\eta_0$ is small enough so that with $\alpha=\sqrt{\mu}\left(1-\frac{\eta_0}{8}\right)$,
% $$
% \frac{\alpha^2}{2\mu}\|\nabla \phi^*(\nabla f(x_{k+1}))\|_{\mathcal D_{k+1}}^2- D_{\phi^*}(\nabla f(x_{k+1}),0)\le 0
% \qquad \text{for all } k.
% $$
% \end{assumption}

% Assumption~\ref{assump:small-deviation} is only used in the analysis. In practice, we simply set $\alpha=\sqrt{\mu}$, which works well in our experiments.

We next bound the change of the Lyapunov function induced by the change of the metric $D$.

\begin{lemma}
    \label{lem:chenge-D}
Let $z_k=(x_k,y_k)$ be the iterates generated by \eqref{eq:scheme} and assume \eqref{eq:boundedness} holds. Then there exists a constant $C>0$ such that for any $k\ge 1$,
$$
\mathcal{E}(z_k;\mathcal D_{k+1})-\mathcal{E}(z_k;\mathcal D_k)
=\frac{\mu}{2}\|y_k-x^\star\|_{\mathcal D_{k+1}-\mathcal D_k}^2
\le C\,\frac{\mu}{2}R^2.
$$
\end{lemma}

% \begin{proof}
% By the definition of $\mathcal{E}$,
% $$
% \mathcal{E}(z_k;\mathcal D_{k+1})-\mathcal{E}(z_k;\mathcal D_k)
% =\frac{\mu}{2}\|y_k-x^\star\|_{\mathcal D_{k+1}-\mathcal D_k}^2.
% $$
% Since $\mathcal D_k$ is diagonal, we have
% $$
% \|y_k-x^\star\|_{\mathcal D_{k+1}-\mathcal D_k}^2\le \|y_k-x^\star\|^2\,\|\mathcal D_{k+1}-\mathcal D_k\|.
% $$
% It remains to bound $\|\mathcal D_{k+1}-\mathcal D_k\|$. Recall that $D(s)\,s=\nabla\phi(s)$ and $\nabla\phi(s)=b.\!*\,\exp(s)-b$. For $s\in\mathbb{R}^m$,
% $$
% D(s)=\mathrm{diag}\!\left(b.\!*\frac{\exp(s)-\mathbf{1}}{s}\right),
% $$
% where the fraction is entrywise (with the continuous extension at $s_j=0$). Moreover, with $s=p(x)=\nabla\phi^*(\nabla f(x))$ and $\nabla\phi^*(\eta)=\log(\mathbf{1}+\eta./b)$, we have
% $$
% s=\log(c_P./b),
% \qquad
% \frac{\exp(s)-\mathbf{1}}{s}=\frac{c_P./b-\mathbf{1}}{\log(c_P./b)}.
% $$
% Define $h(g):=g./\log(g+\mathbf{1})$ for $g>0$ (componentwise). Then
% $$
% \mathcal D_k=\mathrm{diag}\!\bigl(b.\!*\,h(c_{P_k}./b-\mathbf{1})\bigr).
% $$
% Under Assumption~\ref{assump:boundedness}, the \textit{marginal deviation} $c_{P_k}./b-1$ is componentwise uniformly bounded, i.e., there exists $M>m>0$ s.t. $c_{P_k}./b-1\in[m,M]$. Since $h$ is Lipschitz on any compact domain, there exists $L=L(m,M)>0$ such that
% $$
% \|\mathcal D_{k+1}-\mathcal D_k\|
% \le L\,\max_j b_j\,\big\|c_{P_{k+1}}./b-c_{P_k}./b\big\|_\infty
% \le 2L(M-m),
% $$
% where the last inequality uses $c_{P_k}./b\-\mathbf{1}\in[m,M]$ for all $k$. Therefore,
% $$
% \begin{aligned}
% \mathcal{E}(z_k;\mathcal D_{k+1})-\mathcal{E}(z_k;\mathcal D_k)
% &\le \frac{\mu}{2}\|y_k-x^\star\|^2\,\|\mathcal D_{k+1}-\mathcal D_k\|\\
% &\le \frac{\mu}{2}R^2\cdot 2L(M-m).
% \end{aligned}
% $$
% Setting $C:=2L(M-m)$ gives the claim.
% \end{proof}

% \begin{theorem}[Boundedness of Energy Increment with $D_\perp$]
% \label{thm:chenge-D-perp}
% Assume that Assumption~\ref{assump:boundedness} and Assumption~\ref{assumption:boundedness-y} hold. Let $D_\perp(x_\perp)$ be the diagonal matrix associated with $\psi$ as in Proposition~\ref{prop:diag-psi}. Then there exists a constant $C>0$ such that for any $k \ge 1$,
% \[
% \mathcal{E}(z_k;D_{\perp,k+1}) - \mathcal{E}(z_k;D_{\perp,k})
% = \frac{\mu}{2} \|y_k - x^\star\|_{D_{\perp,k+1}-D_{\perp,k}}^2
% \le C \, \frac{\mu}{2} R^2.
% \]
% \end{theorem}

% \begin{proof}
% By the definition of $\mathcal{E}$, we have
% \[
% \mathcal{E}(z_k;D_{\perp,k+1}) - \mathcal{E}(z_k;D_{\perp,k}) 
% = \frac{\mu}{2} \|y_k - x^\star\|_{D_{\perp,k+1}-D_{\perp,k}}^2.
% \]

% Since $D_{\perp,k}$ is diagonal in the coordinates of $N^\perp$, the standard norm inequality gives
% \[
% \|y_k - x^\star\|_{D_{\perp,k+1}-D_{\perp,k}}^2
% \le \|y_k - x^\star\|^2 \, \|D_{\perp,k+1} - D_{\perp,k}\|.
% \]

% It remains to bound $\|D_{\perp,k+1}-D_{\perp,k}\|$. Recall from Proposition~\ref{prop:psi-diagonal} that
% \[
% D_\perp(x_\perp) = \mathrm{diag} \left(\left(\frac{B b.*e^{x_\perp}}{\sum_j b_j e^{x_\perp^j}}-b\right)./{x_\perp}\right),  \quad B = \sum_j b_j.
% \]

% Moreover, with $x_\perp = \nabla \psi^*(\nabla f(x))$ and $\nabla \psi^*(\eta) = P_{N^\perp}\nabla\phi^*(\eta)=P_{N^\perp}(\log(1+\eta./b))$, we have
% \[
% x_\perp = P_{N^\perp} \log(c_P./b), \quad D_\perp(x_\perp) = \mathrm{diag} \left(\frac{B\frac{c_P}{\sum_j c_P^j}-b}{P_{N^\perp}(\log(c_P./b))}\right).
% \]
% Under some Lipschitz condition, we can bound this term. XXX To be checked.XXX

% \end{proof}


% \begin{remark}
% Following Theorem~\ref{thm:alpha-bound}, one may take $\eta_0=\tfrac12$. A direct calculation gives $L(\eta_0)=(1-\log 2)/\log^2 2$. Therefore, we can take $C=(1-\log 2)/\log^2 2$ in the above lemma.
% \end{remark}

Combining the above results, we conclude that the Lyapunov function decays geometrically up to a bounded perturbation term of order $\mu R^2$.

\begin{theorem}[Perturbed exponential decay of the Lyapunov function]\label{coro:perturbed-decay}
Let $z_k=(x_k,y_k)$ be the iterates generated by \eqref{eq:scheme} and assume \eqref{eq:boundedness} holds. Then there exists a constant $C>0$ such that for any $k\ge 1$,
$$
\mathcal{E}(x_{k+1},y_{k+1};\mathcal D_{k+1}) \le \left(\frac{1}{1+\alpha}\right)^{k+1}\mathcal{E}(x_0,y_0;D_0) + C\mu R^2.
$$
\end{theorem}

Then using homotopy argument, we can obtain Theorem \ref{thm:homotopy}. 
% Finally, we can analyze the convergence of Algorithm \ref{alg:A2MD-homotopy}. We adopt the notation $$\mathcal{E}(x,y,\mu;D):=f(x)-f(x^\star)+\frac{\mu}{2}\|y-x\|_D^2$$ to explicitly indicate the dependence of $\mathcal{E}$ on $\mu$.

% \begin{theorem}
% 		Choose $(x_0, y_0)$ and $\mu_0$ satisfying $\mathcal E(x_0, y_0, \mu_0) \leq (R^2+ 1)\mu_{0}$.
% 		Let $(x_{k}, y_{k}, \mu_{k})$ be generated by Algorithm \ref{alg:A2MD-homotopy}. Assume that assumption \ref{assump:boundedness} holds. Then we have
% 		\begin{equation}\label{eq:Ek-ek}
% 			\mathcal E(x_{k}, y_{k}, \mu_{k})\leq (R^2+1)\mu_{k}\quad\forall\,k\geq0,
% 		\end{equation}
% Moreover, let $M_k:=\sum_{i=0}^{k}m_i $ be the overall iteration steps after the $k$-th outer iteration and $C = \frac{(\sqrt{2}-1)}{(\sqrt{2L_F}+\sqrt{2\mu_0})\ln (2(R^2+1))}$, then we have the global convergence rate
% \begin{equation}\label{eq:rate}
% 	\mathcal E(x_k,y_k,\mu_k)\leq \frac{R^2+1}{\left (C M_k+\mu_0^{-1/2}\right )^{2}}\quad\forall\,k\geq 0.
% \end{equation}
% 		Thus, it takes $M_k=O(\sqrt{1/\mu})$ iterations to get the accuracy $\mathcal E(x_{k}, y_{k}, \mu_{k}) = O(\mu)$. 
% 	\end{theorem}



% 	\begin{proof}[Proof of Theorem 2]
% 	We prove \eqref{eq:Ek-ek} by induction. For $k=0$, it holds by choosing $\mu_0 = D_f(x_0, x^{\star})$. Now suppose that $\mathcal E(x_k, y_k, \mu_{k})\leq (R^2+1)\mu_k$ and let us consider the $k+1$-th iteration. Since $0<\alpha\leq \sqrt{\mu_0}$ and $(1+\alpha)^{-1}\leq 1-\alpha/(1+\sqrt{\mu_0})$, the number of inner iterations $m_{k+1}= (1+\sqrt{\mu_{0}})\ln (2(R^2+1))\mu_{k+1}^{-1/2}$
% 	is chosen so that
% 	\[
% 		\begin{aligned}
% 			(1+\alpha)^{-m_{k+1}} \leq{}& (1-\alpha/(1+\sqrt{\mu_0}))^{m_{k+1}}\\
% 			\leq{}& \exp\left(-\alpha m_{k+1}/(1+\sqrt{\mu_0})\right)= (2(R^2+1))^{-1}.
% 		\end{aligned}
% 	\]
% 	Therefore, by Corollary \eqref{coro:perturbed-decay},
% 	$$
% 	\begin{aligned}
% 		\mathcal E(x_{k+1}, y_{k+1}, \mu_{k+1}) &\leq \frac{1}{2(R^2+1)}\mathcal E(x_k, y_k, \mu_{k+1}) + \mu_{k+1} R^2\\
% 		&\leq \frac{1}{2(R^2+1)}\mathcal E(x_k, y_k, \mu_{k}) + \mu_{k+1} R^2\\
% 		&\leq \frac{1}{2}\mu_{k} + \mu_{k+1} R^2 = (R^2+1) \mu_{k+1}.
% 	\end{aligned}
% 	$$
	
% Since $\mu_k = \mu_0 2^{-k}$, after the $k$-th outer iteration, the overall iteration steps is
% 		\[
% 		\begin{aligned}
% 			M_k=\sum_{i=0}^{k}m_i ={}&  (1+\sqrt{\mu_{0}})\ln (2(R^2+1)) \sum_{i=1}^k \mu_{i}^{-1/2}\\
% 			={}& (1+\sqrt{\mu_{0}})\ln (2(R^2+1))\frac{\sqrt{2}}{\sqrt{2} - 1} \left(\mu_k^{-1/2}-\mu_0^{-1/2}\right).
% 		\end{aligned}
% 		\]
% 		Calculating $\mu_k$ from this and plugging it into \eqref{eq:Ek-ek} proves \eqref{eq:rate}. As a result, the iteration complexity bound $O(\sqrt{1/\mu})$ follows easily.
% 		%Hence, to get $\mathcal E(x_k, y_k, \mu_{k})=O(\mu)$.
% 	%$$\sqrt{L_F}\ln (2(R^2+1)) \sum_{k=1}^n \mu_k^{-1/2} \leq \ln (2(R^2+1))\frac{\sqrt{2}}{\sqrt{2} - 1}\sqrt{\frac{L_F}{\mu}}.
% 	%$$
% 	$\Box$\end{proof}




% \begin{proof}
% The proof uses the one-step sufficient decay estimate established earlier for the Sinkhorn map, together with the discrete Lyapunov estimate for the accelerated scheme.

% For a fixed value of $\mu$, the inner iteration satisfies a perturbed contraction of the form
% \begin{equation}\label{eq:inner-contraction}
% (1+\alpha)\mathcal E(x_{j+1},y_{j+1};\mu,D)\le \mathcal E(x_j,y_j;\mu,D)+\mu R^2,
% \qquad \alpha=\sqrt{2\mu}.
% \end{equation}
% Here the negative term coming from the Sinkhorn step is controlled by the sufficient decay lemma, while Assumption~\ref{assump:boundedness} bounds the perturbation term generated by the weighted norm in the Lyapunov function.

% Iterating \eqref{eq:inner-contraction} for $m_k$ inner steps gives
% \begin{equation}\label{eq:mk-step}
% \mathcal E_{k+1}\le (1+\alpha_k)^{-m_k}\mathcal E_k+\mu_{k+1}R^2,
% \end{equation}
% where $\mathcal E_k:=\mathcal E(x_k,y_k;\mu_k,D)$. The choice of $m_k$ in Algorithm~\ref{alg:A2MD-homotopy} ensures that
% $$
% (1+\alpha_k)^{-m_k}\le \frac{1}{2(R^2+1)}.
% $$
% Therefore,
% $$
% \mathcal E_{k+1}\le \frac{1}{2(R^2+1)}\mathcal E_k+\mu_{k+1}R^2.
% $$
% Assume inductively that
% $$
% \mathcal E_k\le (R^2+1)\mu_k.
% $$
% Since $\mu_{k+1}=\mu_k/2$, we obtain
% $$
% \mathcal E_{k+1}\le \frac{1}{2}\mu_k+\mu_{k+1}R^2=(R^2+1)\mu_{k+1},
% $$
% which proves \eqref{eq:Ek-ek}.

% It remains to express $\mu_k$ in terms of the total number $M_k$ of inner steps. Since
% $$
% \mu_k=\mu_0\,2^{-k},
% \qquad
% \alpha_k=\sqrt{2\mu_k},
% $$
% and $m_k$ is chosen proportional to $\mu_k^{-1/2}$, summing the geometric progression yields
% $$
% M_k\asymp \sum_{i=0}^k \mu_i^{-1/2}
% \asymp \mu_k^{-1/2}-\mu_0^{-1/2}.
% $$
% Equivalently,
% $$
% \mu_k\asymp \frac{1}{\bigl(C_*M_k+\mu_0^{-1/2}\bigr)^2},
% $$
% for the constant $C_*$ defined above. Substituting this relation into \eqref{eq:Ek-ek} gives \eqref{eq:rate}. The final complexity bound follows immediately.
% \end{proof}

\paragraph{Discussion on Complexity}
Let $P^\star$ be the solution of the unregularized optimal transport problem \eqref{eq:OT}. The solution $P^\varepsilon$ of the $\varepsilon$-regularized optimal transport problem \eqref{eq:EOT} satisfies  \cite{peyre2019computational}
$$\langle C,P^\varepsilon\rangle-\langle C,P^\star\rangle\le \varepsilon \log nm,$$ where $P^\star$ is the unregularized solution. The suboptimality of the Acc-Sinkhorn iterate $P_k=P(u(v_k),v_k)$ is composed of two parts: optimization error $\langle C,P_k\rangle-\langle C,P^\varepsilon\rangle$ and regularization bias $\langle C, P^\varepsilon-P^\star\rangle$. More precisely,
$$\langle C,P_k\rangle-\langle C,P^\star\rangle = \langle C,P_k-P^\varepsilon\rangle + \langle C, P^\varepsilon-P^\star\rangle \le \|P_k-P^\varepsilon\|_1\|C\|_\infty + \varepsilon \log nm.$$
To reach accuracy $\tau$, we set $\varepsilon=O(\tau/\log nm)$ and run Acc-Sinkhorn until $\|P_k-P^\varepsilon\|_1 \lesssim \tau/\|C\|_\infty$. Because $\|P_k-P^\varepsilon\|_1 = \|P_k^\top 1-b\|_1 = \|\nabla f(v_k)\|_1$, the convergence rate of Acc-Sinkhorn shows this takes $O(\|C\|_\infty/\tau)$ iterations. The total complexity is $O(n^2/\tau)$, a large improvement over the $O(n^2/\tau^2)$ iterations of Sinkhorn. Since each iteration of Acc-Sinkhorn has a similar cost to Sinkhorn, the total complexity is also better.

% LC{but $||C||_{\infty} = 1/\varepsilon$, if we have the scaling. check the discussion. the complexity should be in terms of $\varepsilon$ to be consistent with the Table.}

\section{Numerical Experiments}
All experiments were run in MATLAB R2025b on an Apple M4 laptop with 24 GB of memory. Random seeds were fixed, and the code is publicly available.

\paragraph{Synthetic Datasets}
We first test Algorithm~\ref{alg:A2MD-homotopy} on synthetic datasets. We generate two discrete distributions $a\in\mathbb{R}^n_+$ and $b\in\mathbb{R}^m_+$ by sampling each entry independently from $[0,1]$ and then normalizing each vector to have unit sum. We generate the cost matrix $C$ in the same way, with entries sampled independently from $[0,1]$ and then rescaled so that $\sum_{i,j} C_{ij}=1$. We compare Acc-Sinkhorn with Sinkhorn for several values of $n$, $m$, and $\varepsilon$. 
\begin{figure}[htp]
\centering
\begin{minipage}[t]{0.35\linewidth}
\vspace{0pt}
Figure~\ref{fig:synthetic_combined} shows the error decay with respect to running time. The results show that Algorithm~\ref{alg:A2MD-homotopy} converges more than $10\times$ faster than Sinkhorn on these synthetic problems, while keeping a similar per-iteration running time. The method is also stable with respect to the choices of $\mu_0$ and $m_0$. These results support the practical efficiency of the proposed method for high-accuracy EOT computation.
\end{minipage}
\hfill
\begin{minipage}[t]{0.625\linewidth}
\vspace{0pt}
\centering
\begin{subfigure}[t]{0.5\linewidth}
    \centering
    \includegraphics[width=0.95\linewidth,height=4.15cm]{figures/n100m50time.pdf}
    \caption{$n=100$, $m=50$, $\varepsilon=10^{-8}$.}
    \label{fig:synthetic1}
\end{subfigure}
\hfill
\begin{subfigure}[t]{0.475\linewidth}
    \centering
    \includegraphics[width=\linewidth,height=4.15cm]{figures/n100m100time.pdf}
    \caption{$n=m = 100, \varepsilon=10^{-8}$.}
    \label{fig:synthetic2}
\end{subfigure}
\end{minipage}
\caption{Synthetic dataset results.}
\label{fig:synthetic_combined}
\end{figure}

% \begin{figure}
% \centering
% \includegraphics[width=\textwidth]{figures/n100m50.pdf}
% \caption{Synthetic dataset results for $n=100, m=50$.}
% \label{fig:synthetic1}
% \end{figure}


% \begin{figure}
% \centering
% \includegraphics[width=\textwidth]{figures/n100m100.pdf}
% \caption{Synthetic dataset results for $n=100, m=100$.}
% \label{fig:synthetic2}
% \end{figure}


% \begin{table}[htbp]
%     \centering
%     \caption{Number of iterations and wall-clock time of the accelerated
%     Sinkhorn algorithm on the synthetic dataset ($\varepsilon = 10^{-8}$)
%     for varying tolerance $\tau$, demonstrating the theoretical
%     $O(\tau^{-1/2})$ dependence on the tolerance (equivalently,
%     $O(k^{-2})$ convergence rate on the relative function gap).}
%     \label{tab:synthetic_convergence}
%     \begin{tabular}{cccccc}
%         \toprule
%         $\tau$ & Iters $k$ & $k / \tau^{-1/2}$ & Time (s) & Time $/ \tau^{-1/2}$ \\
%         \midrule
%         $10^{-4}$ & 1989& 19.89& 0.1141& \\
%         $10^{-5}$ & 2440& 7.7160& 0.1535& \\
%         $10^{-6}$ & 4033& 4.0330& 0.2351& \\
%         $10^{-7}$ & 5389& 1.7042& 0.3137& \\
%         $10^{-8}$ & 6183& 0.6183& 0.3594& \\
%         \bottomrule
%     \end{tabular}
% \end{table}

\paragraph{Color Transfer}

We evaluate the proposed Acc-Sinkhorn algorithm on a color transfer
task between two images, which serves as a canonical application of
optimal transport in image processing \citep{reinhard2001color,
ferradans2014regularized}.


We use two RGB images: a source image of size $1000 \times 669$ pixels
(providing the target color palette) and a content image of size
$1000 \times 750$ pixels (whose structure is to be preserved). From each
image we uniformly subsample $n = 1000$ pixels at random, yielding source
and target point clouds
$\{x_i\}_{i=1}^n$ and
$\{y_j\}_{j=1}^n$,
where each point represents an RGB color value normalized to $[0, 1]^3$.


The cost matrix is defined as the squared Euclidean distance in RGB color
space:
\begin{equation}
    C_{ij} = \| x_i - y_j \|_2^2, \quad i, j \in [n],
\end{equation}
normalized by its maximum entry so that $C \in [0,1]^{n \times n}$.
Both marginals are set to the uniform distribution:
$a_i = b_j = 1/n$.

% \paragraph{Regularized OT problem.}
% We solve the entropy-regularized optimal transport problem
% \begin{equation}
%     \min_{T \geq 0} \; \langle C, T \rangle - \varepsilon H(T)
%     \quad \text{s.t.} \quad T\mathbf{1} = a,\; T^\top\mathbf{1} = b,
% \end{equation}
% where $H(T) = -\sum_{ij} T_{ij}(\log T_{ij} - 1)$ is the entropy of the
% transport plan, for regularization parameter $\varepsilon = 0.05$.

% A moderate value of $\varepsilon$ is appropriate for color transfer, as it
% produces smooth color mappings without excessive blurring; very small
% $\varepsilon$ leads to speckle artifacts while very large $\varepsilon$
% produces overly diffuse transport \citep{peyre2019computational}.

% \paragraph{Solvers and convergence criterion.}
%\LC{Matrix $T$ is not defined.}
We compare the proposed accelerated Sinkhorn algorithm against the standard Sinkhorn algorithm, both implemented in the log-domain for numerical stability \citep{chizat2016scaling}. All solvers are initialized at the zero dual variable $v_0=\mathbf{0}\in\mathbb{R}^{2n}$ and are terminated when the marginal violation satisfies
\begin{equation}\label{eq:l1tol}
\|P\mathbf{1}-a\|_1+\|P^\top\mathbf{1}-b\|_1<\tau.
\end{equation}
We use the $\ell_1$ norm because it gives a direct measure of total feasibility error in the row and column marginals. The threshold $\tau=2/n$ means that the average absolute error is about $1/n$ for each marginal vector, which matches the natural scale of the problem when the marginals are normalized probability vectors. 

%The maximum number of iterations is set to $10^4$.

% \paragraph{Color transfer procedure.}
Given the converged transport plan $P^*$, we use barycentric projection to assign a transferred color to each sampled target pixel. These colors are then propagated to all $750{,}000$ full-resolution target pixels by nearest-neighbor lookup, and the final image is clipped to $[0,1]^3$.

Table~\ref{tab:color_transfer} reports the solver statistics. For small regularization parameters, Acc-Sinkhorn is more than $10\times$ faster than Sinkhorn. Figure~\ref{fig:color_transfer} shows the transferred colors as $\varepsilon$ decreases. Smaller $\varepsilon$ gives a sharper transport plan and better matches the source palette, but it also makes Sinkhorn much slower. Acc-Sinkhorn computes this sharper regime more efficiently.

\begin{table*}[t]
\centering

\begin{minipage}[t]{0.45\textwidth}
\centering
\captionof{table}{Color transfer ($n = 1000$).}
\label{tab:color_transfer}

%\scriptsize
\setlength{\tabcolsep}{3pt}

\resizebox{0.9\linewidth}{!}{
\begin{tabular}{ccccccc}
\toprule
& & \multicolumn{2}{c}{Acc-Sinkhorn} & \multicolumn{2}{c}{Sinkhorn} \\
\cmidrule(lr){3-4} \cmidrule(lr){5-6}
$\varepsilon$ & Tol & It & Time & It & Time \\
\midrule
$1$       & 2e-3& 3& 0.037& 1& 0.013\\
$10^{-1}$ & 2e-3& 5& 0.063& 4& 0.049\\
$10^{-2}$ & 2e-3& 18& 0.217& 37& 0.436\\
$10^{-3}$ & 2e-3& 46& 0.539& 359& 4.232\\
$10^{-4}$ & 2e-3& 239& 3.244& 3507& 49.691\\
\bottomrule
\end{tabular}
}
\end{minipage}
\hfill
\begin{minipage}[t]{0.51\textwidth}
\centering
\captionof{table}{Word alignment (En–Fr, $n=500$).}
\label{tab:word_alignment}

%\scriptsize
\setlength{\tabcolsep}{3pt}

\resizebox{0.9\linewidth}{!}{
\begin{tabular}{ccccccc}
\toprule
& & \multicolumn{2}{c}{Acc-Sinkhorn} & \multicolumn{2}{c}{Sinkhorn} \\
\cmidrule(lr){3-4} \cmidrule(lr){5-6}
$\varepsilon$ & Top 1/5 & It & Time & It & Time \\
\midrule
$1$       & 76.7/94.0& 4& 0.013& 1& 0.003\\
$10^{-1}$ & 78.0/94.8& 6& 0.019& 4& 0.012\\
$10^{-2}$ & 82.8/95.3& 234& 0.638& 5085& 13.745\\
$10^{-3}$ & 81.9/95.7& 268& 0.756& 8254& 22.993\\
$10^{-4}$ & 82.8/94.8& 667& 1.494& 20320& 49.360\\
\bottomrule
\end{tabular}
}
\end{minipage}

\end{table*}

% \begin{table}[h]
% \centering
% \caption{Color transfer ($n = 1000$).
% Number of iterations and wall-clock time for each solver.}
% \label{tab:color_transfer}
% \begin{tabular}{ccccccc}
% \toprule
% & & \multicolumn{2}{c}{Accelerated Sinkhorn} & \multicolumn{2}{c}{Sinkhorn} \\
% \cmidrule(lr){3-4} \cmidrule(lr){5-6}
% $\varepsilon$ & Tolerance & Iters & Time (s) & Iters & Time (s) \\
% \midrule
% $1$       & 2e-3& 3& 0.0373& 1& 0.0128\\
% $10^{-1}$ & 2e-3& 5& 0.0626& 4& 0.0490\\
% $10^{-2}$ & 2e-3& 18& 0.2167& 37& 0.4359\\
% $10^{-3}$ & 2e-3& 46& 0.5391& 359& 4.2321\\
% $10^{-4}$ & 2e-3& 239& 3.2436& 3507& 49.6905\\
% \bottomrule
% \end{tabular}
% \end{table}

\begin{figure}[htbp]
    \centering

    % Row 1: source and target
    \begin{subfigure}[b]{0.32\textwidth}
        \includegraphics[width=0.9\textwidth]{figures/ocean_day.jpg}
        \caption{Source}
    \end{subfigure}
    \hfill
    \begin{subfigure}[b]{0.32\textwidth}
        \includegraphics[width=0.85\textwidth]{figures/ocean_sunset.jpg}
        \caption{Target (content)}
    \end{subfigure}

    \vspace{0.5em}

    % Row 2: results from large reg (bad) to small reg (good)
    \begin{subfigure}[b]{0.18\textwidth}
        \includegraphics[width=\textwidth]{figures/ocean_transfer_1e+00_accsinkhorn.jpg}
        \caption{$\varepsilon = 1$}
    \end{subfigure}
    \hfill
    \begin{subfigure}[b]{0.18\textwidth}
        \includegraphics[width=\textwidth]{figures/ocean_transfer_1e-01_accsinkhorn.jpg}
        \caption{$\varepsilon = 10^{-1}$}
    \end{subfigure}
    \hfill
    \begin{subfigure}[b]{0.18\textwidth}
        \includegraphics[width=\textwidth]{figures/ocean_transfer_1e-02_accsinkhorn.jpg}
        \caption{$\varepsilon = 10^{-2}$}
    \end{subfigure}
    \hfill
    \begin{subfigure}[b]{0.18\textwidth}
        \includegraphics[width=\textwidth]{figures/ocean_transfer_1e-03_accsinkhorn.jpg}
        \caption{$\varepsilon = 10^{-3}$}
    \end{subfigure}
    \hfill
    \begin{subfigure}[b]{0.18\textwidth}
        \includegraphics[width=\textwidth]{figures/ocean_transfer_1e-04_accsinkhorn.jpg}
        \caption{$\varepsilon = 10^{-4}$}
    \end{subfigure}

    \caption{Color transfer results for decreasing regularization parameter
    $\varepsilon$ by Acc-Sinkhorn. The source image (top left) provides the color palette and
    the target image (top right) provides the content. As $\varepsilon$
    decreases, the transport plan becomes sharper and the transferred colors
    more faithfully reproduce the source palette, at the cost of increased
    computation.}
    \label{fig:color_transfer}
\end{figure}


\paragraph{Word Embedding Alignment}
We test Acc-Sinkhorn on bilingual word embedding alignment, where a sharp, near-permutation transport plan is needed and small $\varepsilon$ is important.

We use the aligned multilingual word vectors of \citet{conneau2018word}, trained on Wikipedia with fastText \citep{bojanowski2017enriching}. We take the top $n=500$ English and French words, giving normalized embeddings $\{x_i\}_{i=1}^n,\{y_j\}_{j=1}^n\subset\mathbb{R}^{300}$. Ground-truth pairs are taken from the MUSE bilingual dictionary \citep{conneau2018word}; we keep only pairs that appear in both vocabularies and remove duplicates.

The cost is the cosine distance
\begin{equation}
C_{ij}=1-\langle x_i,y_j\rangle,\qquad i,j\in[n],
\end{equation}
with $C$ normalized to $[0,1]^{n\times n}$. Both marginals are uniform: $a_i=b_j=1/n$.

We compare Acc-Sinkhorn with Sinkhorn, both implemented in the log domain for numerical stability at small $\varepsilon$ \citep{chizat2016scaling}. All solvers start from $v_0=\mathbf{0}\in\mathbb{R}^{2n}$ and stop when the marginal violation \eqref{eq:l1tol} holds with $\tau=0.01\times 2/n$, corresponding to a $1\%$ average marginal error.

Given the converged plan $P^*$, we predict the French translation of English word $i$ by
\begin{equation}
\hat{j}(i)=\arg\max_{j\in[n]} P^*_{ij}.
\end{equation}
We report top-1 and top-5 accuracy over valid evaluation pairs.

Table~\ref{tab:word_alignment} reports accuracy and solver statistics. Acc-Sinkhorn is up to $30\times$ faster than Sinkhorn, with larger gains at smaller $\varepsilon$, where Sinkhorn requires many more iterations.

\section{Conclusion}
This paper presents a simple and efficient accelerated Sinkhorn algorithm for computing entropy-regularized optimal transport that maintains the per-iteration cost of Sinkhorn while achieving a $10\times$--$30\times$ speedup in the small $\varepsilon$ regime. Lyapunov convergence analysis proves the $\mathcal{O}(1/k^2)$ rate under a verifiable stability condition and suggests the overall complexity $\widetilde{\mathcal O}(n^2/\varepsilon)$, compared to $\mathcal{O}(1/k)$ and $\widetilde{\mathcal O}(n^2/\varepsilon^2)$ for Sinkhorn. Numerical experiments on synthetic and real datasets support the practical efficiency of the method. Removing the stability condition and proving the limit step size $\alpha=\sqrt{2\mu}$ directly remain important directions for future work.

\bibliography{references}
\bibliographystyle{plainnat}


\appendix
%%% begin appendices
\section{Proofs for Section~\ref{sec:sinkhorn-dmd}}
\label{app:proof-Sinkhorn}
\subsection{Hessian of the reduced objective \texorpdfstring{$f$}{f}}
% First, we prove a lemma that relates the Hessians of a convex function and its conjugate.
First, we have the following lemma that relates the Hessians of a convex function and its conjugate, which is a standard result in convex analysis; see for example \cite{rockafellar1997convex}.
\begin{lemma}[Inverse Hessian identity]
Let $f$ be $C^2$ and strictly convex, and let $v$ be such that $\xi=\nabla f(v)$ is well defined. Then
$$
\nabla^2 f^*(\xi)=\nabla^2 f(v)^{-1}.
$$
\end{lemma}
% \begin{proof}
% Since $f$ is strictly convex and $C^2$, the gradient map $\nabla f$ is invertible on its domain and $(\nabla f)^{-1}=\nabla f^*$. Fix $v$ and set $\xi=\nabla f(v)$. Then $\nabla f^*(\xi)=v$. Differentiating with respect to $\xi$ yields
% $$
% \nabla^2 f^*(\xi)=\frac{dv}{d\xi}.
% $$
% Differentiating $\xi=\nabla f(v)$ with respect to $v$ gives
% $$
% \frac{d\xi}{dv}=\nabla^2 f(v).
% $$
% By the inverse function theorem,
% $$
% \frac{dv}{d\xi}=\left(\frac{d\xi}{dv}\right)^{-1}=\nabla^2 f(v)^{-1}.
% $$
% Combining the two identities gives $\nabla^2 f^*(\xi)=\nabla^2 f(v)^{-1}$.
% \end{proof}

The following lemma gives the closed-form expression of $\nabla^2 f$.

\begin{lemma}[Hessian of $f$]
The Hessian of $f$ is given by the Schur-complement formula
\begin{equation}
\nabla^2 f(v) = \mathrm{diag}(c_P) - P^\top \mathrm{diag}(r_P)^{-1} P.
\end{equation}
\end{lemma}
\begin{proof}
First, \eqref{eq:grad-f} gives $\nabla f(v)=c_P-b$. Differentiating with respect to $v$ gives
\begin{equation}\label{eq:hessianf}
\nabla^2 f(v)=\frac{d}{dv}\,c_P(u(v),v)=\frac{\partial c_P}{\partial u}\,\frac{du}{dv}+\frac{\partial c_P}{\partial v}.
\end{equation}

Since $P_{ij}=\exp(u_i+v_j-C_{ij})$, we have
$$
\frac{\partial P_{ij}}{\partial u_{i'}}=\delta_{ii'}P_{ij},\qquad
\frac{\partial P_{ij}}{\partial v_{j'}}=\delta_{jj'}P_{ij}.
$$
Therefore, for $r_P\in\mathbb{R}^n$ and $c_P\in\mathbb{R}^m$,
$$
\frac{\partial r_P}{\partial u}=\mathrm{diag}(r_P),\qquad
\frac{\partial r_P}{\partial v}=P,
$$
and
$$
\frac{\partial c_P}{\partial u}=P^\top,\qquad
\frac{\partial c_P}{\partial v}=\mathrm{diag}(c_P).
$$

By the implicit function theorem, differentiating $r_P(u(v),v)=a$ with respect to $v$ gives
$$
\frac{\dd u}{\dd v}=-\left(\frac{\partial r_P}{\partial u}\right)^{-1}\left(\frac{\partial r_P}{\partial v}\right)=-\mathrm{diag}(r_P)^{-1}P.
$$


Substituting into \eqref{eq:hessianf} yields
$$
\nabla^2 f(v)=P^\top\bigl(-\mathrm{diag}(r_P)^{-1}P\bigr)+\mathrm{diag}(c_P)
=\mathrm{diag}(c_P)-P^\top \mathrm{diag}(r_P)^{-1}P,
$$
which is the claimed formula.
\end{proof}
% \begin{proof}
% We derive the Schur-complement formula for $\nabla^2 f(v)$ from the reduced objective
% $$
% f(v)=\min_u F(u,v)=F(u(v),v),
% $$
% where $u(v)$ is defined by the first-order condition $\partial_u F(u(v),v)=0$. Recall
% $$
% P_{ij}(u,v)=\exp(u_i+v_j-C_{ij}),\qquad r_P:=P\mathbf{1}_m,\qquad c_P:=P^\top\mathbf{1}_n,
% $$
% and
% $$
% \partial_u F(u,v)=r_P-a,\qquad \nabla_v F(u,v)=c_P-b.
% $$
% By the chain rule and $\partial_u F(u(v),v)=0$,
% $$
% \nabla f(v)=\nabla_v F(u(v),v)=c_P-b.
% $$
% Hence $\nabla^2 f(v)=\frac{d}{dv}\,c_P(u(v),v)$. To compute this derivative, differentiate the optimality condition $r_P(u(v),v)=a$ with respect to $v$.

% \paragraph{Step 1: Jacobians of $r_P$ and $c_P$.}
% Since $P_{ij}=\exp(u_i+v_j-C_{ij})$, we have
% $$
% \frac{\partial P_{ij}}{\partial u_{i'}}=\delta_{ii'}P_{ij},\qquad
% \frac{\partial P_{ij}}{\partial v_{j'}}=\delta_{jj'}P_{ij}.
% $$
% Therefore, for $r_P\in\mathbb{R}^n$ and $c_P\in\mathbb{R}^m$,
% $$
% \frac{\partial r_P}{\partial u}=\mathrm{diag}(r_P),\qquad
% \frac{\partial r_P}{\partial v}=P,
% $$
% and
% $$
% \frac{\partial c_P}{\partial u}=P^\top,\qquad
% \frac{\partial c_P}{\partial v}=\mathrm{diag}(c_P).
% $$

% \paragraph{Step 2: Implicit differentiation for $u(v)$.}
% Differentiate $r_P(u(v),v)=a$:
% $$
% 0=\frac{d}{dv}\,r_P(u(v),v)=\frac{\partial r_P}{\partial u}\,\frac{du}{dv}+\frac{\partial r_P}{\partial v}.
% $$
% Using the Jacobians above, this gives
% $$
% 0=\mathrm{diag}(r_P)\,\frac{du}{dv}+P,
% \qquad\Rightarrow\qquad
% \frac{du}{dv}=-\mathrm{diag}(r_P)^{-1}P.
% $$

% \paragraph{Step 3: Differentiate $c_P(u(v),v)$.}
% Now
% $$
% \nabla^2 f(v)=\frac{d}{dv}\,c_P(u(v),v)=\frac{\partial c_P}{\partial u}\,\frac{du}{dv}+\frac{\partial c_P}{\partial v}.
% $$
% Substituting $\frac{\partial c_P}{\partial u}=P^\top$, $\frac{\partial c_P}{\partial v}=\mathrm{diag}(c_P)$, and $\frac{du}{dv}=-\mathrm{diag}(r_P)^{-1}P$ yields
% $$
% \nabla^2 f(v)=P^\top\bigl(-\mathrm{diag}(r_P)^{-1}P\bigr)+\mathrm{diag}(c_P)
% =\mathrm{diag}(c_P)-P^\top \mathrm{diag}(r_P)^{-1}P,
% $$
% which is the claimed formula.
% \end{proof}

\subsection{Sublinear Convergence rate of the Sinkhorn algorithm}
% Starting from the dual mirror descent formulation of the Sinkhorn iteration, we prove the convergence rates of the Sinkhorn algorithm. 
In this section, we show that the dual relative smoothness inequality holds with constant $L=1$, which gives a sublinear convergence rate of $\mathcal{O}(1/k)$ for the original Sinkhorn iteration. 
\begin{lemma}[$1$-smoothness]
Let $v_k$ be the sequence generated by the Sinkhorn iteration \eqref{eq:md-general}. Then $f$ is dual relatively smooth with respect to $\phi$ with constant $L=1$.
\end{lemma}
\begin{proof}
Since $\nabla \phi(v)=b.\!*\,\exp(v)-b$, differentiating componentwise gives
$$
\nabla^2 \phi(v)=\mathrm{diag}\bigl(b.\!*\,\exp(v)\bigr).
$$
Using $\nabla^2 \phi^*(\eta)=\nabla^2 \phi\bigl(\nabla \phi^*(\eta)\bigr)^{-1}$ and $\nabla \phi^*(\eta)=\log(\mathbf{1}+\eta./b)$, we obtain
$$
\nabla^2 \phi^*(\eta)
=\mathrm{diag}\!\Bigl(b.\!*\,\exp\bigl(\log(\mathbf{1}+\eta./b)\bigr)\Bigr)^{-1}
=\mathrm{diag}(b+\eta)^{-1},
\qquad \eta>-b.
$$

Let $\xi=\nabla f(v)=c_P-b$. Then $b+\xi=c_P$, and hence
$$
\nabla^2 \phi^*(\xi)=\mathrm{diag}(c_P)^{-1}.
$$

On the other hand,
$$
\nabla^2 f^*(\xi)=\nabla^2 f(v)^{-1},
\qquad
\nabla^2 f(v)=\mathrm{diag}(c_P)-P^\top \mathrm{diag}(r_P)^{-1}P.
$$
Since $P^\top \mathrm{diag}(r_P)^{-1}P\succeq 0$, we have 
%$\nabla^2 f(v)\preceq \mathrm{diag}(c_P)$, and hence $\nabla^2 f(v)^{-1}\succeq \mathrm{diag}(c_P)^{-1}$. Therefore,
$$
\nabla^2 \phi^*(\nabla f(v))=\mathrm{diag}(c_P)^{-1}\preceq \nabla^2 f(v)^{-1}=\nabla^2 f^*(\nabla f(v)),
$$
so the dual relative smoothness inequality holds with $L=1$.
\end{proof}

\begin{theorem}[Theorem 3.9, \cite{maddison2021dual}]
Let $f,\phi^* : \mathbb{R}^d \to \mathbb{R}\cup\{\infty\}$ satisfy dual relative smoothness with constant $L$. If $x_0 \in \operatorname{int}(\operatorname{dom} f)$, then for all $i>0$ the iterates of Algorithm~1.1 satisfy
\begin{equation}
\phi^*(\nabla f(x_i))-\phi^*(0)
\le
\frac{L}{i}\bigl(f(x_0)-f(x_{\min})\bigr).
\end{equation}
\end{theorem}
Therefore, Sinkhorn iteration \eqref{eq:sinkhorn-descent} converges with rate $\mathcal{O}(1/k)$: $\phi^*(\nabla f(v_k))-\phi^*(0)=\mathcal{O}(1/k)$. 
% \begin{corollary}[$\mathcal{O}(1/k)$ convergence]
% Let $v_k$ be the sequence generated by the Sinkhorn iteration \eqref{eq:md-general}. Then the function $f$ converges with rate $\mathcal{O}(1/k)$: $\phi^*(\nabla f(v_k))-\phi^*(0)=\mathcal{O}(1/k)$.    
% \end{corollary}


% \begin{lemma*}[Restatement of Lemma~\ref{lem:sinkhorn-descent}]
% Let
% % \begin{equation}\label{eq:sinkhornstep}
% $v^+:=v-\nabla\phi^*(\nabla f(v))$
% % \end{equation}
% be the Sinkhorn update.
% Then
% \begin{equation}\label{eq:sinkhorn-descent-appendix}
% f(v^+)-f(v)\le -D_{\phi^*}(0,\nabla f(v)).
% \end{equation}
% \end{lemma*}
% \begin{proof}
% Since $f$ is dual relatively smooth with respect to $\phi$ with constant $L=1$, the descent inequality for dual mirror descent gives
% $$
% \phi^*(\nabla f(v^+))\le \phi^*(0)+f(v)-f(v^+)-D_{\phi^*}(0,\nabla f(v)).
% $$
% Because $\phi^*(0)=0$ and $\phi^*\ge 0$, we obtain
% $$
% 0\le \phi^*(\nabla f(v^+))\le f(v)-f(v^+)-D_{\phi^*}(0,\nabla f(v)),
% $$
% which implies \eqref{eq:sinkhorn-descent-appendix}.
% %$$
% %f(v^+)-f(v)\le -D_{\phi^*}(0,\nabla f(v)).
% %$$
% %This is exactly .
% \end{proof}


\subsection{Normalized Sinkhorn and linear convergence}
\label{app:normalized-sinkhorn}
In this section, we show the normalized Sinkhorn iteration \eqref{eq:normalized-sinkhorn} converges linearly.
First, we have the following lemma for the Sinkhorn update.

% Similar to Section~\ref{sec:sinkhorn-dmd}, we rewrite the normalized Sinkhorn iteration in the dual mirror framework. Define
% \begin{equation}
%     \psi^*(\xi):=\phi^*(P_{N^\perp}\xi),
% \end{equation}
% so that $\nabla\psi^*(\xi)=P_{N^\perp}\nabla\phi^*(P_{N^\perp}\xi)$, and $\nabla^2\psi^*(\xi)=P_{N^\perp}\nabla^2\phi^*(P_{N^\perp}\xi)P_{N^\perp}$. Since $\nabla^2\phi^*(P_{N^\perp}\xi)\preceq \nabla^2 f^*(P_{N^\perp}\xi)=\nabla^2 f^*(\xi)$, we also have $\nabla^2\psi^*(\xi)\preceq \nabla^2 f^*(P_{N^\perp}\xi)$, so the dual relative smoothness still holds with constant $L=1$. Consider the \textit{normalized Sinkhorn} iteration
% \begin{equation}\label{eq:normalized-sinkhorn-psi}
% \tilde{v}^0 = P_{N^\perp}v^0, \quad \tilde{v}^{k+1}=\tilde{v}^k-\nabla\psi^*\!\bigl(\nabla f(\tilde{v}^k)\bigr),\quad k\geq 0.
% \end{equation}
% From the construction, \eqref{eq:normalized-sinkhorn-psi} is the same as the normalized Sinkhorn iteration \eqref{eq:normalized-sinkhorn} in the main text, and the iterates $\{\tilde{v}^k\}$ are always in $N^\perp$. Moreover, the following lemma shows that the normalized Sinkhorn is exactly equivalent to the original Sinkhorn iteration \eqref{eq:md-general}.
% \begin{lemma}\label{lem:normalized-sinkhorn}
% The normalized Sinkhorn iterates $\{\tilde{v}^k\}$ are related to the original Sinkhorn iterates $\{v^k\}$ by $\tilde{v}^k=P_{N^\perp}v^k.$ Therefore, the normalized Sinkhorn iteration is equivalent to the original Sinkhorn iteration.
% \end{lemma}
% \begin{proof}
% We prove by induction. The base case $k=0$ is true by definition. Assume $\tilde{v}^k=P_{N^\perp}v^k$. Then
% \begin{align*}
% \tilde{v}^{k+1}&=\tilde{v}^k-\nabla\psi^*\!\bigl(\nabla f(\tilde{v}^k)\bigr)\\
% &=P_{N^\perp}v^k-P_{N^\perp}\nabla\phi^*\!\bigl(P_{N^\perp}\nabla f(P_{N^\perp}v^k)\bigr)\\
% &=P_{N^\perp}\Bigl(v^k-\nabla\phi^*\!\bigl(\nabla f(v^k)\bigr)\Bigr)\\
% &=P_{N^\perp}v^{k+1},
% \end{align*}
% where the first equality uses the definition of normalized Sinkhorn, the second equality uses the induction hypothesis, the third equality uses the fact that $P_N\nabla f(\cdot)=0$, and the last equality uses the definition of Sinkhorn.
% \end{proof}

% We also derive the explicit formula for the mirror function $\psi$. Essentially, the function $\psi$ lives on the subspace $N^\perp$.
% \begin{proposition} \label{prop:psi}
% Define $\psi: N^\perp \to \mathbb{R}$ by
% \[
% \psi(x_\perp) := \min_{\alpha \in \mathbb{R}} \sum_{j=1}^n b_j \Big( e^{x_\perp^j + \alpha} - (x_\perp^j + \alpha) \Big), \quad x_\perp \in N^\perp.
% \]
% Then $\psi$ is well-defined, convex, and satisfies
% \[
% \psi(x_\perp) = B \log \sum_{j=1}^n b_j e^{x_\perp^j}, \quad B := \sum_{j=1}^n b_j,
% \]
% up to an additive constant. That is, $\psi$ depends only on the sum-zero component $x_\perp$ of $x$ and is invariant under shifts along $N$.
% \end{proposition}

% \begin{proof}
% By definition of the convex conjugate and the restriction, we have
% \[
% \psi(x) = \sup_{y \in \mathbb{R}^n} \big\{ \langle x, y \rangle - \psi^*(y) \big\} 
% = \sup_{y \in \mathbb{R}^n} \big\{ \langle x, y \rangle - \phi^*(P_{N^\perp} y) \big\}.
% \]

% Using the standard property of convex conjugates under a linear map $A$:
% \[
% (f^* \circ A)^*(x) = \min_{v \in \mathbb{R}^n: A^\top v = x} f(v),
% \]
% and noting that $P_{N^\perp}^\top = P_{N^\perp}$, we obtain
% \[
% \psi(x) = \min_{v \in \mathbb{R}^n : P_{N^\perp} v = x} \phi(v).
% \]

% Observe that the constraint $P_{N^\perp} v = x_\perp$ implies
% \[
% v = x_\perp + \alpha \mathbf{1}, \quad \alpha \in \mathbb{R}.
% \]

% Therefore,
% \[
% \psi(x_\perp) = \min_{\alpha \in \mathbb{R}} \sum_{j=1}^n b_j \Big( e^{x_\perp^j + \alpha} - (x_\perp^j + \alpha) \Big).
% \]

% This is a one-dimensional convex minimization problem in $\alpha$. Define
% \[
% f(\alpha) := \sum_{j=1}^n b_j e^{x_\perp^j} e^{\alpha} - \sum_{j=1}^n b_j (x_\perp^j + \alpha) = e^\alpha \sum_{j=1}^n b_j e^{x_\perp^j} - \sum_{j=1}^n b_j x_\perp^j - B \alpha.
% \]

% Taking the derivative with respect to $\alpha$:
% \[
% f'(\alpha) = e^\alpha \sum_{j=1}^n b_j e^{x_\perp^j} - B.
% \]

% Setting $f'(\alpha^*) = 0$ gives the unique minimizer
% \[
% \alpha^* = \log \frac{B}{\sum_j b_j e^{x_\perp^j}}.
% \]

% Plugging $\alpha^*$ back into $f$, we obtain
% \[
% \psi(x_\perp) = \sum_j b_j e^{x_\perp^j + \alpha^*} - \sum_j b_j x_\perp^j - B \alpha^* = B - 0 - B \log \frac{B}{\sum_j b_j e^{x_\perp^j}} = B \log \sum_j b_j e^{x_\perp^j}-\sum_j b_j x_\perp^j.
% \]

% This shows that $\psi$ is convex, as the logarithm of a positive weighted sum of exponentials is convex. Also, $\psi$ is well-defined as a function on $N^\perp$ and is invariant under shifts along $N$.
% \end{proof}

% \begin{remark}
% Let $x_\perp \in N^\perp$. Then
% \[
% \nabla \psi(x_\perp) = B\frac{b .* e^{x_\perp}}{\sum_j b_j e^{x_\perp^j}} - b\in N^\perp, \quad \text{so that } \sum_j (\nabla \psi(x_\perp))_j = 0.
% \]

% This is useful for the proof of Acc-Sinkhorn, which we will present in section~4.
% \end{remark}

% Therefore, the normalized Sinkhorn iteration is exactly a dual mirror descent iteration for the objective $f|_{P_{N^\perp}}$, which satisfies the mirror P-L inequality and has the same minimizer as $f$ (up to a shift in $N$).

\begin{lemma}\label{lem:sinkhorn-descent}
Let
$
v^+:=v-\nabla\phi^*(\nabla f(v))
$
be the Sinkhorn update. Then
\begin{equation}\label{eq:sinkhorn-descent}
f(v^+)-f(v)
\le
-D_{\phi^*}(0,\nabla f(v))
-D_{\phi^*}(\nabla f(v^+),0).
\end{equation}
\end{lemma}
\begin{proof}
Let $g:=\nabla f(v)$ and $g^+:=\nabla f(v^+)$. Since $v-v^+=\nabla\phi^*(g)-\nabla\phi^*(0)$, expansion at $v^+$ and Bregman duality give
$$
f(v^+)-f(v)
=
\langle g^+,v^+-v\rangle-D_f(v,v^+)
=
-\langle g^+,\nabla\phi^*(g)-\nabla\phi^*(0)\rangle-D_{f^*}(g^+,g).
$$
By the three-point identity,
$$
-\langle g^+-0,\nabla\phi^*(g)-\nabla\phi^*(0)\rangle
=
-D_{\phi^*}(g^+,0) - D_{\phi^*}(0,g) + D_{\phi^*}(g^+,g).
$$
Using dual relative smoothness, $D_{\phi^*}(g^+,g)\le D_{f^*}(g^+,g)$, we obtain \eqref{eq:sinkhorn-descent}.
%$$
%f(v^+)-f(v)\le -D_{\phi^*}(\nabla f(v^+),0)-D_{\phi^*}(0,\nabla f(v)).
%$$
\end{proof}

Then, we show that $f$ is coercive on the subspace $N^\perp$. 
\begin{lemma}\label{lem:coercivity}
    $f(v)\to+\infty$ as $\|v\|\to\infty$ along $N^\perp$.
\end{lemma}
\begin{proof}
The minimizer $u(v)$ satisfies $\partial_u F = 0$, i.e.\ $\sum_j e^{u_i(v)+v_j-C_{ij}} = a_i$
for all $i$, giving $u_i(v) = \log a_i - \log(Ke^v)_i$ where $(Ke^v)_i := \sum_j e^{v_j - C_{ij}}$.
Summing over $i$ shows $\sum_{i,j} e^{u_i(v)+v_j - C_{ij}} = \mathbf{1}^\top a = 1$,
so substituting $u(v)$ into $F$ gives
\[
  f(v)
  = 1 - \langle u(v), a\rangle - \langle v,b\rangle
  = 1 - \sum_i a_i \log a_i + \sum_i a_i\log(Ke^v)_i - \langle v,b\rangle.
\]
Since $(Ke^v)_i \ge e^{v_j - C_{ij}}$ for every $j$, taking $j^*=\arg\max_j v_j$ gives
$\log(Ke^v)_i \ge \max_j v_j - C_{ij^*} \ge \max_j v_j - \|C\|_\infty$.
Summing over $i$ with weights $a_i$:
\begin{equation}\label{eq:lb}
  f(v) \;\ge\; 1 - \sum_i a_i \log a_i - \|C\|_\infty + \max_j v_j - \langle v,b\rangle.
\end{equation}
Since $b>0$ and $\mathbf{1}^\top b=1$, a weighted average cannot exceed its maximum:
$\langle v,b\rangle = \sum_j b_j v_j \le \max_j v_j$,
with equality if and only if all $v_j$ are equal, i.e.\ $v\in N:=\mathrm{span}\{\mathbf{1}\}$.
Hence $g(v) \ge 0$, with $g(v)=0$ only on $N$. On $N^\perp\setminus\{0\}$ we thus
have $g(v)>0$. Since $g$ is continuous and positively $1$-homogeneous, compactness
of $\{v\in N^\perp:\|v\|=1\}$ yields $c>0$ with $g(v)\ge c\|v\|$ for all $v\in N^\perp$.
\end{proof}
Therefore, every sublevel set of $f|_{N^\perp}$ is bounded, thus compact.
The next lemma shows that the function $f$ satisfies the Polyak--\L ojasiewicz (P-L) inequality on sublevel sets of $f|_{P_{N^\perp}}$.
\begin{lemma}\label{prop:PL}
For every sublevel set $\mathcal S\subset N^\perp$, there exists $0<\mu < 1$ such that
\begin{equation}\label{eq:mirror-PL-appendix}
f(v)-f(v^*)
\le
\frac1\mu D_{\phi^*}(0,\nabla f(v)),
\qquad
v\in\mathcal S.
\end{equation}
\end{lemma}

\begin{proof}
It follows from \eqref{eq:HessianL} that
\[
\ker \nabla^2 f(v)=\mathrm{span}\{\mathbf 1\}.
\]
Indeed, for any $h\in\mathbb R^n$,
\[
h^\top \nabla^2 f(v)h
=
\sum_j c_{P,j}h_j^2
-
\sum_i \frac1{r_{P,i}}
\Bigl(\sum_j P_{ij}h_j\Bigr)^2.
\]
By Cauchy--Schwarz,
\[
\Bigl(\sum_j P_{ij}h_j\Bigr)^2
\le
r_{P,i}\sum_j P_{ij}h_j^2,
\]
with equality iff $h_j$ is constant on the support of the $i$-th row of $P$.
Since $P_{ij}>0$, equality holds iff $h$ is constant.
Therefore $\nabla^2 f(v)$ is positive definite on $N^\perp$.

Since $f$ is $C^2$, its Hessian satisfies
\[
\nabla^2 f(v)\succ 0
\qquad\text{on }N^\perp.
\]

By continuity of $\nabla^2 f$ and compactness of the sublevel set $\mathcal S$, there exists $\lambda_{\min}>0$ such that
\[
h^\top \nabla^2 f(v) h
\ge
\lambda_{\min}\|h\|^2,
\qquad
v\in\mathcal S,
\quad
h\in N^\perp.
\]
Hence $f$ is strongly convex on $\mathcal S\subset N^\perp$, and the standard Polyak--\L ojasiewicz inequality gives
\begin{equation}
\label{eq:euclidean-PL}
f(v)-f(v^*)
\le
\frac{1}{2\lambda_{\min}}
\|\nabla f(v)\|^2.
\end{equation}

Next, recall
\[
\nabla^2\phi^*(\xi)
=
\mathrm{diag}(b+\xi)^{-1}.
\]
Since
\[
b+\xi=c_P,
\]
and $v\in\mathcal S$, the coordinates of $c_P$ are uniformly bounded above on $\mathcal S$.
Therefore there exists $M_\phi>0$ such that
\[
\nabla^2\phi^*(\xi)
\succeq
\frac1{M_\phi}I.
\]
By Taylor expansion of the Bregman divergence,
\[
D_{\phi^*}(0,\xi)
=
\frac12
\xi^\top
\nabla^2\phi^*(\tilde\xi)\xi
\ge
\frac1{2M_\phi}\|\xi\|^2,
\]
for some $\tilde\xi$ on the segment joining $0$ and $\xi$.
Substituting $\xi=\nabla f(v)$ yields
\[
\|\nabla f(v)\|^2
\le
2M_\phi D_{\phi^*}(0,\nabla f(v)).
\]
Combining with \eqref{eq:euclidean-PL} gives
\[
f(v)-f(v^*)
\le
\frac{M_\phi}{\lambda_{\min}}
D_{\phi^*}(0,\nabla f(v)).
\]
Absorbing constants proves \eqref{eq:mirror-PL-appendix}.
\end{proof}

Combining the above results gives a linear convergence rate for the normalized Sinkhorn iteration \eqref{eq:normalized-sinkhorn}, and hence for the original Sinkhorn iteration \eqref{eq:md-general}.
\begin{theorem}\label{thm:sinkhorn-linear-conv}
Let $v_k$ be the sequence generated by the Sinkhorn iteration \eqref{eq:md-general}. Then the function $f$ converges with rate 
\[
f(v_k)-f(v^*)
\le
(1-\mu)^k(f(v_0)-f(v^*)).
\]
\end{theorem}
\begin{proof}
By Lemma~\ref{lem:sinkhorn-descent} and Lemma~\ref{lem:coercivity}, the sequence of iterates $\{ v_k \}$ remain in the bounded sublevel set $\mathcal S=\{v:\ f(v)\le f(v_0)\}$. By Proposition~\ref{prop:PL}, there exists $\mu>0$ such that
\[f(v)-f(v^*)
\le
\frac1\mu
D_{\phi^*}(0,\nabla f(v)),
\qquad v\in\mathcal S.
\]
By Lemma~\ref{lem:sinkhorn-descent},
\[
f(v_{k+1})-f(v_k)\le -D_{\phi^*}(0,\nabla f(v_k)).
\]
Combining the two inequalities gives
\[
f(v_{k+1})-f(v^*)\le (1-\mu)(f(v_k)-f(v^*)).
\]
Iterating this inequality yields the claimed linear convergence rate.
\end{proof}

\section{Discussion on Lemma~\ref{assump:boundedness}}
\subsection{Sufficient conditions for the \texorpdfstring{$x$}{x}-part}
The $x$-part of Lemma~\ref{assump:boundedness} follows from more standard uniform bounds on the iterates and gradients. 
\begin{lemma}
suppose there exists a constant $C>0$ such that
$$
\|x_k-x^\star\|_2\le C,
\qquad
\|\nabla f(x_k)\|_\infty\le C
\qquad\text{for all }k,
$$
then the $x$-part of Lemma~\ref{assump:boundedness} holds.
\end{lemma}
\begin{proof}
We have
$$
\|x_k-x^\star\|_{D(x_k^\phi)}^2
\le \lambda_{\max}(D(x_k^\phi))\|x_k-x^\star\|_2^2
\le C^2\,\lambda_{\max}(D(x_k^\phi)).
$$
Since
$$
\nabla\phi(z)=b.*e^z-b,
$$
the matrix $D(z)$ is given by
$$
D(z)=\mathrm{diag}\!\left(b.*\frac{e^z-1}{z}\right),
$$
with the continuous extension at $z_i=0$. Substituting $z=x_k^\phi=\nabla\phi^*(\nabla f(x_k))$ yields
$$
x_k^\phi=\log\!\bigl(\mathbf{1}+\nabla f(x_k)./b\bigr),
$$
and hence
$$
D(x_k^\phi)=\mathrm{diag}\!\left(b.*\frac{e^{x_k^\phi}-1}{x_k^\phi}\right)
=\mathrm{diag}\!\left(\frac{\nabla f(x_k)}{\log(1+\nabla f(x_k)./b)}\right).
$$
If $\|\nabla f(x_k)\|_\infty\le C$, then each coordinate of $x_k^\phi$ lies in a bounded interval, provided $b_i$ is bounded away from $0$. Since the scalar function
$$
h(t):=\frac{e^t-1}{t}
$$
extends continuously to $t=0$ with $h(0)=1$ and is bounded on every bounded interval, there exists a constant $C_D>0$ such that
$$
\lambda_{\max}(D(x_k^\phi))\le C_D
\qquad\text{for all }k.
$$
Therefore,
$$
\|x_k-x^\star\|_{D(x_k^\phi)}^2\le C^2 C_D,
$$
so Lemma~\ref{assump:boundedness} holds with $R=C\sqrt{C_D}$.
\end{proof}
 In this sense, the assumption is mild once the iterates stay in a bounded region and the gradient remains uniformly bounded.

\subsection{Sufficient conditions of Lemma~\ref{assump:boundedness} }
\label{app:boundedness-iterates}
In this section, we verify that the iterates of Acc-Sinkhorn remain bounded under certain mild, verifiable conditions.



% \subsection{Continuous level}

% Let $v = y - x$. The VOS flow can be written as
% \[
% \begin{cases}
% x' = v,\\
% v' = -v - \mu^{-1}\nabla f(x).
% \end{cases}
% \]
% Define the Lyapunov function
% \[
% E(x,v) = f(x) - f^* + \frac{\mu}{2}\|v\|^2.
% \]
% Then
% \[
% \frac{d}{dt}E(x,v) = \langle \nabla f(x), v\rangle + \mu\langle v, -v - \mu^{-1}\nabla f(x)\rangle = -\mu\|v\|^2 \le 0.
% \]
% Hence $E$ is nonincreasing.

% If, in addition, the sublevel sets of $f$ are bounded, for example if $f$ is coercive or strongly
% convex, then $x(t)$ remains bounded. Since
% \[
% v = y - x,
% \]
% we have
% \[
% y - x^* = (x - x^*) + v,
% \]
% and therefore $y(t)$ is also bounded whenever both $x(t)$ and $v(t)$ are bounded.

% Notice that $\mu > 0$ is an arbitrary positive parameter, not the strong convexity constant of
% $f$. In fact, convexity is not needed in the continuous argument; coercivity of $f$ is enough.

% \subsection{Discrete level}

Consider the Acc-Sinkhorn scheme
\begin{align*}
x_{k+1} - x_k &= \alpha(y_k - x_{k+1}) - \frac{1}{\alpha}\nabla \phi^*(\nabla f(x_k)),\\
y_{k+1} - y_k &= \alpha(x_{k+1} - y_{k+1}) - \frac{\alpha}{\mu}\nabla \phi^*(\nabla f(x_{k+1}))
\end{align*}
with $\alpha = \sqrt{c\mu}$ for some $c > 0$, and $\mu>0$. We have the following result on the boundedness of the iterates.
\begin{theorem}\label{thm:condition-boundedness}
Assume $\alpha^2 = \mu/L$, then Lemma~\ref{assump:boundedness} holds if the following conditions hold for all $k$:
% \begin{enumerate}
%     % \item $$\|\nabla f(x_k)\|_{D_k^{-1} D_{k+1} D_k^{-1}}^2 < \frac{4-8\alpha-\alpha^2}{2\alpha(1+\alpha)^2}D_{\phi^*}(0,\nabla f(x_{k}))$$
%     % \item $$\|v_k\|_{D_{k+1}-D_k} < 2\alpha\|v_{k+1}\|_{D_{k+1}}.$$
%         \item $$\|\nabla f(x_k)\|_{D_k^{-1} D_{k+1} D_k^{-1}}^2 < \frac{4-\alpha^2-2\alpha^{1/2}}{2\alpha(1+\alpha)^2}D_{\phi^*}(0,\nabla f(x_{k}))$$
%     \item $$\|v_k\|_{D_{k+1}-D_k} < \alpha^{3/4}\|v_{k+1}\|_{D_{k+1}}.$$
% \end{enumerate}
\begin{equation}
    \begin{aligned}
\|\nabla f(x_k)\|_{D_k^{-1} D_{k+1} D_k^{-1}}^2 <{}& \frac{4-\alpha^2-2\alpha^{1/2}}{2\alpha(1+\alpha)^2}D_{\phi^*}(0,\nabla f(x_{k}))\\
\|v_k\|_{D_{k+1}-D_k} <{}& \alpha^{3/4}\|v_{k+1}\|_{D_{k+1}}.
    \end{aligned}
\end{equation}
\end{theorem}

Several remarks are in order.
\begin{remark}
Both conditions are easy to verify, as they only require the knowledge of the iterates and the gradients.
\end{remark}
\begin{remark}
Condition 1 is satisfied when $\alpha$ is sufficiently small, since the left-hand side approaches $\|\nabla f(x_k)\|_{D_k^{-1}}^2$ while the right-hand side scales as $\mathcal{O}(1/\alpha)$ as $\alpha\to 0$.
Condition 2 is also satisfied when $\alpha$ is sufficiently small. This is because the left-hand side scales as $O(\alpha)$ while the right-hand side scales as $O(\alpha^{3/4})$ as $\alpha\to 0$.
Therefore, Lemma~\ref{assump:boundedness} holds, under a backtracking line search procedure to choose $\alpha$. This implies the boundedness of the iterates.
\end{remark}

\begin{remark}
    In general, for any $c\in [1,2)$, the conditions can be similaring derived, with a different choice of $\xi$. In practice, we simply choose the limiting case $c=2$.
\end{remark}

\subsection{Proof sketch of Theorem~\ref{thm:condition-boundedness}}
Below is a concise proof sketch for an intuitive understanding. The full proof is deferred after that.
\begin{proof}[Proof sketch]
Set $v_k := y_k - x_k$, $\beta:=1/\alpha$ and $\tilde\alpha := \alpha/(1+\alpha)$.
Eliminating $y_k$ from the scheme,
\begin{align}
  x_{k+1} - x_k &= \tilde\alpha\bigl(v_k - \beta\nabla\phi^*(\nabla f(x_k))\bigr),
  \label{eq:xdiff}\\
  v_{k+1} - v_k &= -\alpha(\alpha+2)\,v_{k+1}
    + \alpha\beta\,\nabla\phi^*(\nabla f(x_k))
    - \tfrac{\alpha(1+\alpha)}{\mu}\,\nabla\phi^*(\nabla f(x_{k+1})).
  \label{eq:vdiff}
\end{align}
Define the shifted Lyapunov function
\[
  \widetilde{E}_k
  := \underbrace{\bigl(f(x_k)-f(x^\star)\bigr)}_{\text{potential}}
   + \underbrace{\tfrac{\mu}{2}\|v_k\|_{D_k}^2}_{\text{kinetic}}
   - \tfrac{\theta}{\tilde{L}}\,D_{\phi^*}(0,\nabla f(x_k)),
\]
where $\tilde{L}=L(1+\alpha)$, $D_k := D(p(x_k))$, and $\theta\in(0,1]$ is a free parameter.
The subtracted Bregman term keeps $\widetilde{E}_k$ equivalent to $E_k$ up to constants
(using $f(x)-f(x^\star) \ge \frac{1}{L}D_{\phi^*}(0,\nabla f(x))$) while allowing a larger
step size $\alpha$.

\textbf{One-step decrease.}
Using dual $L$-smoothness, the three-point Bregman identity, and the kinetic energy
update~\eqref{eq:vdiff}, one computes
\[
  \widetilde{E}_{k+1} - \widetilde{E}_k
  \le
  -\Bigl[\mu\alpha(\alpha+2) - \tfrac{\delta^2}{2\eta}\Bigr]\|v_{k+1}\|_{D_{k+1}}^2
  - \tfrac{1-\theta}{\tilde{L}}D_{\phi^*}(0,\nabla f(x_k))
  + R_k,
\]
where $\delta := \alpha^2(\alpha+2)/(1+\alpha)$, $\eta>0$ is a free parameter from Young's
inequality, and $R_k$ collects the two remainder terms
\[
  R_k :=
  \tfrac{\mu}{L}\dual{D_{k+1}v_{k+1},\nabla\phi^*(\nabla f(x_k))}
  + \tfrac{\mu}{2}\|v_k\|_{D_{k+1}-D_k}^2.
\]
Setting $\alpha^2 = \mu/L$ and choosing $\eta = 2\alpha/[L(1+\alpha)^2]$,
$\theta = (1/2+\alpha)/(1+\alpha)$ annihilates the $\|\nabla f(x_{k+1})\|^2$
coefficient and makes the kinetic and potential damping terms explicitly negative.

\textbf{Reduction to two verifiable conditions.}
A final application of Young's inequality to $R_k$ shows that
$\widetilde{E}_{k+1} - \widetilde{E}_k \le 0$ follows from
\[
\begin{aligned}
  &\|\nabla f(x_k)\|_{D_k^{-1}D_{k+1}D_k^{-1}}^2
    \;\le\;
    \xi\,\tfrac{4(1-c/2)}{1+\alpha}\,D_{\phi^*}(0,\nabla f(x_k)),\\[6pt]
  &\|v_k\|_{D_{k+1}-D_k}^2
    \;\le\;
    \kappa(\alpha,c,\xi)\,\|v_{k+1}\|_{D_{k+1}}^2.
\end{aligned}
\]
where $c=1$ and $\kappa(\alpha,c,\xi)>0$ for sufficiently small $\alpha$.
Both conditions measure how much the metric $D_k$ changes between steps;
they hold whenever the iterates and gradients remain uniformly bounded,
which is the case when $\alpha$ is small enough.
Since $\widetilde{E}_k$ is equivalent to $E_k$ and $\widetilde{E}_{k+1}\le\widetilde{E}_k$,
the sequence $\{E_k\}$ is non-increasing, yielding the claimed boundedness
$\|x_k - x^\star\|_{D(p(x_k))} \le R$ and $\|y_k - x^\star\| \le R$.
\end{proof}

\subsection{Full proof of Theorem~\ref{thm:condition-boundedness}}
\begin{proof}
To prove this theorem, we reformulate the Acc-Sinkhorn scheme using a different set of variables $(x,v)$, and then analyze the difference of the a new Lyapunov function $\tilde{E}_k$. The two conditions are derived by ensuring $\tilde{E}_{k+1}-\tilde{E}_k$ is nonpositive, which implies the boundedness of the iterates. 
\paragraph{Reformulation.} Introduce
\[
v_k := y_k - x_k.
\]
From the first equation,
\[
x_{k+1} - x_k = \alpha\bigl[(y_k - x_k) - (x_{k+1} - x_k)\bigr] - \alpha\beta\nabla \phi^*(\nabla f(x_k)),
\]
hence
\[
(1+\alpha)(x_{k+1}-x_k) = \alpha v_k - \alpha\beta\nabla \phi^*(\nabla f(x_k)).
\]
Therefore
\[
x_{k+1} - x_k = \tilde{\alpha}\bigl(v_k - \beta\nabla \phi^*(\nabla f(x_k))\bigr),
\qquad \tilde{\alpha} := \frac{\alpha}{1+\alpha}.
\]
The second equation becomes
\[
y_{k+1} - y_k = -\alpha v_{k+1} - \frac{\alpha}{\mu}\nabla \phi^*(\nabla f(x_{k+1})).
\]
Hence
\[
v_{k+1} - v_k = (y_{k+1}-y_k) - (x_{k+1}-x_k)
= -\alpha v_{k+1} - \frac{\alpha}{\mu}\nabla \phi^*(\nabla f(x_{k+1}))
- \tilde{\alpha}\bigl(v_k - \beta\nabla \phi^*(\nabla f(x_k))\bigr).
\]
Using
\[
v_k = v_{k+1} - (v_{k+1} - v_k),
\]
we obtain
\[
(1-\tilde{\alpha})(v_{k+1}-v_k)
= -(\alpha+\tilde{\alpha})v_{k+1}
+ \tilde{\alpha}\beta\nabla \phi^*(\nabla f(x_k))
- \frac{\alpha}{\mu}\nabla \phi^*(\nabla f(x_{k+1})).
\]
Since
\[
\tilde{\alpha} = \frac{\alpha}{1+\alpha}, \qquad 1 - \tilde{\alpha} = \frac{1}{1+\alpha},
\]
it follows that
\begin{equation}% \label{eq:vdiff}
v_{k+1} - v_k
= -\alpha(\alpha+2)\,v_{k+1}
+ \alpha\beta\,\nabla \phi^*(\nabla f(x_k))
- \frac{\alpha(1+\alpha)}{\mu}\,\nabla \phi^*(\nabla f(x_{k+1})).
\end{equation}

\paragraph{Difference of potential energy.}
Using the $L$-dual smoothness and convexity of $f$ give
\[
f(x_{k+1}) - f(x_k)
\le \tilde{\alpha}\langle\nabla f(x_{k+1}), v_k\rangle
- \tilde{\alpha}\beta\langle\nabla f(x_{k+1}), \nabla \phi^*(\nabla f(x_k))\rangle
- \frac{1}{L}D_{\phi^*}(\nabla f(x_{k+1}), \nabla f(x_k)).
\]
Now assume
\[
\alpha\beta = \frac{1}{L}, \qquad
\tilde{\alpha}\beta = \frac{1}{\tilde{L}} \le \frac{1}{L}, \qquad
\tilde{L} = L(1+\alpha).
\]
Using the three -point identity
\[
D_{\phi^*}(a,b) + D_{\phi^*}(b,c) - D_{\phi^*}(a,c) = \langle a-b, \nabla \phi^*(c) - \nabla \phi^*(b)\rangle,
\]
we obtain
\[
-\tilde{\alpha}\beta\langle\nabla f(x_{k+1}),\nabla \phi^*(\nabla f(x_k))\rangle
- \frac{1}{2L}D_{\phi^*}(\nabla f(x_{k+1}), \nabla f(x_k))
\le -\frac{1}{\tilde{L}}D_{\phi^*}(\nabla f(x_{k+1}), 0) - \frac{1}{\tilde{L}}D_{\phi^*}(0, \nabla f(x_k)).
\]
Therefore,
\begin{equation}\label{eq:fdiff}
f(x_{k+1}) - f(x_k)
\le \tilde{\alpha}\langle\nabla f(x_{k+1}), v_k\rangle
- \frac{1}{\tilde{L}}D_{\phi^*}(\nabla f(x_{k+1}), 0) - \frac{1}{\tilde{L}}D_{\phi^*}(0, \nabla f(x_k)).
\end{equation}

\paragraph{Difference of kinetic energy.} Also,
\[
\frac{\mu}{2}\|v_{k+1}\|_{D_{k+1}}^2 - \frac{\mu}{2}\|v_k\|_{D_{k+1}}^2
= \mu\langle D_{k+1}v_{k+1}, v_{k+1}-v_k\rangle - \frac{\mu}{2}\|v_{k+1}-v_k\|_{D_{k+1}}^2.
\]
Substituting \eqref{eq:vdiff},
we obtain
\begin{align}
\frac{\mu}{2}\|v_{k+1}\|_{D_{k+1}}^2 - \frac{\mu}{2}\|v_k\|_{D_{k+1}}^2
&= -\mu\alpha(\alpha+2)\|v_{k+1}\|_{D_{k+1}}^2
+ \mu\alpha\beta\langle D_{k+1}v_{k+1},\nabla \phi^*(\nabla f(x_k))\rangle \notag\\
&\quad - \alpha(1+\alpha)\langle v_{k+1},\nabla f(x_{k+1})\rangle
- \frac{\mu}{2}\|v_{k+1}-v_k\|_{D_{k+1}}^2. \label{eq:vdiff2}
\end{align}

\paragraph{Cancellation.} Let
\[
\Delta v_k := v_{k+1} - v_k, \qquad v_k = v_{k+1} - \Delta v_k.
\]
Then
\begin{align*}
&\tilde{\alpha}\langle\nabla f(x_{k+1}), v_k\rangle
- \alpha(1+\alpha)\langle v_{k+1}, \nabla f(x_{k+1})\rangle\\
&\quad = -\tilde{\alpha}\langle\nabla f(x_{k+1}), \Delta v_k\rangle
- \delta\langle\nabla f(x_{k+1}), v_{k+1}\rangle,
\end{align*}
where
\[
\delta := \alpha(1+\alpha) - \tilde{\alpha} = \frac{\alpha^2(\alpha+2)}{1+\alpha}.
\]

\paragraph{Larger step size.} Introduce the shifted energy
\[
\widetilde{E}_k := E_k - \frac{\theta}{\tilde{L}}D_{\phi^*}(0,\nabla f(x_k)).
\]
Since
\[
f(x) - f^* \ge \frac{1}{L}D_{\phi^*}(0,\nabla f(x)),
\]
we have for all $\theta \in (0,1+\alpha)$,
\[
\widetilde{E}_k \ge c(\theta)\|\nabla f(x)\|^2.
\]
Moreover,
\begin{align*}
\widetilde{E}_{k+1} - \widetilde{E}_k
&\le -\mu\alpha(\alpha+2)\|v_{k+1}\|_{D_{k+1}}^2
+ \frac{\mu}{L}\langle D_{k+1}v_{k+1}, \nabla\phi^*(\nabla f(x_k))\rangle\\
&\quad - \delta\langle\nabla f(x_{k+1}), v_{k+1}\rangle
- \tilde{\alpha}\langle\nabla f(x_{k+1}), \Delta v_k\rangle\\
&\quad - \frac{\mu}{2}\|\Delta v_k\|_{D_{k+1}}^2
- \frac{1-\theta}{\tilde{L}}D_{\phi^*}(0,\nabla f(x_k))
- \frac{1}{\tilde{L}}D_{\phi^*}(\nabla f(x_{k+1}),0)
- \frac{\theta}{\tilde{L}}D_{\phi^*}(0,\nabla f(x_{k+1})).
\end{align*}
Using Young's inequality again,
% \[
% \frac{\mu}{L}\langle D_{k+1}v_{k+1}, \nabla f(x_k)\rangle
% \le \frac{1}{4\tilde{L}}\|\nabla f(x_k)\|^2
% + \frac{\mu^2(1+\alpha)}{2L}\|v_{k+1}\|_{D_{k+1}}^2,
% \]
\[
-\tilde{\alpha}\dual{\nabla f(x_{k+1}),\Delta v_k}\leq \frac{\tilde{\alpha}^2}{2\mu}\|\nabla f(x_{k+1})\|_{D_{k+1}^{-1}}^2 + \frac{\mu}{2}\|\Delta v_k\|_{D_{k+1}}^2,
\]
and
\[
-\delta\langle\nabla f(x_{k+1}), v_{k+1}\rangle
\le \frac{\eta}{2}\|\nabla f(x_{k+1})\|_{D_{k+1}^{-1}}^2 + \frac{\delta^2}{2\eta}\|v_{k+1}\|_{D_{k+1}}^2.
\]
Therefore
\begin{equation}\label{eq:Etilde}
    \begin{aligned}
\widetilde{E}_{k+1} - \widetilde{E}_k
\le{}& -\left[\mu\alpha(\alpha+2) - \frac{\delta^2}{2\eta}\right]\|v_{k+1}\|_{D_{k+1}}^2\\
&{}-\frac{1}{\tilde{L}}D_{\phi^*}(\nabla f(x_{k+1}),0)-\frac{\theta}{\tilde{L}}D_{\phi^*}(0,\nabla f(x_{k+1}))+\left(\frac{\tilde{\alpha}^2}{2\mu}+\frac{\eta}{2}\right)\|\nabla f(x_{k+1})\|_{D_{k+1}^{-1}}^2\\
&{}-\frac{1-\theta}{\tilde{L}}D_{\phi^*}(0,\nabla f(x_k))+\frac{\mu}{L}\langle D_{k+1}v_{k+1}, \nabla\phi^*(\nabla f(x_k))\rangle + \frac{\mu}{2}\|v_k\|_{D_{k+1}-D_k}^2.
\end{aligned}
\end{equation}

Now choose
\[
\alpha^2 =  c\frac{\mu}{L},~c\in [1,2], \qquad\text{equivalently}\qquad \mu = L\alpha^2/c.
\]
Then
\[
\frac{\tilde{\alpha}^2}{2\mu} = \frac{\alpha^2}{2\mu(1+\alpha)^2} = \frac{c}{2}\frac{1}{L(1+\alpha)^2}.
\]
Hence the second line of equation \ref{eq:Etilde} becomes less than or equal to
\[
-\min(1,\theta)\frac{1}{\tilde{L}}\left(D_{\phi^*}(0,\nabla f(x_{k+1}))+D_{\phi^*}(\nabla f(x_{k+1}),0)\right) + \left(\frac{c}{2}\frac{1}{L(1+\alpha)^2}+\frac{\eta}{2}\right)\|\nabla f(x_{k+1})\|_{D_{k+1}^{-1}}^2.
\]
To eliminate this term, we assume $\theta\leq 1$, and choose $\eta$ such that
\[
\frac{\eta}{2} \leq \frac{\theta}{\tilde{L}} - \frac{c}{2}\frac{1}{L(1+\alpha)^2} = \frac{\theta(1+\alpha)-c/2}{L(1+\alpha)^2}.
\]
which is positive for every $\theta\in (\frac{1}{1+\alpha},1]$. For simplicity, we choose $\eta=2\frac{\alpha}{L(1+\alpha)^2}$, and $\theta$ such that the equality holds above. Then $\theta=\frac{c/2+\alpha}{1+\alpha}$. When $c<2$, $\theta<1$. And the coefficient of $\|\nabla f(x_{k+1})\|_{D_{k+1}^{-1}}^2$ is $0$. % $$-\frac{1}{\tilde{L}}+\left(\frac{c}{2}\frac{1}{L(1+\alpha)^2}+\frac{\eta}{2}\right)=\frac{c/2-1}{L(1+\alpha)^2}$$ which is nonpositive for $c\in [1,2]$.

Using
\[
\mu = L\alpha^2/c, \qquad
\eta = 2\frac{\alpha}{L(1+\alpha)^2}, \qquad
\delta := \alpha(1+\alpha) - \tilde{\alpha} = \frac{\alpha^2(\alpha+2)}{1+\alpha},
\]
we obtain the coefficient of $-\|v_{k+1}\|_{D_{k+1}}^2$ is
\begin{align*}
&\mu\alpha(\alpha+2) - \frac{\delta^2}{2\eta}\\
&= \frac{1}{c}L\alpha^3(\alpha+2)
- \frac{1}{2}\frac{L(1+\alpha)^2}{2\alpha}\cdot\frac{\alpha^4(\alpha+2)^2}{(1+\alpha)^2}\\
&= \frac{1}{c}L\alpha^3(\alpha+2)- \frac{1}{4}L\alpha^3(\alpha+2)^2=\frac{1}{4}L\alpha^3(\alpha+2)\left(\frac{4}{c}-(\alpha+2)\right).
\end{align*}
To make this coefficient positive, we need
\[c<\frac{4}{\alpha+2}.\]
Since $\alpha^2 = c\mu/L$, this is equivalent to
\[\alpha^2 < \frac{4}{\alpha+2}\frac{\mu}{L}.\]
This condition is satisfied for sufficiently small $\alpha$. In particular, if $c=1$, then $\alpha^2 = \mu/L < 4/(\alpha+2)$ is satisfied for all $\alpha\in (0,1)$.


Therefore, the change of the shifted energy is as follows:
\[
\begin{aligned}
\widetilde{E}_{k+1} - \widetilde{E}_k
\le{}& \frac{\mu}{L}\langle D_{k+1}v_{k+1}, \nabla \phi^*(\nabla f(x_k))\rangle + \frac{\mu}{2}\|v_k\|_{D_{k+1}-D_k}^2\\
&{}- \frac{1-c/2}{L(1+\alpha)^2}D_{\phi^*}(0,\nabla f(x_{k}))+ \left( \alpha + 2 - \frac{4}{c} \right)\frac{1}{4}L\alpha^3(\alpha+2)\|v_{k+1}\|_{D_{k+1}}^2.
\end{aligned}
\]

Use Young's inequality again to bound the cross term on the right-hand side:
\[
\frac{\mu}{L}\langle D_{k+1}v_{k+1}, \nabla \phi^*(\nabla f(x_k))\rangle
\le \frac{1}{4\tilde{L}\xi}\|\nabla \phi^*(\nabla f(x_k))\|_{D_{k+1}}^2 + \xi\frac{\mu^2(1+\alpha)}{L}\|v_{k+1}\|_{D_{k+1}}^2.
\]
The first term on the right-hand side can be written as
\[\frac{1}{4\tilde{L}\xi}\|\nabla \phi^*(\nabla f(x_k))\|_{D_{k+1}}^2 = \frac{1}{4\tilde{L}\xi}\|\nabla f(x_k)\|_{D_k^{-1} D_{k+1} D_k^{-1}}^2.\]
Wrapping up, we have
\[\begin{aligned}
\widetilde{E}_{k+1} - \widetilde{E}_k
\le{}& \frac{1}{4\tilde{L}\xi}\|\nabla f(x_k)\|_{D_k^{-1} D_{k+1} D_k^{-1}}^2 + \frac{\mu}{2}\|v_k\|_{D_{k+1}-D_k}^2\\
&{} + \frac{c/2-1}{L(1+\alpha)^2}D_{\phi^*}(0,\nabla f(x_{k}))+ \left(\left( \alpha + 2 - \frac{4}{c} \right)\frac{1}{4}L\alpha^3(\alpha+2)+\xi\frac{\mu^2(1+\alpha)}{L}\right)\|v_{k+1}\|_{D_{k+1}}^2.
\end{aligned}\]
Then, the boundeness is reduced to the following conditions on the sequences $\{x_k\}$, $\{y_k\}$:
\begin{enumerate}
    \item $$\frac{1}{4\tilde{L}\xi}\|\nabla f(x_k)\|_{D_k^{-1} D_{k+1} D_k^{-1}}^2 < \frac{1-c/2}{L(1+\alpha)^2}D_{\phi^*}(0,\nabla f(x_{k}))$$
    \item $$\frac{\mu}{2}\|v_k\|_{D_{k+1}-D_k}^2 < -\left(\left( \alpha + 2 - \frac{4}{c} \right)\frac{1}{4}L\alpha^3(\alpha+2)+\xi\frac{\mu^2(1+\alpha)}{L}\right)\|v_{k+1}\|_{D_{k+1}}^2.$$
\end{enumerate}
Simplifying the conditions gives
\begin{enumerate}
    \item $$\|\nabla f(x_k)\|_{D_k^{-1} D_{k+1} D_k^{-1}}^2 < \xi\frac{4(1-c/2)}{1+\alpha}D_{\phi^*}(0,\nabla f(x_{k}))$$
    \item $$\|v_k\|_{D_{k+1}-D_k}^2 < -\frac{\alpha}{2c}\left[(c^{2}+4\xi)\,\alpha^{2}
+4(c^{2}-c+\xi)\,\alpha+4c(c-2)\right]\|v_{k+1}\|_{D_{k+1}}^2.$$
\end{enumerate}
Note that both conditions are verifiable. They are satisfied when the step size $\alpha$ is sufficiently small, under appropriate choices of $\xi$ and $c$.

Lastly, choose $c=1$, and set $$\xi =\frac{4 - \alpha^2 - 2\alpha^{1/2}}{4\alpha(\alpha + 1)}.$$ In the small $\alpha$ regime, $\xi$ is positive, and it scales as $1/\alpha-1/(2\sqrt{\alpha})+O(\alpha)$ as $\alpha\to 0$.
\end{proof}


\section{Proofs for Section \ref{sec:accelerated}}\label{app:proofs-accelerated}
% \subsection{Continuous-time ODE}
% We first consider the continuous-time ODE of the scheme,
% \begin{equation}
% \begin{aligned}
% x' &= y - x - \beta \nabla \phi^*(\nabla f(x)),\\
% y' &= x - y - \frac{1}{\mu} \nabla\phi^*(\nabla f(x)).
% \end{aligned}
% \end{equation}

% We will use the Lyapunov function
% \begin{equation}
% \mathcal{E}(x,y;D) = f(x)-f(x^\star) + \frac{\mu}{2}\|y-x^\star\|_{D(p(x))}^2,
% \qquad
% p(x) :=\nabla\phi^*(\nabla f(x)),
% \end{equation}
% where $D(\cdot)$ is the positive definite diagonal matrix defined by
% $$
% D(x)\,x=\nabla\phi(x).
% $$

% \begin{lemma*}[Restatement of Lemma 1]
% Let $\phi(x)=(b,e^x)-(b,x)$. Define $D(x)$ by $D(x)\,x=\nabla\phi(x)$. Then
% $$
% D(x)=\mathrm{diag}\!\left(b.\!*\frac{e^x-\mathbf{1}}{x}\right),
% $$
% where the fraction is entrywise (with the continuous extension at $x_j=0$). Consequently, with $p(x)=\nabla\phi^*(\nabla f(x))$,
% $$
% D(p(x))\,p(x)=\nabla f(x).
% $$
% \end{lemma*}

% \begin{proof}
% Since $\nabla\phi$ is $C^1$ and $\nabla\phi(0)=0$, the fundamental theorem of calculus gives
% $$
% \nabla\phi(x)=\int_0^1 \nabla^2\phi(tx)\,x\dd t=\Bigl(\int_0^1 \nabla^2\phi(tx)\dd t\Bigr)x.
% $$
% Because $\phi$ is separable,
% $$
% \nabla\phi(x)=b.\!*\,e^x-b,
% \qquad
% \nabla^2\phi(x)=\mathrm{diag}\bigl(b.\!*\,e^x\bigr).
% $$
% Therefore,
% $$
% \int_0^1 \nabla^2\phi(tx)\dd t
% =\mathrm{diag}\!\Bigl(\int_0^1 b.\!*\,e^{tx}\dd t\Bigr)
% =\mathrm{diag}\!\left(b.\!*\frac{e^x-\mathbf{1}}{x}\right).
% $$
% Defining $D(x):=\int_0^1 \nabla^2\phi(tx)\dd t$ yields $D(x)\,x=\nabla\phi(x)$ and the stated closed form.

% For the second claim, assume $\nabla f(x)>-b$ so that $p(x)=\nabla\phi^*(\nabla f(x))$ is well defined. Then
% $$
% D(p(x))\,p(x)=\nabla\phi(p(x))=\nabla\phi\bigl(\nabla\phi^*(\nabla f(x))\bigr)=\nabla f(x).
% $$
% \end{proof}

% When $x_j\to 0$, the $j$th diagonal entry of $D(x)$ satisfies
% $$
% b_j\frac{e^{x_j}-1}{x_j} = b_j(x_j + x_j^2 ...)/x_j \to b_j,
% $$
% which equals the $j$th diagonal entry of $\nabla^2\phi(0)$. Thus $D(x)$ admits a continuous extension at $x=0$ by setting $D(0)=\nabla^2\phi(0)$.


% Let $z=(x,y)$ and let $\mathcal{G}(z)$ denote the right-hand side of the ODE. 

\begin{lemma*}[Restatement of Lemma 2]
Let $z(t)=(x(t),y(t))$ be a trajectory of $z'=\mathcal{G}(z)$. Define
$$
p(x(t)):=\nabla\phi^*(\nabla f(x(t))),\qquad \mathcal D(t):=D(p(x(t))),
$$
where $D(\cdot)$ is the diagonal map satisfying $D(s)\,s=\nabla\phi(s)$. Then, for all $t\ge 0$, the following identity holds:
$$
\begin{aligned}
\mathcal{E}'(x,y;\mathcal D)={}&\langle \nabla \mathcal{E}(x,y;\mathcal D),\mathcal{G}(z)\rangle\\
={}& -\mathcal{E}(x,y;\mathcal D) - D_f(x^\star,x) -\beta\|p(x)\|_{\mathcal D(t)}^2 -\frac{\mu}{2}\|x-y\|_{\mathcal D(t)}^2
\\
&\quad + \frac{\mu}{2}\|x-x^\star\|_{\mathcal D(t)}^2+\frac{\mu}{2}\|y-x^\star\|_{D'(t)}^2.
\end{aligned}
$$
\end{lemma*}
% \begin{lemma}
% Let $z(t)=(x(t),y(t))$ be a trajectory of $z'=\mathcal{G}(z)$. Define
% $$
% p(x(t)):=\nabla\phi^*(\nabla f(x(t))),\qquad \mathcal D(t):=D(p(x(t))),
% $$
% where $D(\cdot)$ is the diagonal map satisfying $D(s)\,s=\nabla\phi(s)$. Then, for all $t\ge 0$, the following identity holds:
% $$
% \begin{aligned}
% \mathcal{E}'(x,y;\mathcal D)={}&\langle \nabla \mathcal{E}(x,y;\mathcal D),\mathcal{G}(z)\rangle\\
% ={}& -\mathcal{E}(x,y;\mathcal D) - D_f(x^\star,x) -\beta\|p(x)\|_{\mathcal D(t)}^2 -\frac{\mu}{2}\|x-y\|_{\mathcal D(t)}^2
% \\
% &\quad + \frac{\mu}{2}\|x-x^\star\|_{\mathcal D(t)}^2+\frac{\mu}{2}\|y-x^\star\|_{D'(t)}^2.
% \end{aligned}
% $$
% \end{lemma}

\begin{proof}
Since $D=\mathcal D(t)$ varies with $t$, differentiating the Lyapunov function along the trajectory gives
$$
\mathcal{E}'(x,y;\mathcal D)
=\langle \nabla f(x),x'\rangle+\mu\langle y-x^\star,y'\rangle_{\mathcal D(t)}+\frac{\mu}{2}\|y-x^\star\|_{D'(t)}^2.
$$
Substituting $x'=y-x-\beta p(x)$ and $y'=x-y-\tfrac{1}{\mu}p(x)$, we obtain
$$
\begin{aligned}
{}& \langle \nabla f(x),y-x-\beta p(x)\rangle + \mu\langle y-x^\star,x-y-\tfrac{1}{\mu}p(x)\rangle_{\mathcal D(t)}
\\
={}& \langle \nabla f(x),y-x\rangle-\beta\langle \nabla f(x),p(x)\rangle+\mu\langle y-x^\star,x-y\rangle_{\mathcal D(t)}
-\langle y-x^\star, \mathcal D(t)p(x)\rangle.
\end{aligned}
$$
By definition of $D(\cdot)$ and $p(x)=\nabla\phi^*(\nabla f(x))$, we have
$$
\mathcal D(t)p(x(t))=D(p(x(t)))\,p(x(t)))=\nabla\phi(p(x(t)))=\nabla\phi\bigl(\nabla\phi^*(\nabla f(x(t)))\bigr)=\nabla f(x(t)),
$$
and hence
$$
\langle \nabla f(x),y-x\rangle-\beta\langle \nabla f(x),p(x)\rangle-\langle y-x^\star, \mathcal D(t)p(x)\rangle
=-\beta\|p(x)\|_{\mathcal D(t)}^2-\langle \nabla f(x),x-x^\star\rangle.
$$
Moreover, the polarization identity yields
$$
\mu\langle y-x^\star,x-y\rangle_{\mathcal D(t)}
=-\frac{\mu}{2}\|y-x^\star\|_{\mathcal D(t)}^2-\frac{\mu}{2}\|x-y\|_{\mathcal D(t)}^2+\frac{\mu}{2}\|x-x^\star\|_{\mathcal D(t)}^2.
$$
Using
$$
\langle \nabla f(x),x-x^\star\rangle = f(x)-f(x^\star) + D_f(x^\star,x),
$$
and recalling $\mathcal{E}(x,y;\mathcal D)=f(x)-f(x^\star)+\frac{\mu}{2}\|y-x^\star\|_{\mathcal D(t)}^2$, we conclude the claim.
\end{proof}

\subsection{Discretization and convergence analysis}

The continuous-time dynamics show that the Lyapunov function $\mathcal{E}$ decays exponentially, up to a positive perturbation term $\frac{\mu}{2}\|x-x^\star\|_{D}^2$ and a negative term $\beta\|p(x)\|_{D}^2$. To control the positive perturbation term, we impose Assumption \ref{assump:boundedness}.

% \begin{assumption}\label{assump:boundedness}
% There exists a constant $R>0$ such that for all iterates $x_k$, it holds that % $\|x_k-x^\star\|\le C$ and $\|\nabla f(x_k)\|\le C$.
% $\|x_k-x_\star\|_{D(x_k^\phi)}\leq R$.
% \end{assumption}

% Note that Assumption \ref{assump:boundedness} is a mild assumption that can be derived from the uniform boundedness of $\|x-x_\star\|_{2}$ and $\|\nabla f(x)\|_{\infty}$. Actually, assume both quantities are bounded by a constant $C$. Then we have
% $$
% \|x_k-x^\star\|_{D(x_k^\phi)}^2
% \le \lambda_{\max}(D(x_k^\phi))\|x_k-x^\star\|^2
% \le C^2\lambda_{\max}(D(x_k^\phi)).
% $$
% Here,
% $$D=D(x_k^{\phi})=\mathrm{diag}\left(b.*\frac{e^{x_k^{\phi}}-1}{x_k^{\phi}}\right)=\mathrm{diag}\left(\frac{\nabla f(x_k)}{\log(1+\nabla f(x_k)./b)}\right).$$
% Under the assumption that $\|\nabla f(x_k)\|_\infty\leq C$, it holds $$\lambda_{\max}(D)\leq \frac{C}{\log(1+C/\max_i b_i)}=\mathcal{O}(C).$$
% Therefore, $\|x_k-x^\star\|_{D(x_k^\phi)}^2\leq R^2$ for some constant $R=\mathcal{O}(C^{3/2})$.

\begin{theorem*}[Restatement of Theorem \ref{thm:single-step}]
    % \label{thm:single-step}
Let $z_k=(x_k,y_k)$ be the iterates generated by Algorithm \ref{alg:A2MD} with $\alpha\beta=1$. Assume that Assumption \ref{assump:boundedness} holds. Then
$$
\begin{aligned}
\mathcal{E}(z_{k+1};D_{k+1}) - \mathcal{E}(z_k;D_{k+1})
\le{}& -\alpha \mathcal{E}(z_{k+1};D_{k+1}) + \frac{\alpha\mu}{2}R^2 -  D_{\phi^*}(0,\nabla f(x_{k}))
\\
&\quad + \frac{\alpha^2}{2\mu}\|\nabla \phi^*(\nabla f(x_{k+1}))\|_{D_{k+1}}^2 -  D_{\phi^*}(\nabla f(x_{k+1}),0).
\end{aligned}
$$
\end{theorem*}
\begin{proof}
First, consider the implicit Euler scheme
\begin{equation}
\label{eq:implicit-euler}
\frac{z_{k+1}-z_k}{\alpha}=\mathcal{G}(z_{k+1}).
\end{equation}
Expanding the Lyapunov difference at $z=z_{k+1}$ yields
$$
\begin{aligned}
\mathcal{E}(z_{k+1};D_{k+1}) - \mathcal{E}(z_k;D_{k+1})
={}& \langle \nabla \mathcal{E}(z_{k+1};D_{k+1}), z_{k+1}-z_k\rangle - D_{\mathcal{E}}(z_{k},z_{k+1};D_{k+1})
\\ 
={}& \alpha \langle \nabla \mathcal{E}(z_{k+1};D_{k+1}), \mathcal{G}(z_{k+1})\rangle - D_{f}(x_{k},x_{k+1}) - \frac{\mu}{2}\|y_{k+1}-y_k\|_{D_{k+1}}^2.
\end{aligned}
$$
Comparing \eqref{eq:implicit-euler} with the actual scheme \eqref{eq:scheme}, the $x$-update differs by two terms. This produces additional error terms in the Lyapunov difference:
$$
\begin{aligned}
\mathcal{E}(z_{k+1};D_{k+1}) - \mathcal{E}(z_k;D_{k+1})
\leq{}& \alpha \langle \nabla \mathcal{E}(z_{k+1};D_{k+1}), \mathcal{G}(z_{k+1})\rangle - D_{f}(x_{k},x_{k+1}) - \frac{\mu}{2}\|y_{k+1}-y_k\|_{D_{k+1}}^2
\\
&\quad + \alpha\langle \nabla f(x_{k+1}), y_k-y_{k+1}\rangle
+ \alpha\beta\langle \nabla f(x_{k+1}), \nabla \phi^*(\nabla f(x_{k+1})) - \nabla \phi^*(\nabla f(x_{k}))\rangle.
\end{aligned}
$$
By the continuous-time analysis,
$$
\langle \nabla \mathcal{E}(z_{k+1};D_{k+1}), \mathcal{G}(z_{k+1})\rangle \leq - \mathcal{E}(z_{k+1};D_{k+1}) +\frac{\mu}{2}R^2-\beta\|\nabla\phi^*(\nabla f(x_{k+1}))\|_{D_{k+1}}^2.
$$
For the first cross term, the Cauchy--Schwarz and Young inequalities give
$$
\begin{aligned}
\alpha\langle \nabla f(x_{k+1}), y_k-y_{k+1}\rangle
\leq{}& \frac{\alpha}{\sqrt{\mu}}\|\nabla f(x_{k+1})\|_{D_{k+1}^{-1}} \cdot \sqrt{\mu} \|y_k-y_{k+1}\|_{D_{k+1}}
\\
\leq{}& \frac{\alpha^2}{2\mu}\|\nabla f(x_{k+1})\|_{D_{k+1}^{-1}}^2 + \frac{\mu}{2}\|y_k-y_{k+1}\|_{D_{k+1}}^2.
\end{aligned}
$$
Since $D_{k+1}\nabla\phi^*(\nabla f(x_{k+1}))=\nabla\phi(\nabla\phi^*(\nabla f(x_{k+1})))=\nabla f(x_{k+1})$, we have
$$
\|\nabla f(x_{k+1})\|_{D_{k+1}^{-1}}^2 = \|\nabla \phi^*(\nabla f(x_{k+1}))\|_{D_{k+1}}^2.
$$
For the second cross term, the three-point identity for the Bregman divergence yields
$$
\begin{aligned}
&\langle \nabla f(x_{k+1}), \nabla \phi^*(\nabla f(x_{k+1})) - \nabla \phi^*(\nabla f(x_{k}))\rangle
\\
={}& D_{\phi^*}(\nabla f(x_{k+1}),\nabla f(x_{k})) - D_{\phi^*}(\nabla f(x_{k+1}),0) - D_{\phi^*}(0,\nabla f(x_{k})).
\end{aligned}
$$
Using $\phi^*(0)=0$ and the definition of the Bregman divergence,
$$
D_{\phi^*}(\nabla f(x_{k+1}),0)=\phi^*(\nabla f(x_{k+1}))-\langle \nabla f(x_{k+1}),0\rangle-\phi^*(0)=\phi^*(\nabla f(x_{k+1})).
$$
Combining the above bounds gives
$$
\begin{aligned}
\mathcal{E}(z_{k+1};D_{k+1}) - \mathcal{E}(z_k;D_{k+1})
\leq{}& -\alpha \mathcal{E}(z_{k+1};D_{k+1}) + \frac{\alpha\mu}{2}R^2  - \alpha\beta D_{\phi^*}(0,\nabla f(x_{k}))
\\
&\quad + \frac{\alpha^2}{2\mu}\|\nabla \phi^*(\nabla f(x_{k+1}))\|_{D_{k+1}}^2 - \alpha\beta D_{\phi^*}(\nabla f(x_{k+1}),0)
\\
&\quad +\alpha\beta D_{\phi^*}(\nabla f(x_{k+1}),\nabla f(x_{k})) - D_{f}(x_{k},x_{k+1}).
\end{aligned}
$$
By dual relative smoothness, $ D_{\phi^*}(\nabla f(x_{k+1}),\nabla f(x_{k})) \leq D_f(x_{k},x_{k+1})$. Therefore, it suffices to choose $\alpha$ and $\beta$ such that $\alpha\beta=1$.
\end{proof}

% In the above Lyapunov difference, there is one positive perturbation term and two negative perturbation terms
% $$
% \frac{\alpha^2}{2\mu}\|\nabla \phi^*(\nabla f(x_{k+1}))\|_{D_{k+1}}^2
% \qquad\text{and}\qquad
% - D_{\phi^*}(\nabla f(x_{k+1}),0) - D_{\phi^*}(0,\nabla f(x_{k})).
% $$
% We use the negative terms to offset the positive term. 
% The following lemma makes the relation between these two terms explicit.

% \begin{lemma}
% For any $k\ge 1$, it holds that
% $$
% \frac{1}{2}\|\nabla \phi^*(\nabla f(x_{k}))\|_{D_{k}}^2
% =\frac{1}{2}\mathrm{KL}(c_{P_k}\|b)+\frac{1}{2}\mathrm{KL}(b\|c_{P_k}),
% $$
% and
% $$
% D_{\phi^*}(\nabla f(x_{k}),0)=\mathrm{KL}(c_{P_k}\|b),
% $$
% where $c_{P_k}:=P_k^\top\mathbf{1}_n$ is the column-sum vector of $P_k$, and $P_k$ is the optimal transport matrix computed by Sinkhorn at $x_k$.
% \end{lemma}

% \begin{proof}
% Let $g_k:=\nabla f(x_k)=c_{P_k}-b$. Since $\nabla \phi^*(\eta)=\log\!\bigl(\mathbf{1}+\eta./b\bigr)$, we have
% $$
% \nabla \phi^*(g_k)=\log(c_{P_k}./b).
% $$
% By the definition of $D_k$ and $x^\phi=\nabla\phi^*(\nabla f(x))$, we have
% $$
% D_k\,\nabla \phi^*(g_k)=g_k.
% $$
% Therefore,
% $$
% \|\nabla \phi^*(g_k)\|_{D_k}^2=\big\langle \nabla \phi^*(g_k),\,g_k\big\rangle.
% $$
% Using the Bregman three-point identity with $(\eta,\xi)=(g_k,0)$,
% $$
% D_{\phi^*}(g_k,0)+D_{\phi^*}(0,g_k)
% =\big\langle \nabla\phi^*(g_k)-\nabla\phi^*(0),\,g_k-0\big\rangle
% =\big\langle \nabla\phi^*(g_k),\,g_k\big\rangle,
% $$
% since $\nabla\phi^*(0)=0$. Hence
% $$
% \frac{1}{2}\|\nabla \phi^*(\nabla f(x_k))\|_{D_k}^2
% =\frac{1}{2}\Big(D_{\phi^*}(g_k,0)+D_{\phi^*}(0,g_k)\Big).
% $$
% By the identification $g_k=c_{P_k}-b$ and equation \eqref{eq:Dphistar-KL-special}, we have
% $$
% D_{\phi^*}(g_k,0)=\mathrm{KL}(c_{P_k}\|b),
% \qquad
% D_{\phi^*}(0,g_k)=\mathrm{KL}(b\|c_{P_k}),
% $$
% which gives
% $$
% \frac{1}{2}\|\nabla \phi^*(\nabla f(x_k))\|_{D_k}^2
% =\frac{1}{2}\mathrm{KL}(c_{P_k}\|b)+\frac{1}{2}\mathrm{KL}(b\|c_{P_k}).
% $$
% \end{proof}

% To ensure
% $$
% \frac{\alpha^2}{2\mu}\|\nabla \phi^*(\nabla f(x_{k}))\|_{D_{k}}^2 -  D_{\phi^*}(\nabla f(x_{k}),0) \leq 0,
% $$
% it suffices to choose $\alpha$ such that
% $$
% \alpha=\sqrt{\mu \frac{\mathrm{KL}(c_{P_k}\|b)+\mathrm{KL}(b\|c_{P_k})}{2\,\mathrm{KL}(c_{P_k}\|b)}}.
% $$
% The ratio
% $$
% \frac{\mathrm{KL}(c_{P_k}\|b)+\mathrm{KL}(b\|c_{P_k})}{2\,\mathrm{KL}(c_{P_k}\|b)}
% $$
% measures the relative asymmetry of the KL divergence. In practice, we set $\alpha=\sqrt{\mu}$, which amounts to ignoring this asymmetry. When the marginal deviation of $c_{P_k}$ from $b$ is small, the asymmetry is also small, so $\alpha=\sqrt{\mu}$ is a good approximation.

% \begin{remark}
%     In the above analysis, the negative term $-D_{\phi^*}(0,\nabla f(x_k))$ is not used. Inspired by \citet{chen2025hnagsuperfastacceleratedgradient}, we can enlarge the step size $\alpha$, and offset the extra positivity using this term. More precisely, we want to choose $\alpha$ such that
%     $$\frac{\alpha^2}{2\mu}\|\nabla \phi^*(\nabla f(x_{k}))\|_{D_{k}}^2 -  D_{\phi^*}(\nabla f(x_{k}),0) \leq (1+\alpha) D_{\phi^*}(0,\nabla f(x_{k})).$$
%     Omit the extra positivity $\alpha D_{\phi^*}(0,\nabla f(x_{k}))$ on the right-hand side, we can choose 
%     $$\alpha=\sqrt{2\mu}\sqrt{\frac{D_{\phi^*}(0,\nabla f(x_{k}))+D_{\phi^*}(\nabla f(x_{k}),0)}{\|\nabla \phi^*(\nabla f(x_{k}))\|_{D_{k}}^2}}=\sqrt{2\mu},$$
%     where the last equality is due to the KL formulation of the three terms.
%     This gives a cleaner step size of $\alpha$.
    
% \end{remark}

% \begin{theorem}[Local bound of $\alpha$ in terms of marginal deviations]\label{thm:alpha-bound}
% Let $P_k$ be the optimal transport matrix computed by Sinkhorn at iterate $x_k$, with column sum $c_{P_k}:=P_k^\top \mathbf{1}_n$ and target marginal $b>0$. Define
% $$
% g_k:=c_{P_k}-b,\qquad
% \eta:=\left\|g_k./b\right\|_\infty\le \tfrac{1}{2},
% $$
% and set
% $$
% \alpha:=\sqrt{\mu\,\frac{\mathrm{KL}(c_{P_k}\|b)+\mathrm{KL}(b\|c_{P_k})}{2\,\mathrm{KL}(c_{P_k}\|b)}}.
% $$
% Then
% $$
% \sqrt{\mu}\!\left(1-\frac{\eta}{8}\right)\le \alpha \le \sqrt{\mu}\!\left(1+\frac{\eta}{8}\right).
% $$
% \end{theorem}

% \begin{proof}
% Write $c_{P_k,j}=b_j(1+t_j)$ with $t_j:=(g_k)_j/b_j$ and $\eta=\max_j|t_j|\le \tfrac12$. Define
% $$
% A:=\sum_j b_j t_j^2>0,\qquad B:=\sum_j b_j t_j^3.
% $$

% \noindent\textbf{KL expansions.}
% Using $\log(1+t)=t-\tfrac12 t^2+\tfrac13 t^3-\cdots$ and $\sum_j b_j t_j=(\mathbf{1},g_k)=0$, we obtain
% $$
% \mathrm{KL}(c_{P_k}\|b)=\tfrac12 A-\tfrac16 B+O(\eta^2 A),
% \qquad
% \mathrm{KL}(b\|c_{P_k})=\tfrac12 A-\tfrac13 B+O(\eta^2 A).
% $$

% \noindent\textbf{Ratio.}
% $$
% \frac{\mathrm{KL}(c_{P_k}\|b)+\mathrm{KL}(b\|c_{P_k})}{2\,\mathrm{KL}(c_{P_k}\|b)}
% =\frac{A-\tfrac12 B}{A-\tfrac13 B}+O(\eta^2)
% =1-\frac{\tfrac16 B}{A-\tfrac13 B}+O(\eta^2).
% $$

% \noindent\textbf{Bounding the correction.}
% Since $|B|\le \eta A$ and $A-\tfrac13|B|\ge A(1-\eta/3)>0$, we have
% $$
% \left|\frac{\tfrac16 B}{A-\tfrac13 B}\right|
% \le \frac{\eta/6}{1-\eta/3}\le \frac{\eta}{4}
% \qquad (\eta\le\tfrac12),
% $$
% and therefore
% $$
% 1-\frac{\eta}{4}\le \frac{\mathrm{KL}(c_{P_k}\|b)+\mathrm{KL}(b\|c_{P_k})}{2\,\mathrm{KL}(c_{P_k}\|b)}\le 1+\frac{\eta}{4}.
% $$

% \noindent\textbf{Conclusion.}
% Applying $\sqrt{1-x}\ge 1-\tfrac12 x$ and $\sqrt{1+x}\le 1+\tfrac12 x$ for $x\in[0,1]$, and then multiplying by $\sqrt{\mu}$, yields
% $$
% \sqrt{\mu}\!\left(1-\frac{\eta}{8}\right)\le \alpha \le \sqrt{\mu}\!\left(1+\frac{\eta}{8}\right). \qedhere
% $$
% \end{proof}

% To summarize, when the marginal deviation $\eta$ is small, we can choose $\alpha$ close to $\sqrt{\mu}$ so that the negative term $- D_{\phi^*}(\nabla f(x_{k+1}),0)$ offsets the positive term $\frac{\alpha^2}{2\mu}\|\nabla \phi^*(\nabla f(x_{k+1}))\|_{D_{k+1}}^2$. This also explains the common practice in Sinkhorn-type methods of assuming $\min_j b_j>0$: a strictly positive target marginal limits the relative deviation $g_k./b$ and thus permits a stable step size as in Theorem~\ref{thm:alpha-bound}. We formalize this in the following assumption.

% \begin{assumption}[Marginal deviation]\label{assump:small-deviation}
% There exists a constant $\eta_0>0$ such that for all iterates $x_k$,
% $$
% \eta_k:=\left\|(c_{P_k}-b)./b\right\|_\infty\le \eta_0,
% $$
% where $c_{P_k}:=P_k^\top\mathbf{1}_n$ and $P_k$ is the optimal transport matrix computed by Sinkhorn at $x_k$. Moreover, $\eta_0$ is small enough so that with $\alpha=\sqrt{\mu}\left(1-\frac{\eta_0}{8}\right)$,
% $$
% \frac{\alpha^2}{2\mu}\|\nabla \phi^*(\nabla f(x_{k+1}))\|_{D_{k+1}}^2- D_{\phi^*}(\nabla f(x_{k+1}),0)\le 0
% \qquad \text{for all } k.
% $$
% \end{assumption}

% Assumption~\ref{assump:small-deviation} is only used in the analysis. In practice, we simply set $\alpha=\sqrt{\mu}$, which works well in our experiments.

We next bound the change of the Lyapunov function induced by the change of the metric $D$.

\begin{lemma*}[Restatement of Lemma~\ref{lem:chenge-D}]
    % \label{lem:chenge-D}
Assume that Assumption~\ref{assump:boundedness} holds. Then there exists a constant $C>0$ such that for any $k\ge 1$,
$$
\mathcal{E}(z_k;D_{k+1})-\mathcal{E}(z_k;D_k)
=\frac{\mu}{2}\|y_k-x^\star\|_{D_{k+1}-D_k}^2
\le C\,\frac{\mu}{2}R^2.
$$
\end{lemma*}

\begin{proof}
By the definition of $\mathcal{E}$,
$$
\mathcal{E}(z_k;D_{k+1})-\mathcal{E}(z_k;D_k)
=\frac{\mu}{2}\|y_k-x^\star\|_{D_{k+1}-D_k}^2.
$$
Since $D_k$ is diagonal, we have
$$
\|y_k-x^\star\|_{D_{k+1}-D_k}^2\le \|y_k-x^\star\|^2\,\|D_{k+1}-D_k\|.
$$
It remains to bound $\|D_{k+1}-D_k\|$. Recall that $D(s)\,s=\nabla\phi(s)$ and $\nabla\phi(s)=b.\!*\,\exp(s)-b$. For $s\in\mathbb{R}^m$,
$$
D(s)=\mathrm{diag}\!\left(b.\!*\frac{\exp(s)-\mathbf{1}}{s}\right),
$$
where the fraction is entrywise (with the continuous extension at $s_j=0$). Moreover, with $s=p(x)=\nabla\phi^*(\nabla f(x))$ and $\nabla\phi^*(\eta)=\log(\mathbf{1}+\eta./b)$, we have
$$
s=\log(c_P./b),
\qquad
\frac{\exp(s)-\mathbf{1}}{s}=\frac{c_P./b-\mathbf{1}}{\log(c_P./b)}.
$$
Define $h(g):=g./\log(g+\mathbf{1})$ for $g>0$ (componentwise). Then
$$
D_k=\mathrm{diag}\!\bigl(b.\!*\,h(c_{P_k}./b-\mathbf{1})\bigr).
$$
Under Assumption~\ref{assump:boundedness}, the \textit{marginal deviation} $c_{P_k}./b-1$ is componentwise uniformly bounded, i.e., there exists $M>m>0$ s.t. $c_{P_k}./b-1\in[m,M]$. Since $h$ is Lipschitz on any compact domain, there exists $L=L(m,M)>0$ such that
$$
\|D_{k+1}-D_k\|
\le L\,\max_j b_j\,\big\|c_{P_{k+1}}./b-c_{P_k}./b\big\|_\infty
\le 2L(M-m),
$$
where the last inequality uses $c_{P_k}./b\-\mathbf{1}\in[m,M]$ for all $k$. Therefore,
$$
\begin{aligned}
\mathcal{E}(z_k;D_{k+1})-\mathcal{E}(z_k;D_k)
&\le \frac{\mu}{2}\|y_k-x^\star\|^2\,\|D_{k+1}-D_k\|\\
&\le \frac{\mu}{2}R^2\cdot 2L(M-m).
\end{aligned}
$$
Setting $C:=2L(M-m)$ gives the claim.
\end{proof}

% \begin{remark}
% Following Theorem~\ref{thm:alpha-bound}, one may take $\eta_0=\tfrac12$. A direct calculation gives $L(\eta_0)=(1-\log 2)/\log^2 2$. Therefore, we can take $C=(1-\log 2)/\log^2 2$ in the above lemma.
% \end{remark}

Combining the above results, we conclude that the Lyapunov function decays geometrically up to a bounded perturbation term of order $\mu R^2$.

\begin{theorem*}[Restatement of Theorem~\ref{coro:perturbed-decay}] % \label{coro:perturbed-decay}
Let $z_k=(x_k,y_k)$ be the iterates generated by Algorithm~\ref{alg:A2MD}. Assume that Assumption~\ref{assump:boundedness} holds. Then there exists $C>0$ such that
$$
\mathcal{E}(x_{k+1},y_{k+1};D_{k+1}) \le \left(\frac{1}{1+\alpha}\right)^{k+1}\mathcal{E}(x_k,y_k;D_k) + C\mu R^2.
$$
\end{theorem*}

\begin{proof}
By Theorem \ref{thm:single-step} and Lemma \ref{lem:chenge-D}, there exists $C>0$ such that for all $k$,
$$
\begin{aligned}
(1+\alpha)\mathcal{E}(z_{k+1};D_{k+1})\leq {}&  \mathcal{E}(z_{k};D_{k}) +C\mu R^2 -  D_{\phi^*}(0,\nabla f(x_{k}))
\\
&\quad + \frac{\alpha^2}{2\mu}\|\nabla \phi^*(\nabla f(x_{k+1}))\|_{D_{k+1}}^2 -  D_{\phi^*}(\nabla f(x_{k+1}),0).
\end{aligned}
$$

Take a towering sum and then re-arrange the items, and we have
$$
\begin{aligned}
&(1+\alpha)^{k+1}\mathcal{E}(z_{k+1};D_{k+1})\\ \leq{}& \mathcal{E}(z_{0};D_{0})+\sum_{i=0}^k \left(1+\alpha\right)^{i} C\mu R^2 + \frac{\alpha^2}{2\mu} \|\nabla\phi^*(x_{k+1})\|_{D_{k+1}}^2 -D_{\phi^*}(\nabla f(x_{k+1}),0)- (1+\alpha)^k D_{\phi^*}(0,\nabla f(x_0))  \\&{}+ \sum_{i=1}^{k} (1+\alpha)^{i-1} \left( \frac{\alpha^2}{2\mu}\|\nabla \phi^*(\nabla f(x_{i}))\|_{D_{i}}^2 -  D_{\phi^*}(\nabla f(x_{i}),0) - (1+\alpha)D_{\phi^*}(0,\nabla f(x_{i}))\right).
\end{aligned}
$$
As for all $i$, $$\|\nabla \phi^*(\nabla f(x_{i}))\|_{D_{i}}^2 = D_{\phi^*}(\nabla f(x_{i}),0) + D_{\phi^*}(0,\nabla f(x_{i})),$$
choosing $\alpha=\sqrt{2\mu}$ makes $\alpha^2/(2\mu)=1$, so the intermediate terms are all non-positive. The conclusion follows from the fact that $\sum_{i=0}^k(1+\alpha)^i/(1+\alpha)^{k+1}$ is finite.
\end{proof}

% \begin{proof}
% By Theorem \ref{thm:single-step} and Lemma \ref{lem:chenge-D}, there exists $C>0$ such that for all $k$,
% $$
% \begin{aligned}
% (1+\alpha)\mathcal{E}(z_{k+1};D_{k+1})\leq {}&  \mathcal{E}(z_{k};D_{k}) +C\mu R^2 -  D_{\phi^*}(0,\nabla f(x_{k}))
% \\
% &\quad + \frac{\alpha^2}{2\mu}\|\nabla \phi^*(\nabla f(x_{k+1}))\|_{D_{k+1}}^2 -  D_{\phi^*}(\nabla f(x_{k+1}),0).
% \end{aligned}
% $$

% Take a towering sum and then re-arrange the items, and we have
% $$
% \begin{aligned}
% &(1+\alpha)^{k+1}\mathcal{E}(z_{k+1};D_{k+1})\\ \leq{}& \mathcal{E}(z_{0};D_{0})+\sum_{i=0}^k \left(1+\alpha\right)^{i} C\mu R^2 + \frac{\alpha^2}{2\mu} \|\nabla\phi^*(x_{k+1})\|_{D_{k+1}}^2 -D_{\phi^*}(\nabla f(x_{k+1}),0)- (1+\alpha)^k D_{\phi^*}(0,\nabla f(x_0))  \\&{}+ \sum_{i=1}^{k} (1+\alpha)^{i-1} \left( \frac{\alpha^2}{2\mu}\|\nabla \phi^*(\nabla f(x_{i}))\|_{D_{i}}^2 -  D_{\phi^*}(\nabla f(x_{i}),0) - (1+\alpha)D_{\phi^*}(0,\nabla f(x_{i}))\right).
% \end{aligned}
% $$
% As for all $i$, $$\|\nabla \phi^*(\nabla f(x_{i}))\|_{D_{i}}^2 = D_{\phi^*}(\nabla f(x_{i}),0) + D_{\phi^*}(0,\nabla f(x_{i})),$$
% choosing $\alpha=\sqrt{2\mu}$ makes $\alpha^2/(2\mu)=1$, so the intermediate terms are all non-positive. The conclusion follows from the fact that $\sum_{i=0}^k(1+\alpha)^i/(1+\alpha)^{k+1}$ is finite  .
% \end{proof}

Finally, we can analyze the convergence of Algorithm \ref{alg:A2MD-homotopy}. We adopt the notation $$\mathcal{E}(x,y,\mu;D):=f(x)-f(x^\star)+\frac{\mu}{2}\|y-x\|_D^2$$ to explicitly indicate the dependence of $\mathcal{E}$ on $\mu$.

% \begin{theorem}
% 		Choose $(x_0, y_0)$ and $\mu_0$ satisfying $\mathcal E(x_0, y_0, \mu_0) \leq (R^2+ 1)\mu_{0}$.
% 		Let $(x_{k}, y_{k}, \mu_{k})$ be generated by Algorithm \ref{alg:A2MD-homotopy}. Assume that assumption \ref{assump:boundedness} holds. Then we have
% 		\begin{equation}\label{eq:Ek-ek}
% 			\mathcal E(x_{k}, y_{k}, \mu_{k})\leq (R^2+1)\mu_{k}\quad\forall\,k\geq0,
% 		\end{equation}
% Moreover, let $M_k:=\sum_{i=0}^{k}m_i $ be the overall iteration steps after the $k$-th outer iteration and $C = \frac{(\sqrt{2}-1)}{(\sqrt{2L_F}+\sqrt{2\mu_0})\ln (2(R^2+1))}$, then we have the global convergence rate
% \begin{equation}\label{eq:rate}
% 	\mathcal E(x_k,y_k,\mu_k)\leq \frac{R^2+1}{\left (C M_k+\mu_0^{-1/2}\right )^{2}}\quad\forall\,k\geq 0.
% \end{equation}
% 		Thus, it takes $M_k=O(\sqrt{1/\mu})$ iterations to get the accuracy $\mathcal E(x_{k}, y_{k}, \mu_{k}) = O(\mu)$. 
% 	\end{theorem}

\begin{theorem}[Restatement of Theorem~\ref{thm:homotopy}] % \label{thm:homotopy}
Choose $(x_0,y_0)$ and $\mu_0$ such that
$$
\mathcal E(x_0,y_0;\mu_0,D)\le (R^2+1)\mu_0.
$$
Let $(x_k,y_k,\mu_k)$ be generated by Algorithm~\ref{alg:A2MD-homotopy}. Assume that Assumption~\ref{assump:boundedness} holds. Then
\begin{equation}% \label{eq:Ek-ek}
\mathcal E(x_k,y_k;\mu_k,D)\le (R^2+1)\mu_k,
\qquad\forall\,k\ge 0.
\end{equation}
Moreover, let
$
M_k:=\sum_{i=0}^k m_i
$
be the total number of inner iterations after the $k$th outer loop, and
$
C_*:=\frac{\sqrt{2}-1}{(\sqrt{2L_F}+\sqrt{2\mu_0})\ln\bigl(2(R^2+1)\bigr)}
$
be a constant. Then
\begin{equation}\label{eq:rate}
\mathcal E(x_k,y_k;\mu_k,D)\le \frac{R^2+1}{\bigl(C_*M_k+\mu_0^{-1/2}\bigr)^2}
\qquad\forall\,k\ge 0.
\end{equation}
In particular, it takes $M_k=O(\mu^{-1/2})$ iterations to reach accuracy
$
\mathcal E(x_k,y_k;\mu_k,D)=O(\mu).
$
\end{theorem}
\begin{proof}[Proof of Theorem~\ref{thm:homotopy}]
The proof uses the one-step sufficient decay estimate established earlier for the Sinkhorn map, together with the discrete Lyapunov estimate for the accelerated scheme.

For a fixed value of $\mu$, the inner iteration satisfies a perturbed contraction of the form
\begin{equation}\label{eq:inner-contraction}
(1+\alpha)\mathcal E(x_{j+1},y_{j+1};\mu,D)\le \mathcal E(x_j,y_j;\mu,D)+\mu R^2,
\qquad \alpha=\sqrt{2\mu}.
\end{equation}
Here the negative term coming from the Sinkhorn step is controlled by the sufficient decay lemma, while Assumption~\ref{assump:boundedness} bounds the perturbation term generated by the weighted norm in the Lyapunov function.

Iterating \eqref{eq:inner-contraction} for $m_k$ inner steps gives
\begin{equation}\label{eq:mk-step}
\mathcal E_{k+1}\le (1+\alpha_k)^{-m_k}\mathcal E_k+\mu_{k+1}R^2,
\end{equation}
where $\mathcal E_k:=\mathcal E(x_k,y_k;\mu_k,D)$. The choice of $m_k$ in Algorithm~\ref{alg:A2MD-homotopy} ensures that
$$
(1+\alpha_k)^{-m_k}\le \frac{1}{2(R^2+1)}.
$$
Therefore,
$$
\mathcal E_{k+1}\le \frac{1}{2(R^2+1)}\mathcal E_k+\mu_{k+1}R^2.
$$
Assume inductively that
$$
\mathcal E_k\le (R^2+1)\mu_k.
$$
Since $\mu_{k+1}=\mu_k/2$, we obtain
$$
\mathcal E_{k+1}\le \frac{1}{2}\mu_k+\mu_{k+1}R^2=(R^2+1)\mu_{k+1},
$$
which proves \eqref{eq:Ek-ek}.

It remains to express $\mu_k$ in terms of the total number $M_k$ of inner steps. Since
$$
\mu_k=\mu_0\,2^{-k},
\qquad
\alpha_k=\sqrt{2\mu_k},
$$
and $m_k$ is chosen proportional to $\mu_k^{-1/2}$, summing the geometric progression yields
$$
M_k\asymp \sum_{i=0}^k \mu_i^{-1/2}
\asymp \mu_k^{-1/2}-\mu_0^{-1/2}.
$$
Equivalently,
$$
\mu_k\asymp \frac{1}{\bigl(C_*M_k+\mu_0^{-1/2}\bigr)^2},
$$
for the constant $C_*$ defined above. Substituting this relation into \eqref{eq:Ek-ek} gives \eqref{eq:rate}. The final complexity bound follows immediately.
\end{proof}




%%%%%%%%%%%%%%%%%%%%%%%%%%%%%%%%%%%%%%%%%%%%%%%%%%%%%%%%%%%%

% \newpage
% \section*{NeurIPS Paper Checklist}


% \begin{enumerate}

% \item {\bf Claims}
%     \item[] Question: Do the main claims made in the abstract and introduction accurately reflect the paper's contributions and scope?
%     \item[] Answer: \answerYes{}
%     \item[] Justification: The abstract and introduction clearly state the two main contributions: (1) interpreting Sinkhorn as dual mirror descent with relative smoothness constant $L=1$, and (2) proposing Acc-Sinkhorn with a provable $\mathcal{O}(1/k^2)$ convergence rate and $O(n^2)$ per-iteration cost. These claims are supported by Theorems 1, 2 and the numerical experiments in Section 5.

%     \item[] Guidelines:
%     \begin{itemize}
%         \item The answer NA means that the abstract and introduction do not include the claims made in the paper.
%         \item The abstract and/or introduction should clearly state the claims made, including the contributions made in the paper and important assumptions and limitations. A No or NA answer to this question will not be perceived well by the reviewers. 
%         \item The claims made should match theoretical and experimental results, and reflect how much the results can be expected to generalize to other settings. 
%         \item It is fine to include aspirational goals as motivation as long as it is clear that these goals are not attained by the paper. 
%     \end{itemize}

% \item {\bf Limitations}
%     \item[] Question: Does the paper discuss the limitations of the work performed by the authors?
%     \item[] Answer: \answerYes{} % Replace by \answerYes{}, \answerNo{}, or \answerNA{}.
%     \item[] Justification: The paper discusses the limitations of the work in Section 1.
%     \item[] Guidelines:
%     \begin{itemize}
%         \item The answer NA means that the paper has no limitation while the answer No means that the paper has limitations, but those are not discussed in the paper. 
%         \item The authors are encouraged to create a separate "Limitations" section in their paper.
%         \item The paper should point out any strong assumptions and how robust the results are to violations of these assumptions (e.g., independence assumptions, noiseless settings, model well-specification, asymptotic approximations only holding locally). The authors should reflect on how these assumptions might be violated in practice and what the implications would be.
%         \item The authors should reflect on the scope of the claims made, e.g., if the approach was only tested on a few datasets or with a few runs. In general, empirical results often depend on implicit assumptions, which should be articulated.
%         \item The authors should reflect on the factors that influence the performance of the approach. For example, a facial recognition algorithm may perform poorly when image resolution is low or images are taken in low lighting. Or a speech-to-text system might not be used reliably to provide closed captions for online lectures because it fails to handle technical jargon.
%         \item The authors should discuss the computational efficiency of the proposed algorithms and how they scale with dataset size.
%         \item If applicable, the authors should discuss possible limitations of their approach to address problems of privacy and fairness.
%         \item While the authors might fear that complete honesty about limitations might be used by reviewers as grounds for rejection, a worse outcome might be that reviewers discover limitations that aren't acknowledged in the paper. The authors should use their best judgment and recognize that individual actions in favor of transparency play an important role in developing norms that preserve the integrity of the community. Reviewers will be specifically instructed to not penalize honesty concerning limitations.
%     \end{itemize}

% \item {\bf Theory assumptions and proofs}
%     \item[] Question: For each theoretical result, does the paper provide the full set of assumptions and a complete (and correct) proof?
%     \item[] Answer: \answerYes{}
%     \item[] Justification: All theorems and lemmas state their assumptions explicitly. Proofs for Section 3 are in Appendix A, the normalized Sinkhorn analysis is in Appendix B, Assumption~1 is discussed in Appendices C--D, and the proofs for Section 4 are in Appendix E.
%     \item[] Guidelines:
%     \begin{itemize}
%         \item The answer NA means that the paper does not include theoretical results. 
%         \item All the theorems, formulas, and proofs in the paper should be numbered and cross-referenced.
%         \item All assumptions should be clearly stated or referenced in the statement of any theorems.
%         \item The proofs can either appear in the main paper or the supplemental material, but if they appear in the supplemental material, the authors are encouraged to provide a short proof sketch to provide intuition. 
%         \item Inversely, any informal proof provided in the core of the paper should be complemented by formal proofs provided in appendix or supplemental material.
%         \item Theorems and Lemmas that the proof relies upon should be properly referenced. 
%     \end{itemize}

%     \item {\bf Experimental result reproducibility}
%     \item[] Question: Does the paper fully disclose all the information needed to reproduce the main experimental results of the paper to the extent that it affects the main claims and/or conclusions of the paper (regardless of whether the code and data are provided or not)?
%     \item[] Answer: \answerYes{} % Replace by \answerYes{}, \answerNo{}, or \answerNA{}.
%     \item[] Justification: Section 5 describes the experimental setup in detail, including dataset generation, cost matrix construction, marginal definitions, stopping criteria, and hardware. The paper states that code is publicly available and random seeds were fixed.

%     \item[] Guidelines:
%     \begin{itemize}
%         \item The answer NA means that the paper does not include experiments.
%         \item If the paper includes experiments, a No answer to this question will not be perceived well by the reviewers: Making the paper reproducible is important, regardless of whether the code and data are provided or not.
%         \item If the contribution is a dataset and/or model, the authors should describe the steps taken to make their results reproducible or verifiable. 
%         \item Depending on the contribution, reproducibility can be accomplished in various ways. For example, if the contribution is a novel architecture, describing the architecture fully might suffice, or if the contribution is a specific model and empirical evaluation, it may be necessary to either make it possible for others to replicate the model with the same dataset, or provide access to the model. In general. releasing code and data is often one good way to accomplish this, but reproducibility can also be provided via detailed instructions for how to replicate the results, access to a hosted model (e.g., in the case of a large language model), releasing of a model checkpoint, or other means that are appropriate to the research performed.
%         \item While NeurIPS does not require releasing code, the conference does require all submissions to provide some reasonable avenue for reproducibility, which may depend on the nature of the contribution. For example
%         \begin{enumerate}
%             \item If the contribution is primarily a new algorithm, the paper should make it clear how to reproduce that algorithm.
%             \item If the contribution is primarily a new model architecture, the paper should describe the architecture clearly and fully.
%             \item If the contribution is a new model (e.g., a large language model), then there should either be a way to access this model for reproducing the results or a way to reproduce the model (e.g., with an open-source dataset or instructions for how to construct the dataset).
%             \item We recognize that reproducibility may be tricky in some cases, in which case authors are welcome to describe the particular way they provide for reproducibility. In the case of closed-source models, it may be that access to the model is limited in some way (e.g., to registered users), but it should be possible for other researchers to have some path to reproducing or verifying the results.
%         \end{enumerate}
%     \end{itemize}


% \item {\bf Open access to data and code}
%     \item[] Question: Does the paper provide open access to the data and code, with sufficient instructions to faithfully reproduce the main experimental results, as described in supplemental material?
%     \item[] Answer: \answerYes{} % Replace by \answerYes{}, \answerNo{}, or \answerNA{}.
%     \item[] Justification: The paper states that code is publicly available. Synthetic data is generated procedurally as described in Section 5.1, and the word embedding experiments use publicly available MUSE vectors from \citet{conneau2018word}.
%     \item[] Guidelines:
%     \begin{itemize}
%         \item The answer NA means that paper does not include experiments requiring code.
%         \item Please see the NeurIPS code and data submission guidelines (\url{https://nips.cc/public/guides/CodeSubmissionPolicy}) for more details.
%         \item While we encourage the release of code and data, we understand that this might not be possible, so “No” is an acceptable answer. Papers cannot be rejected simply for not including code, unless this is central to the contribution (e.g., for a new open-source benchmark).
%         \item The instructions should contain the exact command and environment needed to run to reproduce the results. See the NeurIPS code and data submission guidelines (\url{https://nips.cc/public/guides/CodeSubmissionPolicy}) for more details.
%         \item The authors should provide instructions on data access and preparation, including how to access the raw data, preprocessed data, intermediate data, and generated data, etc.
%         \item The authors should provide scripts to reproduce all experimental results for the new proposed method and baselines. If only a subset of experiments are reproducible, they should state which ones are omitted from the script and why.
%         \item At submission time, to preserve anonymity, the authors should release anonymized versions (if applicable).
%         \item Providing as much information as possible in supplemental material (appended to the paper) is recommended, but including URLs to data and code is permitted.
%     \end{itemize}


% \item {\bf Experimental setting/details}
%     \item[] Question: Does the paper specify all the training and test details (e.g., data splits, hyperparameters, how they were chosen, type of optimizer, etc.) necessary to understand the results?
%     \item[] Answer: \answerYes{} % Replace by \answerYes{}, \answerNo{}, or \answerNA{}.
%     \item[] Justification: Section 5 provides dataset sizes, cost matrix construction, regularization parameters $\varepsilon$, stopping criteria, initialization, and hardware details. Algorithm hyperparameters ($\mu_0$, $m_0$) are discussed and noted to be stable.

%     \item[] Guidelines:
%     \begin{itemize}
%         \item The answer NA means that the paper does not include experiments.
%         \item The experimental setting should be presented in the core of the paper to a level of detail that is necessary to appreciate the results and make sense of them.
%         \item The full details can be provided either with the code, in appendix, or as supplemental material.
%     \end{itemize}

% \item {\bf Experiment statistical significance}
%     \item[] Question: Does the paper report error bars suitably and correctly defined or other appropriate information about the statistical significance of the experiments?
%     \item[] Answer: \answerNo{} % Replace by \answerYes{}, \answerNo{}, or \answerNA{}.
%     \item[] Justification: This is a theoretical paper on a deterministic optimization algorithm. The experiments are meant to illustrate the theoretical results and are not intended to be a comprehensive empirical evaluation of the method.

%     \item[] Guidelines:
%     \begin{itemize}
%         \item The answer NA means that the paper does not include experiments.
%         \item The authors should answer "Yes" if the results are accompanied by error bars, confidence intervals, or statistical significance tests, at least for the experiments that support the main claims of the paper.
%         \item The factors of variability that the error bars are capturing should be clearly stated (for example, train/test split, initialization, random drawing of some parameter, or overall run with given experimental conditions).
%         \item The method for calculating the error bars should be explained (closed form formula, call to a library function, bootstrap, etc.)
%         \item The assumptions made should be given (e.g., Normally distributed errors).
%         \item It should be clear whether the error bar is the standard deviation or the standard error of the mean.
%         \item It is OK to report 1-sigma error bars, but one should state it. The authors should preferably report a 2-sigma error bar than state that they have a 96\% CI, if the hypothesis of Normality of errors is not verified.
%         \item For asymmetric distributions, the authors should be careful not to show in tables or figures symmetric error bars that would yield results that are out of range (e.g. negative error rates).
%         \item If error bars are reported in tables or plots, The authors should explain in the text how they were calculated and reference the corresponding figures or tables in the text.
%     \end{itemize}

% \item {\bf Experiments compute resources}
%     \item[] Question: For each experiment, does the paper provide sufficient information on the computer resources (type of compute workers, memory, time of execution) needed to reproduce the experiments?
%     \item[] Answer: \answerYes{} % Replace by \answerYes{}, \answerNo{}, or \answerNA{}.
%     \item[] Justification: Section 5 specifies the hardware: a MacBook Pro (14-inch, 2024) with Apple Silicon M4, 24 GB unified memory, 1 TB SSD, running MATLAB R2025b. Wall-clock times are reported in Tables 1 and 2.

%     \item[] Guidelines:
%     \begin{itemize}
%         \item The answer NA means that the paper does not include experiments.
%         \item The paper should indicate the type of compute workers CPU or GPU, internal cluster, or cloud provider, including relevant memory and storage.
%         \item The paper should provide the amount of compute required for each of the individual experimental runs as well as estimate the total compute. 
%         \item The paper should disclose whether the full research project required more compute than the experiments reported in the paper (e.g., preliminary or failed experiments that didn't make it into the paper). 
%     \end{itemize}
    
% \item {\bf Code of ethics}
%     \item[] Question: Does the research conducted in the paper conform, in every respect, with the NeurIPS Code of Ethics \url{https://neurips.cc/public/EthicsGuidelines}?
%     \item[] Answer: \answerYes{} % Replace by \answerYes{}, \answerNo{}, or \answerNA{}.
%     \item[] Justification: The paper presents foundational algorithmic research on optimal transport with no foreseeable ethical concerns. It does not involve human subjects, sensitive data, or dual-use risks.

%     \item[] Guidelines:
%     \begin{itemize}
%         \item The answer NA means that the authors have not reviewed the NeurIPS Code of Ethics.
%         \item If the authors answer No, they should explain the special circumstances that require a deviation from the Code of Ethics.
%         \item The authors should make sure to preserve anonymity (e.g., if there is a special consideration due to laws or regulations in their jurisdiction).
%     \end{itemize}


% \item {\bf Broader impacts}
%     \item[] Question: Does the paper discuss both potential positive societal impacts and negative societal impacts of the work performed?
%     \item[] Answer: \answerNA{} % Replace by \answerYes{}, \answerNo{}, or \answerNA{}.
%     \item[] Justification: This paper presents foundational research on accelerating a classical optimization algorithm with no direct path to negative societal impact.

%     \item[] Guidelines:
%     \begin{itemize}
%         \item The answer NA means that there is no societal impact of the work performed.
%         \item If the authors answer NA or No, they should explain why their work has no societal impact or why the paper does not address societal impact.
%         \item Examples of negative societal impacts include potential malicious or unintended uses (e.g., disinformation, generating fake profiles, surveillance), fairness considerations (e.g., deployment of technologies that could make decisions that unfairly impact specific groups), privacy considerations, and security considerations.
%         \item The conference expects that many papers will be foundational research and not tied to particular applications, let alone deployments. However, if there is a direct path to any negative applications, the authors should point it out. For example, it is legitimate to point out that an improvement in the quality of generative models could be used to generate deepfakes for disinformation. On the other hand, it is not needed to point out that a generic algorithm for optimizing neural networks could enable people to train models that generate Deepfakes faster.
%         \item The authors should consider possible harms that could arise when the technology is being used as intended and functioning correctly, harms that could arise when the technology is being used as intended but gives incorrect results, and harms following from (intentional or unintentional) misuse of the technology.
%         \item If there are negative societal impacts, the authors could also discuss possible mitigation strategies (e.g., gated release of models, providing defenses in addition to attacks, mechanisms for monitoring misuse, mechanisms to monitor how a system learns from feedback over time, improving the efficiency and accessibility of ML).
%     \end{itemize}
    
% \item {\bf Safeguards}
%     \item[] Question: Does the paper describe safeguards that have been put in place for responsible release of data or models that have a high risk for misuse (e.g., pretrained language models, image generators, or scraped datasets)?
%     \item[] Answer: \answerNA{} % Replace by \answerYes{}, \answerNo{}, or \answerNA{}.
%     \item[] Justification: The paper does not release pretrained models, scraped datasets, or any assets that pose misuse risks.

%     \item[] Guidelines:
%     \begin{itemize}
%         \item The answer NA means that the paper poses no such risks.
%         \item Released models that have a high risk for misuse or dual-use should be released with necessary safeguards to allow for controlled use of the model, for example by requiring that users adhere to usage guidelines or restrictions to access the model or implementing safety filters. 
%         \item Datasets that have been scraped from the Internet could pose safety risks. The authors should describe how they avoided releasing unsafe images.
%         \item We recognize that providing effective safeguards is challenging, and many papers do not require this, but we encourage authors to take this into account and make a best faith effort.
%     \end{itemize}

% \item {\bf Licenses for existing assets}
%     \item[] Question: Are the creators or original owners of assets (e.g., code, data, models), used in the paper, properly credited and are the license and terms of use explicitly mentioned and properly respected?
%     \item[] Answer: \answerYes{} % Replace by \answerYes{}, \answerNo{}, or \answerNA{}.
%     \item[] Justification: The paper cites the original sources for all datasets and code used, including the MUSE word vectors \citep{conneau2018word} and fastText embeddings \citep{bojanowski2017enriching}.

%     \item[] Guidelines:
%     \begin{itemize}
%         \item The answer NA means that the paper does not use existing assets.
%         \item The authors should cite the original paper that produced the code package or dataset.
%         \item The authors should state which version of the asset is used and, if possible, include a URL.
%         \item The name of the license (e.g., CC-BY 4.0) should be included for each asset.
%         \item For scraped data from a particular source (e.g., website), the copyright and terms of service of that source should be provided.
%         \item If assets are released, the license, copyright information, and terms of use in the package should be provided. For popular datasets, \url{paperswithcode.com/datasets} has curated licenses for some datasets. Their licensing guide can help determine the license of a dataset.
%         \item For existing datasets that are re-packaged, both the original license and the license of the derived asset (if it has changed) should be provided.
%         \item If this information is not available online, the authors are encouraged to reach out to the asset's creators.
%     \end{itemize}

% \item {\bf New assets}
%     \item[] Question: Are new assets introduced in the paper well documented and is the documentation provided alongside the assets?
%     \item[] Answer: \answerYes{} % Replace by \answerYes{}, \answerNo{}, or \answerNA{}.
%     \item[] Justification: The paper releases code for Acc-Sinkhorn. The paper states that the code is publicly available.

%     \item[] Guidelines:
%     \begin{itemize}
%         \item The answer NA means that the paper does not release new assets.
%         \item Researchers should communicate the details of the dataset/code/model as part of their submissions via structured templates. This includes details about training, license, limitations, etc. 
%         \item The paper should discuss whether and how consent was obtained from people whose asset is used.
%         \item At submission time, remember to anonymize your assets (if applicable). You can either create an anonymized URL or include an anonymized zip file.
%     \end{itemize}

% \item {\bf Crowdsourcing and research with human subjects}
%     \item[] Question: For crowdsourcing experiments and research with human subjects, does the paper include the full text of instructions given to participants and screenshots, if applicable, as well as details about compensation (if any)? 
%     \item[] Answer: \answerNA{} % Replace by \answerYes{}, \answerNo{}, or \answerNA{}.
%     \item[] Justification: The paper does not involve crowdsourcing or research with human subjects.

%     \item[] Guidelines:
%     \begin{itemize}
%         \item The answer NA means that the paper does not involve crowdsourcing nor research with human subjects.
%         \item Including this information in the supplemental material is fine, but if the main contribution of the paper involves human subjects, then as much detail as possible should be included in the main paper. 
%         \item According to the NeurIPS Code of Ethics, workers involved in data collection, curation, or other labor should be paid at least the minimum wage in the country of the data collector. 
%     \end{itemize}

% \item {\bf Institutional review board (IRB) approvals or equivalent for research with human subjects}
%     \item[] Question: Does the paper describe potential risks incurred by study participants, whether such risks were disclosed to the subjects, and whether Institutional Review Board (IRB) approvals (or an equivalent approval/review based on the requirements of your country or institution) were obtained?
%     \item[] Answer: \answerNA{} % Replace by \answerYes{}, \answerNo{}, or \answerNA{}.
%     \item[] Justification: The paper does not involve human subjects and therefore requires no IRB approval.

%     \item[] Guidelines:
%     \begin{itemize}
%         \item The answer NA means that the paper does not involve crowdsourcing nor research with human subjects.
%         \item Depending on the country in which research is conducted, IRB approval (or equivalent) may be required for any human subjects research. If you obtained IRB approval, you should clearly state this in the paper. 
%         \item We recognize that the procedures for this may vary significantly between institutions and locations, and we expect authors to adhere to the NeurIPS Code of Ethics and the guidelines for their institution. 
%         \item For initial submissions, do not include any information that would break anonymity (if applicable), such as the institution conducting the review.
%     \end{itemize}

% \item {\bf Declaration of LLM usage}
%     \item[] Question: Does the paper describe the usage of LLMs if it is an important, original, or non-standard component of the core methods in this research? Note that if the LLM is used only for writing, editing, or formatting purposes and does not impact the core methodology, scientific rigorousness, or originality of the research, declaration is not required.
%     %this research? 
%     \item[] Answer: \answerNA{} % Replace by \answerYes{}, \answerNo{}, or \answerNA{}.
%     \item[] Justification: LLMs are not used as a component of the core methods in this research.
%     \item[] Guidelines:
%     \begin{itemize}
%         \item The answer NA means that the core method development in this research does not involve LLMs as any important, original, or non-standard components.
%         \item Please refer to our LLM policy (\url{https://neurips.cc/Conferences/2025/LLM}) for what should or should not be described.
%     \end{itemize}

% \end{enumerate}




\end{document}
